\chapter[The Deligne series IV: final theorems and named inputs]{The Deligne series IV: the final theorems and the named inputs}\label{ch:L27}
\begingroup\setlength{\emergencystretch}{3em}\sloppy

\providecommand{\ltwentyseveng}{\mathfrak g}
\providecommand{\ltwentysevenIns}{I_2^{\mathrm{ns}}}
\providecommand{\ltwentysevenHilb}{\operatorname{Hilb}}
\providecommand{\ltwentysevenvac}{\mathbf 1}
\providecommand{\ltwentysevenJ}{J_\infty}
\providecommand{\ltwentysevenF}{\mathcal F}
\providecommand{\ltwentysevengr}{\operatorname{gr}}
\providecommand{\ltwentysevendg}{d_{\mathfrak g}}
\ifdefined\conjecture\else\newtheorem{conjecture}[definition]{Conjecture}\fi

This chapter was written against an earlier state of the LeanVOA library; the statements of
the literature inputs and the partial results of Section~\ref{l27:sec:partial} were revised against a later state of the
source (28 October 2026), following the project's audit of the literature inputs
(\texttt{Notes/Textbook/referee/DELIGNE\_INPUTS\_PRECISION.md}). What the machine confirmed, and when, is stated once, in
Section~\ref{l27:sec:verification}. For this edition, \texttt{scripts/textbook/sync.py} checks that every Lean quotation below occurs verbatim (up to whitespace and comments) in the source files of that release.

\paragraph{What this chapter is about.} The \emph{Deligne series} is the list of simple Lie algebras
$A_1\subset A_2\subset G_2\subset D_4\subset F_4\subset E_6\subset E_7\subset E_8$. At the level $k=-1-h^\vee/6$
($h^\vee$ the dual Coxeter number) the simple affine vertex algebra $L_k(\ltwentyseveng)$ is expected to be
\emph{classically free}: its graded algebra for Li's filtration is the arc algebra of its $C_2$-algebra, and the
$C_2$-algebra is $\operatorname{Sym}\ltwentyseveng$ modulo an explicit space $\ltwentysevenIns$ of quadrics.
Chapters~\ref{ch:L24}--\ref{ch:L26} prove the pieces. This chapter states the result: for each of the seven
higher-rank members $X\in\{A_2,G_2,D_4,F_4,E_6,E_7,E_8\}$ the library has one theorem \lean{X_classicallyFree_final}
with that conclusion, and this chapter states all seven exactly, checks them against the intended mathematics, and
records, for each type, which inputs are proved and which are still hypotheses.

\paragraph{The result in one sentence.} For $A_2$ the final theorem has \emph{no hypotheses}: the library proves that
$L_{-3/2}(\mathfrak{sl}_3)$ is classically free and that $R_V=\operatorname{Sym}\mathfrak{sl}_3/(\ltwentysevenIns)$ for eight
explicit quadrics (\lean{VOA.DeligneAudit.A2.classicallyFree}, re-exported as \lean{A2_classicallyFree_final}). Machine
status: confirmed. The aggregate build was green on 2026-10-28, and \texttt{\#print axioms} on both declarations reports
\lean{[propext, Classical.choice, Quot.sound]} (Section~\ref{l27:sec:verification}). For the
other six types the final theorems are proved \emph{conditionally}. For $D_4$ one published input remains: Input~T,
the Kus--Littelmann fusion-product bound. Input~R is proved in Lean for $D_4$, so the theorem
\lean{D4_classicallyFree_of_inputT} (Theorem~\ref{l27:thm:d4inputT}) has Input~T as its only hypothesis (machine-confirmed:
aggregate build green on 3 November 2026, standard axioms). $G_2$, $F_4$, $E_6$, $E_7$, $E_8$ rest on three inputs. Two of
them (T and R) are published theorems for every type, the first with one proof step flagged by the project; the library
has formalized R only for $A_2$ and $D_4$, and T only for $A_2$. The third (S), needed for $G_2$, $F_4$, $E_6$, $E_7$,
$E_8$, has no published proof known to us. \textbf{These five final theorems are therefore conditional on a conjecture},
supported by numerical evidence only.

\paragraph{The three inputs, by type.} The hypotheses of the final theorems are of three kinds, which we call
Input~T, Input~R and Input~S throughout (they are defined in Section~\ref{l27:sec:inputs}).
\begin{itemize}
\item \textbf{Input T} (the \emph{total-dimension bound}): every module with the relations of the fusion module
  $\ltwentysevenF_m$ has dimension at most $(\dim\ltwentyseveng)^m$. Lean: \lean{FusionTotalDimG}, and for $A_2$
  \lean{SectorWeyl.FusionTotalDim}.
\item \textbf{Input R} (the \emph{Demazure realizations}): for every $j$, some module with the relations of the
  Demazure module $D(1,j\theta)$ is, in every degree, at least as large as the Demazure character. Lean:
  \lean{DemazureRealizationG}.
\item \textbf{Input S} (the \emph{star bound}): the Demazure sum is bounded by the Arakawa--Kawasetsu coefficients,
  $\mathcal D_N\le c_N$. Lean: \lean{XStarBound}.
\end{itemize}
\begin{center}\small
\begin{tabular}{p{0.07\textwidth}p{0.25\textwidth}p{0.25\textwidth}p{0.25\textwidth}}
\toprule
type & Input T & Input R & Input S\\
\midrule
$A_2$ & proved (\lean{CVUniv.fusionTotalDim}) & proved (\lean{EvalModel.}\allowbreak\lean{demazureRealization_all}) & proved (\lean{A2.starIdentity})\\
$D_4$ & hypothesis & proved (\lean{D4Char.}\allowbreak\lean{d4_demazureRealization}) & proved (\lean{D4Wp.d4SchurEqE4} with the unclamped Demazure $=$ Schur identity; item (vi))\\
$G_2$, $F_4$ & hypothesis & hypothesis & hypothesis\\
$E_6$, $E_7$, $E_8$ & hypothesis & hypothesis & hypothesis\\
\bottomrule
\end{tabular}
\par\smallskip\footnotesize ``Proved'' for $A_2$: machine-confirmed; aggregate build green on 2026-10-28, standard axioms (Section~\ref{l27:sec:verification}).
``Proved'' for Input R of $D_4$: machine-confirmed; aggregate build green on 2026-11-03, standard axioms.
``Hypothesis'' means \emph{not proved in Lean}; it says nothing about the literature. In the literature, T and R are
theorems for every type and S is conjectural for $G_2$, $F_4$, $E_6$, $E_7$, $E_8$ (paragraph below, and
Table~\ref{l27:tab:precise}).
\end{center}
Everything else in the argument is proved for all seven types (Chapters~\ref{ch:L24}--\ref{ch:L26}). The long form of
this table, with sources and numerical evidence, is Section~\ref{l27:sec:ledger}.

\paragraph{The literature behind the inputs, stated once.} We state the literature position here and refer to this
paragraph elsewhere. A citation never discharges a Lean hypothesis; it only bears on whether the hypothesis is believed
true.
\begin{enumerate}
\item[(i)] \textbf{Input T is a published theorem for every simple $\ltwentyseveng$, with one flagged proof step.} The
  universal module with the Lean relations is the \emph{truncated Weyl module} $W(m\theta,m)$ of Kus and Littelmann (their
  Definition~4.1): this is an elementary identification, proved on paper in Lemma~\ref{l27:lem:truncweyl} below. Kus and
  Littelmann, Pacific J.\ Math.\ 278 (2015) 427--445, Theorem~4.2 (proved through their Theorem~3.5 and
  Corollary~3.6(3)), taken at $\lambda=\theta$, $\lambda_0=0$, $N=m$, gives $W(m\theta,m)\cong\ltwentyseveng^{*m}$, of dimension
  $(\dim\ltwentyseveng)^m$. For $\lambda_0=0$ the paper places \emph{no restriction on $\ltwentyseveng$} (the restriction to
  classical types and $G_2$ concerns only $\lambda_0\neq0$), and $\theta$ satisfies their hypothesis $\lambda\in L^+$ in every
  type, since it is supported on long simple roots (or is $2\omega_1$ in type $C$). So the specialization is the same for
  all six types. \emph{Caveat:} the written argument on p.~437 of Kus--Littelmann, which passes from ``the difference lies in
  the left ideal generated by the relations'' to ``its symbol kills the fusion vector'', controls only the initial forms of
  the generators of the ideal, not of all its elements. The project's source audit flags this step as unjustified as written
  (Section~\ref{l27:sec:feigin}). The caveat applies uniformly to all six types. For $D_4$ it has been sharpened by
  computation (Python, not Lean): at $m=2$ the step as written is \emph{insufficient}, since the symbols of the defining
  relations generate a strictly smaller left ideal than the symbols of all elements of the annihilator, with an explicit
  witness $\tilde a$ (end of Section~\ref{l27:sec:feigin}). This is a gap in the published proof, not a counterexample: the
  conclusion of the step was confirmed for $D_4$ at $m=2$ (exactly) and $m=3$ (modulo two primes), and Theorem~4.2 itself is not
  known to fail for any type.
\item[(ii)] \textbf{Feigin's earlier proof.} The statement first appears as E.~Feigin, \emph{SIGMA} 4 (2008) 070,
  arXiv:0806.4851, Proposition~1 (the bigraded isomorphism) and Corollary~2 ($\ltwentyseveng^{*N}\cong
  D(1,N\theta)/(\ltwentyseveng\otimes t^N\C[t])D(1,N\theta)$, for every simple $\ltwentyseveng$). In arXiv v1 the same results are
  numbered Proposition~2.1 and Corollary~2.2. The printed proof of Proposition~1 is labelled ``the idea of the proof''.
  Kus and Littelmann report a gap in it, by personal communication from the author, and the natural reading of the sketch is
  false: for two adjoint factors of $\mathfrak{sl}_3$ a filtration piece of dimension $55$ would have to split into pieces of
  total dimension $63$ (Section~\ref{l27:sec:feigin}).
\item[(iii)] \textbf{Ravinder is not independent evidence for T.} B.~Ravinder, \emph{SIGMA} 10 (2014) 110,
  arXiv:1409.0274, states his results for every simple $\ltwentyseveng$; only his Theorem~2 and Corollary~3 assume
  $\ltwentyseveng$ simply laced. His Corollary~2 at $a=0$ \emph{is} Feigin's Corollary~2, imported as his Proposition~5, so it
  inherits Feigin's gap.
\item[(iv)] \textbf{The $A_2$ route.} Setia, Rani and Khandai (arXiv:2405.11944, \S8) give an independent proof for
  $\mathfrak{sl}_3$ through Chari--Venkatesh modules. The library's proof of Input T for $A_2$ follows this route
  (Chapter~\ref{ch:L26}, Section~\ref{l26:sec:cv}), so it depends neither on Feigin nor on the flagged step of
  Kus--Littelmann.
\item[(v)] \textbf{Input R is a published theorem for every type; it is formalized for $A_2$ and $D_4$.} Two facts are needed.
  The relation half, that the extremal vector of $D(1,j\theta)$ satisfies the Lean relations, follows from Ravinder's
  Proposition~1 (a presentation of $D(1,j\theta)$, stated for every simple $\ltwentyseveng$, after Chari--Venkatesh,
  Corollary~3.5) by the short argument of Lemma~\ref{l27:lem:truncweyl}; it includes the short roots of $G_2$ and $F_4$. The
  graded half, that the graded dimension of $D(1,j\theta)$ is the output of the Lean algorithm, is the Demazure character
  formula (Kumar, Invent.\ Math.\ 89 (1987); Mathieu, Ast\'erisque 159--160 (1988)), valid for all Kac--Moody algebras. For
  $D_4$, $E_6$, $E_7$, $E_8$ the Frenkel--Kac construction also realizes $D(1,j\theta)$ inside the lattice VOA. For $D_4$
  the library proves Input R in every degree by its own route (\lean{D4Char.d4_demazureRealization},
  Theorem~\ref{l27:thm:d4R}). For $G_2$, $F_4$, $E_6$, $E_7$, $E_8$ the library lacks a Lean model of $D(1,j\theta)$ with
  its graded dimension, and a proof that the Lean algorithm implements the Demazure character formula.
\item[(vi)] \textbf{Input S is a conjecture for $G_2$, $F_4$, $E_6$, $E_7$, $E_8$.} No published source is known to us. The
  evidence is numerical only: equality $\mathcal D_N=c_N$ is checked through $q^{60}$ for $G_2$, $q^{25}$ for $F_4$, $q^{20}$
  for $E_6$, $q^{16}$ for $E_7$ and $q^{14}$ for $E_8$ (scoping pass of 28 October 2026; Section~\ref{l27:sec:starinput}). For $G_2$ the library reduces it to the
  single identity \lean{G2DemazureEqKW}, the other half (\lean{G2KWEqMLDE}) being proved. $D_4$ has no Input S: the needed
  identity is proved as a pair. \lean{D4Wp.d4SchurEqE4} is the Schur $=$ $E_4$ half (the series \lean{d4SchurQ} is the MLDE
  solution); the Demazure $=$ Schur half is \lean{d4DemazureEqSchur_of_positivity} (from JTP's finite identity
  \lean{d4FiniteIdentity_holds}), whose positivity hypothesis is discharged on the unclamped route by
  \lean{D4Wp.d4DemazurePositivity_of_U5Z}, i.e.\ from U5 and hence from Inputs T and R, of which R is proved for $D_4$
  (Theorem~\ref{l25:thm:d4}; dictionary of Definition~\ref{l27:def:realization}).
\end{enumerate}
Table~\ref{l27:tab:precise} in Section~\ref{l27:sec:ledger} states each input type by type, with its precise literature
status and the library's progress toward it.

\paragraph{Plan of the chapter.}
Section~\ref{l27:sec:chain} recalls what classical freeness is and how the proof is organised.
Section~\ref{l27:sec:fusionprimer} is a primer on the current-algebra modules in which the inputs are stated, with a
worked $\mathfrak{sl}_3$ example. Section~\ref{l27:sec:conclusion} states the common conclusion and lists the seven
VOAs. Section~\ref{l27:sec:inputs} states Inputs T, R and S. Section~\ref{l27:sec:finals} states the seven final
theorems, starting with the unconditional $A_2$ theorem. Section~\ref{l27:sec:ledger} is the type-by-type ledger, and
Section~\ref{l27:sec:cores} says what is known about the inputs that remain open. Section~\ref{l27:sec:ladders}
presents the per-type headline theorems from which the final theorems are assembled. Section~\ref{l27:sec:deadends}
records approaches that are known to fail. The chapter ends with the verification record, remarks on the Lean sources,
the referee's checklist and the references.

\paragraph{Prerequisites.} From Chapter~\ref{ch:L24}: the predicate \lean{CaseClassicallyFree}
(Definition~\ref{l24:def:packages}), the seven explicit presentations (Theorem~\ref{l24:thm:sevenpres}), sector data
(Definition~\ref{l24:def:sectordata}), current modules and the input U5 (Definition~\ref{l24:def:U5}), and the logical
arc of Remark~\ref{l24:rem:arc}. From Chapter~\ref{ch:L25}: the Demazure sum (Definition~\ref{l25:def:dsum}), the
coefficients $c_N$ (Definition~\ref{l25:def:mlde}), the Demazure algorithm (Definition~\ref{l25:def:demazure}) and the
target lower bound $c_N\le\dim L_N$ (Theorem~\ref{l25:thm:alltypes}). From Chapter~\ref{ch:L26}: the reduction of U5
and U6 to Inputs T and R (universal modules, Section~\ref{l26:sec:univ}; the kernel lemma, Section~\ref{l26:sec:kernelG};
the squeeze, Section~\ref{l26:sec:squeeze}; U6 by construction, Section~\ref{l26:sec:u6}) and the proofs of Inputs T and
R for $A_2$ (Sections~\ref{l26:sec:u7} and~\ref{l26:sec:cv}). Each object is recalled where it first appears here.

\subsection{Conventions}

\begin{itemize}
\item \emph{Omitted proofs.} A Lean quotation shows a statement, not its proof. A block that ends in \lean{:=} or
  \lean{:= by} with nothing after it is a statement whose proof is omitted. A one-line proof term is shown when it is
  informative, for example when it is a composition of earlier theorems.
\item \emph{``Axioms standard''} means that \texttt{\#print axioms} reports a subset of
  \lean{[propext, Classical.choice, Quot.sound]}, the three standard axioms of Lean's logic (Chapter~\ref{ch:01},
  Section~\ref{ch01:sec:trust}). The library has no project axiom (\lean{VOA.pbw} is a theorem, Chapter~\ref{ch:L02}).
  ``Axioms standard'' does \emph{not} mean ``unconditional'': \texttt{\#print axioms} does not list the explicit
  hypotheses of a theorem (Section~\ref{ch01:sec:uncond}). Every status box therefore lists the hypotheses.
\item \emph{Where a status claim comes from.} Unless a status box says otherwise, an axiom or build claim is a
  \emph{campaign record} (the project's builds and the \texttt{\#print axioms} lines at the end of each Lean file),
  described in Section~\ref{l27:sec:verification}. ``Re-checked'' marks a claim the author of this chapter re-ran.
\end{itemize}

\subsection{Reading the Lean in this chapter}\label{l27:sec:readlean}

Chapter~\ref{ch:01}, Section~\ref{ch01:sec:lean}, is the general primer; Chapter~\ref{ch:L26},
Section~\ref{l26:sec:readlean}, covers the current-module idioms. The constructs that recur here are these.
\begin{itemize}
\item \textbf{Fully qualified names.} The final theorems are stated with full names such as
  \lean{VOA.DeligneAudit.A2.W}; the file \leanpath{VOA/DeligneAudit/DeligneSeries.lean} has no \lean{variable}
  declarations, so there are no hidden section hypotheses. Elsewhere short names such as \lean{W} or
  \lean{explicitPresentation} resolve through the file's \lean{namespace} and \lean{open} lines; the dictionaries give the
  full name whenever a short name is ambiguous (\lean{explicitPresentation}, \lean{classicallyFree_of_fusion} and
  \lean{mldeCoeff} each denote several declarations).
\item \textbf{\lean{.A2}, \lean{.D4}, \dots} are the constructors of the enumeration \lean{HigherRankDeligneCase}; Lean
  infers the type from the position, so \lean{CaseClassicallyFree .A2 W P} reads ``case $A_2$, VOA $W$, presentation $P$''.
\item \textbf{Named statements.} \lean{def P : Prop := …}, \lean{abbrev P : Prop := …} and \lean{structure P : Prop where …}
  only \emph{name} a statement; nothing is asserted. A theorem with an argument \lean{(h : P)} is proved assuming
  \lean{P}. A named statement becomes a theorem when the library proves a term of it, as
  \lean{CVUniv.fusionTotalDim : SectorWeyl.FusionTotalDim} does. \lean{abbrev} is a definition that Lean unfolds
  automatically.
\item \textbf{Existentials over modules.} \lean{∃ (V : Type) (_ : AddCommGroup V) (_ : Module ℂ V) (ρ : …) (w : V), …} asks
  for a vector space $V$ (the two anonymous arguments \lean{_} are its additive group and $\C$-vector-space structures), a
  current-module structure \lean{ρ} on it and a vector \lean{w}.
\item \textbf{Current modules.} \lean{CurrentModule f V} is a family of linear maps \lean{ρ.op x r}$=x(r)=x\otimes t^r$ on $V$
  satisfying the current-algebra bracket for the structure constants \lean{f}; \lean{wordOp ρ l w} applies a word;
  \lean{cycDeg ρ w d} is the degree-$d$ part of the cyclic span of $w$ (Chapter~\ref{ch:L24}, Definition~\ref{l24:def:U5}).
\item \textbf{Sums and series.} \lean{∑ b ∈ Finset.range B, f b} is $\sum_{b<B}f(b)$; \lean{PowerSeries ℚ} is $\Q[[q]]$;
  \lean{coeff N F} is $[q^N]F$; \lean{Module.finrank ℂ U} is $\dim U$ (with the convention that it is $0$ when $U$ is
  infinite-dimensional, which is why the finiteness of the relevant spaces is always proved separately).
\item \textbf{Kernels and ideals.} \lean{MvPolynomial.aeval x} is the algebra map $\C[X_0,\dots,X_{n-1}]\to A$ sending $X_i$
  to $x_i$; \lean{RingHom.ker (…).toRingHom} is its kernel, an ideal; \lean{Ideal.span} is the ideal generated by a set.
\end{itemize}

\subsection{Notation}

We follow Chapters~\ref{ch:L24}--\ref{ch:L26}. Several letters are overloaded in the sources; in the text of this
chapter we use the following, and keep the Lean names inside quotations and dictionaries.
\begin{center}\small
\begin{longtable}{p{0.22\textwidth}p{0.70\textwidth}}
\toprule
symbol & meaning\\
\midrule\endhead
$\ltwentyseveng$, $\ltwentysevendg$, $h^\vee$, $k$ & a simple Lie algebra of the series, its dimension
  $\ltwentysevendg=\dim\ltwentyseveng$ (Lean: \lean{S.n}, and in some displays of earlier chapters $n$), its dual Coxeter
  number, the level $k=-1-h^\vee/6$.\\
$\theta$; $e_\theta$, $f_\theta$ & the highest root; its root vectors. For $\mathfrak{sl}_3$ the simple lowering vectors are
  $f_1,f_2$ and $f_\theta=f_{12}$.\\
$\ltwentyseveng[t]$, $x(r)$ & the current algebra $\ltwentyseveng\otimes\C[t]$; $x(r)=x\otimes t^r$, $r\ge0$. The power of $t$
  is the \emph{degree}, recorded by the variable $q$.\\
$L_k(\ltwentyseveng)$, $L_N$, $W$ & the simple affine VOA; its weight-$N$ space; $W$ is the Lean name of the VOA.\\
$R_W$, $\ltwentysevenIns$ & Zhu's $C_2$-algebra; the explicit space of quadrics (Chapter~\ref{ch:L24}).\\
$[q^N]f$, $(q)_j$ & the coefficient of $q^N$; $(q)_j=\prod_{i=1}^j(1-q^i)$.\\
$V(\lambda)$ & the irreducible $\ltwentyseveng$-module of highest weight $\lambda$; $V(\theta)=\ltwentyseveng$.\\
$D(1,\lambda)$; $\mathrm{Dem}_j$ & the level-one Demazure module (Section~\ref{l27:sec:fusionprimer}); its graded
  character for $\lambda=j\theta$. The Demazure index is $j$ (Lean writes \lean{n}).\\
$\mathcal D_N$ & the Demazure sum $[q^N]\sum_{j\ge0}q^j\mathrm{Dem}_j/(q)_j$.\\
$c_N$ & the Arakawa--Kawasetsu coefficient, Lean \lean{DeligneCharacter.mldeCoeff h N}.\\
$M_{a,b}$; $\ltwentysevenF_m$ & Ravinder's modules; the fusion module $\ltwentysevenF_m=M_{0,m}$.\\
$m$ & the number of adjoint factors in a fusion product; the modes $x(r)$ with $r\ge m$ are the ones cut off in
  $\ltwentysevenF_m$.\\
$\mathcal I_m$ & the ideal $\ltwentyseveng\otimes t^m\C[t]$ of $\ltwentyseveng[t]$ (not to be confused with
  $\ltwentysevenIns$).\\
$\Phi_d$ & the degree filtration of a fusion product.\\
$u$ & the number of trivial tensor factors (Section~\ref{l27:sec:feigin}).\\
\bottomrule
\end{longtable}
\end{center}

%======================================================================
\section{Classical freeness and the shape of the proof}\label{l27:sec:chain}
%======================================================================

\paragraph{Classical freeness in five sentences.} For a VOA $W$, Zhu's $C_2$-algebra is $R_W=W/C_2(W)$, where $C_2(W)$
is spanned by the states $a_{(-2)}b$; it is a commutative algebra, the ``classical limit'' of $W$ (for a 4d
$\mathcal N=2$ theory, the Higgs branch chiral ring). Li's filtration $F^0W\supseteq F^1W\supseteq\dots$ makes the
associated graded $\ltwentysevengr^{\mathrm{Li}}W$ a commutative algebra with a derivation, and its degree-zero part is
$R_W$ (Chapters~\ref{ch:L11} and~\ref{ch:L24}). The \emph{arc algebra} $\ltwentysevenJ R_W$ is the universal commutative
algebra with a derivation generated by $R_W$: if $R_W=\C[X]/I$, it is $\C[X^{(0)},X^{(1)},\dots]$ modulo all derivatives of
all elements of $I$. By the universal property there is a canonical surjection $\ltwentysevenJ R_W\to\ltwentysevengr^{\mathrm{Li}}W$;
$W$ is \emph{classically free} when it is injective, so that every state of $W$ is a classical polynomial in the generators
and their derivatives, subject only to the derivatives of the relations of $R_W$. Counting by weight, classical freeness is
the equality $\dim W_N=\dim(\ltwentysevenJ R_W)_N$ for every $N$.

\paragraph{The carrier.} For an affine VOA the currents $x_{(-1)}\ltwentysevenvac$ generate, so
$R_W=\operatorname{Sym}\ltwentyseveng/\ker$. The library proves that the explicit quadrics $\ltwentysevenIns$ lie in that
kernel (Chapter~\ref{ch:L24}), so $\ltwentysevenJ R_W$ is a quotient of
$\ltwentysevenJ\bigl(\operatorname{Sym}\ltwentyseveng/(\ltwentysevenIns)\bigr)$. The latter has an explicit linear model, the
\emph{carrier}, a space of polynomial functions in which the weight-$N$ part is finite-dimensional; write
$\mathrm{carrierDim}(N)$ for its dimension. Since Li's map is onto, $\dim L_N\le\mathrm{carrierDim}(N)$ always. Classical
freeness, together with $\ker=(\ltwentysevenIns)$, follows from the reverse inequality.

\begin{remark}[The logical arc]\label{l27:rem:arc}
Chapter~\ref{ch:L24} (Remark~\ref{l24:rem:arc}) reduces classical freeness to three statements, for every weight $N$:
\begin{align*}
\text{(A)}\quad & \bigl(\mathrm{carrierDim}(N)\le\dim L_N\ \text{for all }N\bigr)\ \Longrightarrow\ \text{classical freeness and}
\ \ker=(\ltwentysevenIns);\\
\text{(B)}\quad & \mathrm{carrierDim}(N)\le\mathcal D_N;\\
\text{(C)}\quad & \mathcal D_N\le c_N\le\dim L_N .
\end{align*}
(A) is proved with no hypothesis (Chapter~\ref{ch:L24}). (C): $c_N\le\dim L_N$ is proved for every type
(Chapter~\ref{ch:L25}); $\mathcal D_N\le c_N$ is Input S. (B) is the \emph{carrier side}. Chapter~\ref{ch:L24} proves it
from two statements about current-algebra modules:
\begin{itemize}
\item \textbf{U5} (fusion module bound): there is a family of power series $F_m(q)$ such that every module with the
  relations of the fusion module $\ltwentysevenF_m$ has degree-$d$ piece of dimension at most $[q^d]F_m$, and total dimension at most
  $\ltwentysevendg^{\,m}$;
\item \textbf{U6} (mixed recursion): the family $M_{a,b}(q)$ generated from $F$ by $M_{0,b}=F_b$ and
  $M_{a+1,b}=M_{a,b+1}+q^{2a+b+1}M_{a,b}$ satisfies $M_{j,0}=\mathrm{Dem}_j$ for every $j$.
\end{itemize}
The carrier side also uses a \emph{divisor comparison}, a statement of commutative algebra about the carrier of each
type (a containment of two explicit modules over the symmetric polynomials, after localising at the primes that contain a
discriminant), proved for each type by integer \emph{certificates}: finite lists of structure-constant identities, each
checked by Lean's kernel, with Lean-proved lemmas saying that a passing check implies the statement
(Chapter~\ref{ch:L24}, Section~\ref{l24:sec:divisor}). Chapter~\ref{ch:L26} proves U6 by \emph{defining} $F$ as the solution
of the recursion, and reduces U5 to Inputs T and R. So the inputs that remain are exactly T, R and S, as in the table
above (for $D_4$, only T).

When all inputs hold the chain closes:
$\dim L_N\le\mathrm{carrierDim}(N)\le\mathcal D_N\le c_N\le\dim L_N$, so all four numbers are equal. In particular
$\dim L_N=c_N$, a \emph{consequence} of the inputs. This equality is not itself stated as a Lean theorem for the six
conditional types; for $A_2$ it is a Lean theorem proved independently (\lean{A2.finrank_eq_mldeCoeff'},
Chapter~\ref{ch:L25}).
\end{remark}

%======================================================================
\section{Fusion products and Demazure modules: a primer}\label{l27:sec:fusionprimer}
%======================================================================

Inputs T and R are statements about graded modules for the current algebra $\ltwentyseveng[t]$. This section defines
the modules involved and computes the smallest nontrivial case for $\mathfrak{sl}_3$.

\paragraph{Graded cyclic modules and presentations.} $\ltwentyseveng[t]$ is graded by the power of $t$. A graded module
generated by one vector $w$ of degree $0$ is \emph{cyclic}; its degree-$d$ piece is spanned by the words
$x_1(r_1)\cdots x_s(r_s)w$ with $r_1+\dots+r_s=d$, and its \emph{graded character} is $\sum_d\dim M_d\,q^d$. A cyclic module is
often given by a \emph{presentation}: a list of words that kill $w$. Every module whose generator satisfies the same relations
is a quotient of the \emph{universal} one, so its graded character is coefficientwise smaller. Consequently an \emph{upper}
bound for a presentation must be proved for the universal module, i.e.\ ``for every module with the relations'', while a
\emph{lower} bound needs only ``some module with the relations''.

\paragraph{Local Weyl modules.} For a dominant weight $\lambda$, the local Weyl module $W_{\mathrm{loc}}(\lambda)$ of Chari
and Pressley is the cyclic module with relations
\begin{equation*}
\mathfrak n^+[t]\,w=0,\qquad h(s)\,w=\delta_{s,0}\,\lambda(h)\,w,\qquad f_i(0)^{\lambda(h_i)+1}w=0 .
\end{equation*}
It is finite-dimensional.

\paragraph{Evaluation modules and fusion products.} A $\ltwentyseveng$-module $V$ and a point $c\in\C$ give the
\emph{evaluation module} $V(c)$, on which $x(r)$ acts by $c^r x$. Given cyclic $\ltwentyseveng$-modules $V_1,\dots,V_m$
and distinct points $c_1,\dots,c_m$, the tensor product $V_1(c_1)\otimes\dots\otimes V_m(c_m)$ is cyclic on the tensor
product of the highest-weight vectors; filtering it by degree in $t$ and taking the associated graded gives a graded
$\ltwentyseveng[t]$-module $V_1*\dots*V_m$ of the same total dimension, the \emph{fusion product} of Feigin and Loktev. (For
the modules used here the result is known not to depend on the points, by Ravinder's presentation theorem; independence of
the points is not known in general and is not needed.) The basic case is $\ltwentyseveng^{*m}=V(\theta)^{*m}$, of total
dimension $\ltwentysevendg^{\,m}$.

\paragraph{Level-one Demazure modules.} Let $L(\Lambda_0)$ be the basic (level-one vacuum) representation of the affine
Lie algebra $\widehat{\ltwentyseveng}$. For a dominant $\lambda$ in the coroot lattice the extremal vector of
$L(\Lambda_0)$ of weight $\lambda$ generates, under the Borel subalgebra containing $\ltwentyseveng[t]$, a
finite-dimensional graded $\ltwentyseveng[t]$-module $D(1,\lambda)$, the \emph{level-one Demazure module}. Its graded
character is given by the Demazure character formula: apply the Demazure operators along a reduced word of the
translation by $\lambda$ to the extremal weight. Chapter~\ref{ch:L25} implements this as an algorithm
(Definition~\ref{l25:def:demazure}); we write $\mathrm{Dem}_j$ for the result with $\lambda=j\theta$. The smallest case is
\begin{gather*}
D(1,\theta)=\ltwentyseveng\oplus\C,\qquad x(0)\ \text{acts on }\ltwentyseveng\text{ by the adjoint action and kills }\C,\\
x(1)\,y=\kappa(x,y)\cdot1\ (y\in\ltwentyseveng),\qquad x(1)\,1=0,\qquad x(r)=0\ (r\ge2),
\end{gather*}
with $\kappa$ the normalised invariant form; so $\mathrm{Dem}_1=\ltwentysevendg+q$. More generally
$D(1,j\theta)\cong D(1,\theta)^{*j}$, of total dimension $(\ltwentysevendg+1)^j$ (Fourier--Littelmann 2007 for simply laced
$\ltwentyseveng$; for all $\ltwentyseveng$ Kus and Littelmann cite Chari--Shereen--Venkatesh--Wand, arXiv:1408.4090, together
with Fourier--Littelmann. We did not verify the latter attribution; nothing in this chapter depends on it). On an evaluation
module $D(1,\theta)(c)$ at a point $c$, the action is $x(r)\,v=c^r\,x(0)v+r\,c^{r-1}\,x(1)v$: the value of $t^r$ at $c$ acts
through the $\ltwentyseveng$-action, and its derivative through the map $\ltwentyseveng\to\C$.

\paragraph{Ravinder's modules $M_{a,b}$.} Fix the highest root vector $f_\theta$ and a \emph{fusion set} of negative root
vectors containing $f_\theta$ (Ravinder takes $\{f_\alpha:\langle\alpha,\theta^\vee\rangle=1\}\cup\{f_\theta\}$, where
$\langle\alpha,\theta^\vee\rangle=2(\alpha,\theta)/(\theta,\theta)$; with long roots of squared length $2$ this is $(\alpha,\theta)=1$.
Since $\theta$ is long, every positive root $\alpha\neq\theta$ has $\langle\alpha,\theta^\vee\rangle\in\{0,1\}$, and the
condition includes the \emph{short} roots of $G_2$ and $F_4$ with $(\alpha,\theta)=1$. For a short root of $G_2$ the other
pairing $(\theta,\alpha^\vee)$ equals $3$, so writing the condition as $(\theta,\alpha^\vee)=1$ would wrongly exclude
them). For $a,b\ge0$, $M_{a,b}$
is the cyclic module with generator $w$ of weight $(a+b)\theta$ and relations
\begin{gather*}
e_\alpha(r)\,w=0\ (\alpha>0,\ r\ge0),\qquad h(r)\,w=0\ (r\ge1),\qquad h(0)\,w=(a+b)\,\theta(h)\,w,\\
f_i(0)^{(a+b)\theta(h_i)+1}w=0,\\
f_\alpha(a+b)\,w=0\ \ (f_\alpha\text{ in the fusion set},\ \alpha\ne\theta),\qquad f_\theta(2a+b)\,w=0 .
\end{gather*}
The first two lines are the local Weyl module $W_{\mathrm{loc}}((a+b)\theta)$; the third cuts the modes. By \emph{Garland's
lemma} (a consequence of the integrability of each $\mathfrak{sl}_2$-triple: for $\mathfrak{sl}_2$, $e(1)^sf(0)^{r+s}w$ is a
nonzero multiple of an explicit polynomial in the $f(k)$ applied to $w$) the cut on one mode implies cuts on all higher modes.
Two cases matter: $M_{0,m}$, with all modes cut at $t^m$, is the \emph{fusion module} $\ltwentysevenF_m$, and the fusion
product $\ltwentyseveng^{*m}$ is a quotient of it, so $\dim\ltwentysevenF_m\ge\ltwentysevendg^{\,m}$; and $M_{j,0}$ is a
presentation of $D(1,j\theta)$ (Ravinder's Proposition~1). Ravinder's theorem (his Theorem~1 and Corollaries~1--2,
stated for every simple $\ltwentyseveng$; the case $a=0$ is imported from Feigin, see item~(iii) of the literature paragraph)
is that $M_{a,b}\cong D(1,\theta)^{*a}*\ltwentyseveng^{*b}$, of dimension $(\ltwentysevendg+1)^a\ltwentysevendg^{\,b}$, and
that the family is linked by short exact sequences
\begin{equation}\label{l27:eq:ravinder}
0\longrightarrow \tau_{2a+b-1}\,M_{a-1,b}\longrightarrow M_{a,b}\longrightarrow M_{a-1,b+1}\longrightarrow0
\qquad(a\ge1),
\end{equation}
where $\tau_s$ raises degrees by $s$ and the left map sends $w$ to $f_\theta(2a+b-1)w$. Taking characters gives the mixed
recursion U6, $\mathrm{ch}\,M_{a,b}=\mathrm{ch}\,M_{a-1,b+1}+q^{2a+b-1}\mathrm{ch}\,M_{a-1,b}$, which determines
$\mathrm{ch}\,M_{0,m}$ from the Demazure characters $\mathrm{ch}\,M_{j,0}=\mathrm{Dem}_j$.

\begin{example}[$\mathfrak{sl}_3$ with two adjoint factors]\label{l27:ex:sl3}
Take $\ltwentyseveng=\mathfrak{sl}_3$, $\ltwentysevendg=8$. Then $\mathrm{ch}\,M_{0,1}=8$ and $\mathrm{ch}\,M_{1,0}=\mathrm{Dem}_1=8+q$,
and the recursion with $(a,b)=(1,0)$ reads $8+q=8+q\cdot1$ (totals $9=8+1$). The Demazure algorithm gives
$\mathrm{Dem}_2=27+28q+17q^2+8q^3+q^4$, of total $81=9^2$. Running the recursion downwards:
\begin{align*}
\mathrm{ch}\,M_{1,1}&=\mathrm{ch}\,M_{2,0}-q^3\,\mathrm{ch}\,M_{1,0}=27+28q+17q^2 &&(81-9=72=9\cdot8),\\
\mathrm{ch}\,M_{0,2}&=\mathrm{ch}\,M_{1,1}-q^2\,\mathrm{ch}\,M_{0,1}=27+28q+9q^2 &&(72-8=64=8^2).
\end{align*}
The degree-zero piece $27$ is $V(2\theta)$, and $27+28q+9q^2$ is the graded character of the fusion product $\mathfrak{sl}_3*\mathfrak{sl}_3$
(Chapter~\ref{ch:L26}, Example~\ref{l26:ex:fus88}).

\emph{The squeeze at $m=2$.} Suppose we know: (R) a module with the relations of $M_{2,0}$ whose character is at least
$\mathrm{Dem}_2$ in every degree (Input R); (K) the \emph{spanning half} of each exact sequence,
$\mathrm{ch}\,M_{a,b}\le\mathrm{ch}\,M_{a-1,b+1}+q^{2a+b-1}\mathrm{ch}\,M_{a-1,b}$ coefficientwise (this is the elementary half:
see the warning below); and (T) $\dim M_{0,2}\le64$ (Input T). Applying (K) twice and the known $\mathrm{ch}\,M_{0,1}=8$,
$\mathrm{ch}\,M_{1,0}\le8+q$,
\begin{equation*}
\mathrm{Dem}_2\ \le\ \mathrm{ch}\,M_{2,0}\ \le\ \mathrm{ch}\,M_{0,2}+8q^2+q^3(8+q),
\qquad\text{so}\qquad \mathrm{ch}\,M_{0,2}\ \ge\ 27+28q+9q^2
\end{equation*}
coefficientwise. The right-hand side has total $64$, and by (T) the total of $\mathrm{ch}\,M_{0,2}$ is at most $64$; a
coefficientwise lower bound with the same total forces equality in every degree. So $\mathrm{ch}\,M_{0,2}=27+28q+9q^2$: the
graded upper bound U5 for $m=2$ follows from a lower bound (R), an elementary spanning inequality (K) and a total upper bound
(T). This is Ravinder's squeeze (Chapter~\ref{ch:L26}, Theorem~\ref{l26:thm:squeeze}) in its smallest case; note that the
subtractive step works only because the total bound (T) closes it.
\end{example}

\begin{warning}[title={The surjection half is free; the injection half is not}]
In (\ref{l27:eq:ravinder}) the right map is surjective for elementary reasons: the relations of $M_{a-1,b+1}$ are those of
$M_{a,b}$ plus the one relation $f_\theta(2a+b-1)w=0$. Its kernel is the cyclic span of $f_\theta(2a+b-1)w$, which is a quotient of
$\tau_{2a+b-1}M_{a-1,b}$ (Ravinder's kernel lemma: that vector satisfies the relations of $M_{a-1,b}$). These facts give
\emph{upper} bounds, inequality (K) above. What is hard is the other direction: that the kernel is \emph{as large as}
$\tau_{2a+b-1}M_{a-1,b}$, or that $D(1,j\theta)$ is \emph{as large as} its character formula says. Upper bounds come from spanning
sets; lower bounds need a model in which vectors can be shown to be independent.

This explains why Input T, an \emph{upper} bound on $\dim M_{0,m}$, is hard. A spanning set for $M_{0,m}$ obtained from the
relations is far larger than $\ltwentysevendg^{\,m}$ (the naive count gives $72$ instead of $64$ already for
$\mathfrak{sl}_3$, $m=2$; Section~\ref{l27:sec:deadends}). Counting through the exact sequences converts the missing upper bound
for $M_{0,m}$ into a missing \emph{lower} bound for the kernels, and a lower bound can in principle be proved by exhibiting
independent vectors in a model. For $A_2$ the library proves Input T by enlarging the family of modules until an induction
of this kind closes (Chapter~\ref{ch:L26}, Section~\ref{l26:sec:cv}).
\end{warning}

%======================================================================
\section{The common conclusion and the seven VOAs}\label{l27:sec:conclusion}
%======================================================================

Every final theorem has the same conclusion, a conjunction of two statements about one VOA $W$ and one explicit
presentation $P$.

\begin{definition}[Classical freeness of a case, and the kernel statement]\label{l27:def:conclusion}
Let $W$ be a VOA with weight-one states $\mathrm{gen}_i$, $i<n$, and let $I\subseteq\C[X_0,\dots,X_{n-1}]$ be an ideal
contained in the kernel of the map $X_i\mapsto[\mathrm{gen}_i]$ into Li's associated graded algebra (a \emph{current
$C_2$ presentation} $P$). The case is \emph{classically free} for $P$ if the infinite Li jet map
$\ltwentysevenJ(\C[X]/I)\to\ltwentysevengr^{\mathrm{Li}}W$ is bijective. The \emph{kernel statement} is
$\ker\bigl(\C[X]\to\ltwentysevengr^{\mathrm{Li}}W\bigr)=I$.
\end{definition}
\noindent File \leanpath{VOA/\allowbreak Examples/\allowbreak DeligneHigherRankConditionalCapstone.lean}.
\begin{leancode}
def CaseClassicallyFree (P : CurrentC2Presentation case W) : Prop :=
  LiPresentationClassicallyFree W
    (fun i => liDegreeZeroClass W 1 ⟨P.gen i, P.gen_weight_one i⟩)
    P.ideal P.relations_vanish
\end{leancode}

\begin{dictionary}
\lean{CurrentC2Presentation case W} is the structure of Definition~\ref{l24:def:packages}: the fields \lean{gen},
\lean{gen_weight_one}, \lean{ideal}, and \lean{relations}, which says that \lean{ideal} is \emph{contained} in the
kernel (\lean{P.relations_vanish} is a renaming of that field). \lean{liDegreeZeroClass W 1 x} is the class of the
weight-one state $x$ in $F^0W_1/F^1W_1$, the degree-zero part of Li's associated graded algebra.
\lean{LiPresentationClassicallyFree} is bijectivity of the infinite jet map (Definition~\ref{l24:def:lpcf}): its source is
the presented jet algebra, the polynomial ring in variables $X_i^{(r)}$ modulo all derivatives of all elements of the ideal,
and $X_i^{(r)}$ is sent to the $r$-th divided derivative of the class of $\mathrm{gen}_i$.
\lean{CaseClassicallyFree} has, from the \lean{variable} lines of its file, an implicit carrier \lean{\{V : Type uV\}} with
its vector-space structure and two explicit arguments \lean{case : HigherRankDeligneCase} (which of the seven types,
fixing the number of currents \lean{case.currentDimension}) and \lean{W : VertexOperatorAlgebra V}.

In every final theorem the kernel statement appears as a second conjunct,
\lean{RingHom.ker (MvPolynomial.aeval fun i => liDegreeZeroClass W 1 ⟨P.gen i, P.gen_weight_one i⟩).toRingHom = I}, an equality
of ideals of $\C[X_0,\dots,X_{n-1}]$. Since the image of $\C[X]$ is $R_W=\ltwentysevengr^{\mathrm{Li}}_0W$ (the currents
strongly generate $W$), this is the statement $R_W=\operatorname{Sym}\ltwentyseveng/I$.

\textbf{Why the conclusion is not weak.} Classical freeness of a presentation forces the ideal to equal the kernel
(Theorem~\ref{l24:thm:kereq}), so the second conjunct follows from the first; the library states both so that a reader
sees the presentation of $R_W$ directly. The conclusion cannot be satisfied by a degenerate choice: the ideal is fixed in
advance, and $W$ is a specific nonzero VOA.
\end{dictionary}

\begin{physics}
For the chiral algebra of a four-dimensional $\mathcal N=2$ theory, the first conjunct says that every Schur operator
is a classical object: a polynomial in the flavour currents and their derivatives, subject only to the derivatives of
the Higgs-branch chiral-ring relations. The second conjunct says that the chiral-ring relations among the moment maps
are exactly the quadrics $\ltwentysevenIns$. Together they say that the Schur index equals the Hilbert series of the arc
space of the \emph{scheme} $\operatorname{Spec}R_W$: no ``quantum'' cancellation between classical states.
\emph{Heuristic, not Lean:} the reduced variety underneath this scheme is the associated variety of $L_k(\ltwentyseveng)$,
which Arakawa and Moreau identify with the minimal nilpotent orbit closure $\overline{\mathbb O_{\min}}$ for the Deligne
series at these levels; the scheme itself is non-reduced (the Casimir is nilpotent in $R_W$ but not zero; \emph{checked outside Lean}: the Lean statement is \lean{sym_X_mul_casPoly}, that $X_a\cdot\mathrm{Cas}$ vanishes in the $C_2$-algebra).
\end{physics}

\begin{status}
Definition (Chapter~\ref{ch:L24}). The kernel statement is always a conclusion, never a hypothesis.
\end{status}

\paragraph{The seven VOAs.} Each final theorem is about a specific VOA built in Chapters~\ref{ch:L23}
and~\ref{ch:L24}. The table lists the Lean name, the Lie data from which it is built, the level, and the number of quadrics in
the library's relation list.
\begin{center}\small
\begin{longtable}{p{0.06\textwidth}p{0.30\textwidth}p{0.20\textwidth}p{0.06\textwidth}p{0.05\textwidth}p{0.05\textwidth}p{0.08\textwidth}}
\toprule
type & VOA (Lean) & Lie data & $k$ & $\ltwentysevendg$ & $h^\vee$ & quadrics\\
\midrule\endhead
$A_2$ & \lean{A2.W} & bespoke $\mathfrak{sl}_3$ & $-3/2$ & 8 & 3 & 8\\
$G_2$ & \lean{G2Lie.g2W} & \lean{g2CertData} & $-5/3$ & 14 & 4 & 27\\
$D_4$ & \lean{D4Lie.W4} & \lean{d4CertData} & $-2$ & 28 & 6 & 105\\
$F_4$ & \lean{F4Lie.f4W} & \lean{f4CertData} & $-5/2$ & 52 & 9 & 324\\
$E_6$ & \lean{DeligneCarrier.E6.W6} & \lean{e6LieData} & $-3$ & 78 & 12 & 650\\
$E_7$ & \lean{E7Lie.W7} & \lean{e7CertData} & $-4$ & 133 & 18 & 1539\\
$E_8$ & \lean{E8Lie.W8} & \lean{e8CertData} & $-6$ & 248 & 30 & 3875\\
\bottomrule
\end{longtable}
\end{center}
All seven are simple quotients of universal affine VOAs: $A_2$ through the $\mathfrak{sl}_3$ construction of
Chapter~\ref{ch:L24}, the other six as \lean{simpleVOA D k _} of Chapter~\ref{ch:L23}. Their central charges are
$c=k\,\ltwentysevendg/(k+h^\vee)=-2(h^\vee+1)$, namely $-8,-10,-14,-20,-26,-38,-62$.

\begin{warning}[title={What the relation lists are, and what Lean proves about them}]
For $A_2$ the ideal \lean{A2.I2ns} is generated by eight explicit quadrics $P_0,\dots,P_7$ (Chapter~\ref{ch:L24},
Definition~\ref{l24:def:A2rel}). Lean proves that they are homogeneous, linearly independent, and in the kernel; that they
span the adjoint summand $V(\theta)\subset\operatorname{Sym}^2\mathfrak{sl}_3$ was checked outside Lean and is \emph{not} a Lean
statement. The ideal does not contain the Casimir (only its products $X_a\cdot\mathrm{Cas}$); this non-membership is checked outside Lean and is not a Lean statement (the Lean statement is \lean{sym_X_mul_casPoly}). So the $A_2$ theorem identifies $R_V$
with $\operatorname{Sym}\mathfrak{sl}_3$ \emph{modulo these eight quadrics}; it does not mention, and must not be described as
identifying $R_V$ with, the coordinate ring of the minimal nilpotent orbit closure or a ``Joseph ideal'' (type $A$ has no Joseph
ideal in the sense of Chapter~\ref{ch:L24}, and $R_V$ is non-reduced, a fact checked outside Lean). For the other six types Lean proves that each listed quadric
lies in the kernel; that the lists ($27,105,324,650,1539,3875$ quadrics) span exactly $\ltwentysevenIns$ is checked outside Lean.
The conditional final theorems conclude $\ker=$ the ideal generated by the list.
\end{warning}

\begin{warning}[title={Scope: which VOA is meant for $G_2$, $D_4$, $F_4$, $E_6$, $E_7$, $E_8$}]\label{l27:warn:scope}
For $D_4$, $E_7$ and $E_8$ the VOA of the final theorem is \lean{simpleVOA} of the \emph{certificate-basis} Lie data
(\lean{d4CertData}, \lean{e7CertData}, \lean{e8CertData}): the lattice Lie algebra of Chapter~\ref{ch:L23}, written in
the basis in which the relation lists were computed, with the form Killing$/12$, $/36$, $/60$. The library also builds
VOAs \lean{d4SimpleVOA}, and so on, from a basis chosen by \lean{Module.finBasis}. Both are built from the same Lie
algebra, form and $h^\vee$, but from different \lean{LieData} (different bases), and the isomorphism of the resulting
VOAs (independence of the affine construction from the basis) is \emph{not} proved in Lean.

The same split exists for $G_2$ and $F_4$. The final theorems use \lean{G2Lie.g2W}, which is
\lean{simpleVOA g2CertData (-5/3) g2_level_ne} (\leanpath{VOA/DeligneAudit/G2LieData.lean}), and \lean{F4Lie.f4W}, which is
\lean{simpleVOA f4CertData (-5/2) f4_level_ne} (\leanpath{VOA/DeligneAudit/F4LieData.lean}), Lie data on the
structure-constant algebras \lean{g2Lie}, \lean{f4Lie} in the table basis, with the forms Killing$/8$ and Killing$/18$. The
library also defines \lean{g2SimpleVOA} and \lean{f4SimpleVOA} (\leanpath{VOA/AffineGeneral/G2Lie.lean},
\leanpath{VOA/AffineGeneral/F4Lie.lean}) from the \lean{standardData} of the same algebras with the same forms; no isomorphism
\lean{g2W} $\cong$ \lean{g2SimpleVOA} or \lean{f4W} $\cong$ \lean{f4SimpleVOA} is proved. Likewise three models of
$\mathfrak e_6$ coexist (a structure-constant table, the lattice $V_1$, and the \lean{E6Lat} certificate basis) and are not
identified; the $E_6$ final theorem uses the structure-constant model \lean{e6LieData}.

Each theorem is nevertheless about a genuine $L_k(\ltwentyseveng)$ for an explicit model of $\ltwentyseveng$: the Lie data carry
their own proofs of the Lie algebra axioms and of the invariant form (Chapter~\ref{ch:L23}). This is a remark on scope, not an
error. It does not concern $A_2$, whose VOA is the bespoke $\mathfrak{sl}_3$ construction.
\end{warning}

%======================================================================
\section{The named inputs}\label{l27:sec:inputs}
%======================================================================

This section defines Inputs T, R and S and the relations in which they are stated. For each input we give the mathematical
statement, the Lean definition and a fidelity check; the literature is the paragraph ``The literature behind the inputs'' at the
start of the chapter.

\subsection{Relation specifications and Ravinder's family}\label{l27:sec:relspec}

\begin{definition}[Relation specifications; Ravinder's relations]\label{l27:def:relspec}
A \emph{relation specification} for a Lie algebra with basis indexed by $\{0,\dots,\ltwentysevendg-1\}$ is a set of
\emph{tails} (words $x_1(r_1)\cdots x_s(r_s)$ in the current generators) and a set of \emph{eigenvalue conditions} $(h,\mu)$.
A vector $w$ of a current module \emph{satisfies} it if every tail kills $w$ and $h(0)w=\mu w$ for every eigenvalue
condition. The relations of $M_{a,b}$ of Section~\ref{l27:sec:fusionprimer} form the specification
$\mathrm{rav}(a+b,\,2a+b)$, where $\mathrm{rav}(k,d)$ consists of: $e(r)w=0$ for raising $e$ and all $r$; $h(r)w=0$ for Cartan
$h$ and $r\ge1$; $f_i(0)^{k\theta(h_i)+1}w=0$ for simple lowering $f_i$; $x(k)w=0$ for $x$ in the fusion set other than
$f_\theta$, and $f_\theta(d)w=0$; and the eigenvalues $h(0)w=k\theta(h)w$. The basis, the structure constants, the fusion set
and the numbers $\theta(h)$ are packaged as \emph{sector data} $S$ (Definition~\ref{l24:def:sectordata}). For $\mathfrak{sl}_3$ the
library also has a hand-written version, \lean{MabRel ρ w a b}.
\end{definition}
\noindent File \leanpath{VOA/\allowbreak DeligneAudit/\allowbreak DeligneFusionUnivG.lean}.
\begin{leancode}
structure RelSpec (n : ℕ) where
  tails : Set (List (Ltr n))
  eig : Set (Fin n × ℂ)

/-- The generator `w` of `ρ` satisfies `R`. -/
structure Sat {V : Type*} [AddCommGroup V] [Module ℂ V] (ρ : CurrentModule f V) (w : V)
    (R : RelSpec n) : Prop where
  tail : ∀ t ∈ R.tails, wordOp ρ t w = 0
  eig : ∀ p ∈ R.eig, ρ.op p.1 0 w = p.2 • w
\end{leancode}
\noindent File \leanpath{VOA/\allowbreak DeligneAudit/\allowbreak D4FusionKernelGeneric.lean}.
\begin{leancode}
def ravSpec (k d : ℕ) : RelSpec S.n where
  tails := {t | (∃ e r, S.raising e ∧ t = [(e, r)]) ∨ (∃ h r, S.cartan h ∧ 1 ≤ r ∧ t = [(h, r)]) ∨
    (∃ i, S.lowering i ∧ t = List.replicate (k * S.lowExp i + 1) (i, 0)) ∨
    (∃ x, S.fusion x ∧ t = [(x, if x = S.c₀ then d else k)]) ∨
    (∃ h r, S.cartanN h ∧ 1 ≤ r ∧ t = [(h, r)])}
  eig := {q | (∃ h, S.cartan h ∧ q = (h, ((k * S.thetaH h : ℕ) : ℂ))) ∨
    (∃ h, S.cartanN h ∧ q = (h, -((k * S.thetaN h : ℕ) : ℂ)))}
\end{leancode}
\begin{leancode}
def ravFam (a b : ℕ) : RelSpec S.n := ravSpec S (a + b) (2 * a + b)
\end{leancode}
\noindent File \leanpath{VOA/\allowbreak DeligneAudit/\allowbreak DeligneFusionRel.lean}.
\begin{leancode}
structure MabRel {V : Type} [AddCommGroup V] [Module ℂ V]
    (ρ : CurrentModule (bracketTable) V) (w : V) (a b : ℕ) : Prop where
  raising : ∀ (e : Fin 8), e.val < 3 → ∀ r, ρ.op e r w = 0
  cartan_pos : ∀ (h : Fin 8), 6 ≤ h.val → ∀ r, 1 ≤ r → ρ.op h r w = 0
  cartan_zero : ∀ (h : Fin 8), 6 ≤ h.val → ρ.op h 0 w = ((a + b : ℕ) : ℂ) • w
  lowering_f1 : (ρ.op 3 0)^[a + b + 1] w = 0
  lowering_f2 : (ρ.op 4 0)^[a + b + 1] w = 0
  trunc_f1 : ρ.op 3 (a + b) w = 0
  trunc_f2 : ρ.op 4 (a + b) w = 0
  trunc_theta : ρ.op 5 (2 * a + b) w = 0
\end{leancode}

\begin{dictionary}
\textbf{Words and letters.} \lean{Ltr n} is a pair (basis index, mode) $=(x,r)$, standing for $x(r)$; a tail is a list of
letters, and \lean{wordOp ρ t w} applies the word to $w$ (the rightmost letter acts first; for the one-letter and
repeated-letter tails used here the order does not matter). \lean{List.replicate k (i, 0)} is the word $f_i(0)^k$, and
\lean{(ρ.op 3 0)^[a + b + 1]} is the $(a+b+1)$-fold iterate of $f_1(0)$.

\textbf{The $f_\theta$ clause.} In \lean{ravSpec} the cut $f_\theta(d)w=0$ is the case \lean{x = S.c₀} of the fusion clause
\lean{[(x, if x = S.c₀ then d else k)]}, imposed only for \lean{x} with \lean{S.fusion x}. So it belongs to the specification
exactly when $c_0=f_\theta$ is in the fusion set. That membership is the field \lean{c0_fusion : S.fusion S.c₀} of the kernel
tables \lean{KernelTable S} and \lean{KernelTableSL S} (\leanpath{VOA/DeligneAudit/D4FusionKernelGeneric.lean}), and every sector
used by a final theorem comes with such a table.

\textbf{Section variables.} \lean{Sat} has an implicit \lean{\{n : ℕ\}} and an explicit structure-constant table
\lean{(f : Fin n → Fin n → Fin n → ℂ)}, so \lean{Sat S.f ρ w R} reads ``$w$ satisfies $R$ in the current module $\rho$ over the
table \lean{S.f}''. \lean{ravSpec} and \lean{ravFam} take \lean{S} as an explicit first argument, which is why the final
theorems write \lean{ravFam G2.G2sector}. The named inputs quantify over \lean{V : Type}, the lowest universe; every module
that occurs (universal modules, evaluation and lattice models) lives there.

\textbf{The two Cartan clauses.} Sector data split the Cartan basis vectors into those with $\theta(h)\ge0$
(\lean{S.cartan}) and those with $\theta(h)<0$ (\lean{S.cartanN}, stored with the nonnegative number \lean{S.thetaN h}); the
negative eigenvalues are written as $-(k\,\theta_N(h))$ because a natural number cannot be negative.

\textbf{$\mathfrak{sl}_3$ indices.} In \lean{bracketTable} the basis is $e_1,e_2,e_\theta$ (indices $0,1,2$), $f_1,f_2,f_\theta$
(indices $3,4,5$) and $h_1,h_2$ (indices $6,7$), with $\theta(h_1)=\theta(h_2)=1$. So \lean{MabRel ρ w a b} is
$\mathrm{rav}(a+b,2a+b)$ with fusion set $\{f_1,f_2,f_\theta\}$. The relation $f_\theta(0)^{2(a+b)+1}w=0$ is not listed; it follows
from the others ($\mathfrak{sl}_2$-integrability).

\textbf{Same as the literature?} For $(a,b)=(0,m)$ the library proves \lean{ravFam S 0 m = lwSpec S m} (\lean{ravFam_zero}), the
local Weyl and fusion relations. For $(a,b)=(j,0)$ these are the relations of Ravinder's Proposition~1 for $D(1,j\theta)$.
The sector data of $E_6$, $E_7$, $E_8$ use exactly Ravinder's fusion set; those of $A_2$, $D_4$, $G_2$, $F_4$ use all negative
root vectors. The extra vectors are the $f_\alpha$ with $\langle\alpha,\theta^\vee\rangle=0$ (every other positive root
$\alpha\neq\theta$ has $\langle\alpha,\theta^\vee\rangle=1$ and is already in Ravinder's set, short roots of $G_2$ and $F_4$
included). For these, $f_\alpha(r)w=0$ holds for \emph{every} $r$ in any module with the local Weyl relations, by the
$h(r)$-bracket argument of Lemma~\ref{l27:lem:truncweyl}. So the two readings impose the same relations, for every type. For
Input T this is in any case harmless, since extra relations make the hypothesis weaker; for Input R it means that the
Demazure module $D(1,j\theta)$ satisfies the Lean relations, short roots included.
\end{dictionary}

\begin{physics}
\emph{Heuristic.} A relation specification is a list of selection rules for a composite state: which current modes
annihilate it. $M_{a,b}$ describes $a+b$ fused adjoint ``particles'', of which $a$ carry the extra one-dimensional excitation
of $D(1,\theta)=\ltwentyseveng\oplus\C$ and $b$ do not; the mode cutoffs encode how much energy the fused system can absorb.
\end{physics}

\begin{status}
Definitions. \lean{ravFam_zero} (Chapter~\ref{ch:L26}) identifies the $a=0$ case with the local Weyl and fusion relations;
axioms standard (campaign record).
\end{status}

\subsection{Input T: the total-dimension bound}\label{l27:sec:totaldim}

\begin{definition}[Input T]\label{l27:def:totaldim}
For sector data $S$ of $\ltwentyseveng$, \lean{FusionTotalDimG S} is the statement: for every $m\ge0$, every current module $V$
and every $w\in V$ satisfying the local Weyl and fusion relations for $m\theta$ (that is, $\mathrm{rav}(m,m)$), and every $B$,
\begin{equation*}
\sum_{b<B}\dim\bigl(\text{degree-}b\text{ part of the cyclic span of }w\bigr)\ \le\ \ltwentysevendg^{\,m} .
\end{equation*}
Equivalently, $\dim\ltwentysevenF_m\le\ltwentysevendg^{\,m}$. For $A_2$ the library has the same statement over the
hand-written $\mathfrak{sl}_3$ relations, \lean{SectorWeyl.FusionTotalDim}, with bound $8^m$.
\end{definition}
\noindent File \leanpath{VOA/\allowbreak DeligneAudit/\allowbreak DeligneFusionFiniteG.lean}.
\begin{leancode}
def FusionTotalDimG : Prop :=
  ∀ (m : ℕ) (V : Type) [AddCommGroup V] [Module ℂ V] (ρ : CurrentModule S.f V) (w : V),
    LocalWeylRelG S ρ w m → ∀ B : ℕ, ∑ b ∈ Finset.range B, Module.finrank ℂ (cycDeg ρ w b) ≤ S.n ^ m
\end{leancode}
\noindent File \leanpath{VOA/\allowbreak DeligneAudit/\allowbreak DeligneFusionFinite.lean}.
\begin{leancode}
def FusionTotalDim : Prop :=
  ∀ (m : ℕ) (V : Type) [AddCommGroup V] [Module ℂ V] (ρ : CurrentModule (bracketTable) V) (w : V),
    LocalWeylRel ρ w m → ∀ B : ℕ, ∑ b ∈ Finset.range B, Module.finrank ℂ (cycDeg ρ w b) ≤ 8 ^ m
\end{leancode}

\begin{dictionary}
\textbf{Decoding.} \lean{LocalWeylRelG S ρ w m} is the structure of Definition~\ref{l24:def:U5}: $w$ satisfies the local Weyl
relations of $W_{\mathrm{loc}}(m\theta)$ and the fusion relations $x(m)w=0$ for $x$ in the fusion set. The $A_2$ predicate
\lean{LocalWeylRel ρ w m} is the same for the $\mathfrak{sl}_3$ table (the library proves
\lean{MabRel ρ w 0 m ↔ LocalWeylRel ρ w m}, \lean{mabRel_zero_iff}). The identification of \lean{LocalWeylRelG S ρ w m} with $\mathrm{rav}(m,m)$ is the Lean theorem \lean{UnivG.localWeylRelG_iff} (equivalence with \texttt{Sat S.f ρ w (lwSpec S m)}), together with \lean{ravFam_zero} (\texttt{ravFam S 0 m = lwSpec S m}); it is not a rewrite, since \lean{ravSpec} writes the powers as \texttt{List.replicate} while \lean{LocalWeylRelG} uses iterates. \lean{S.n} is $\ltwentysevendg$. The quantifier ``$\forall B$''
bounds every finite partial sum, which is the same as bounding the total dimension.

\textbf{An upper bound over all modules.} The statement quantifies over \emph{all} modules and generators with the
relations, so it is a statement about the universal module $M_{0,m}$, of which all the others are quotients.

\textbf{The finiteness clause.} \lean{Module.finrank} is $0$ on an infinite-dimensional space, so read alone the bound would
say nothing about a module with an infinite-dimensional degree piece. The library proves separately, for every type, that
every degree piece is finite-dimensional (\lean{fusionTotalBoundG_of_dim}; Chapter~\ref{ch:L26},
Section~\ref{l26:sec:finiteG}), so the input that remains is only the bare bound displayed above.

\textbf{Same as the literature?} With the finiteness clause it is exactly ``$\dim\ltwentysevenF_m\le\ltwentysevendg^{\,m}$'',
and by Lemma~\ref{l27:lem:truncweyl} below $\ltwentysevenF_m$ is the truncated Weyl module $W(m\theta,m)$ of Kus--Littelmann,
whichever fusion set the sector uses. The published statement is the equality $W(m\theta,m)\cong\ltwentyseveng^{*m}$:
first in Feigin, \emph{SIGMA} 4 (2008) 070, Corollary~2, whose proof has a gap; then in Kus--Littelmann, Theorem~4.2 at
$\lambda=\theta$, $\lambda_0=0$, $N=m$, a repair valid for every simple $\ltwentyseveng$, with one step (p.~437) that the
project flags as unjustified as written, and that for $D_4$ is shown by computation to be insufficient as written at $m=2$
(item~(i) of the literature paragraph; the theorem is not known to fail). Only the upper half of the equality is used;
the lower half is free.

\textbf{Could it be vacuous or false?} Modules with the relations exist (the universal one, and the evaluation product
$\ltwentyseveng(c_1)\otimes\dots\otimes\ltwentyseveng(c_m)$ of dimension exactly $\ltwentysevendg^{\,m}$), so the bound, where true, is sharp.
If it were false for a type, the conditional final theorem of that type would be vacuously true.
\end{dictionary}

\begin{physics}
\emph{Heuristic.} $\ltwentysevenF_m$ is the space of states of $m$ adjoint particles brought together, keeping only the
excitations allowed by the relations. The bound says that bringing the particles together creates no new states.
\end{physics}

\begin{status}
Definitions. For $A_2$, \lean{SectorWeyl.FusionTotalDim} is \textbf{proved}: \lean{CVUniv.fusionTotalDim} (Chapter~\ref{ch:L26},
Theorem~\ref{l26:thm:cvtotal}; \texttt{\#print axioms}-clean per the targeted build record). For $G_2$, $D_4$, $F_4$, $E_6$, $E_7$,
$E_8$, \lean{FusionTotalDimG S} is a \textbf{named hypothesis} of the final theorem, not proved in Lean. In the literature it is
\textbf{published} for all six types (Kus--Littelmann, Theorem~4.2), with one proof step flagged by the project; for $D_4$
that step is insufficient as written at $m=2$ (item~(i) of the literature paragraph). Library progress:
for $D_4$ the cases $m=1,2$ are proved (Section~\ref{l27:sec:partial}); toward all $m$, Lean reduces the bound to
$\mathfrak{sl}_4$ core bounds with at least two adjoint sites, and on paper the levels with two adjoint sites are reduced further
to two one-node bounds, both open in general (end of Section~\ref{l27:sec:feigin}); for the other types only the finiteness clause.
\end{status}

\begin{lemma}[The fusion module is a truncated Weyl module]\label{l27:lem:truncweyl}
Let $m\ge0$ and let $w$ be a vector of a $\ltwentyseveng[t]$-module satisfying the local Weyl relations of
$W_{\mathrm{loc}}(m\theta)$, together with $x(m)w=0$ for every $x$ in a set $\mathcal S$ of negative root vectors with
$\{f_\alpha:\langle\alpha,\theta^\vee\rangle=1\}\cup\{f_\theta\}\subseteq\mathcal S$. Then $(\ltwentyseveng\otimes t^m\C[t])$
annihilates the cyclic module $U(\ltwentyseveng[t])w$. Consequently the universal module with these relations,
$\ltwentysevenF_m=M_{0,m}$, is the \emph{truncated Weyl module} $W(m\theta,m)$ of Kus--Littelmann (Definition~4.1): the
quotient of $W_{\mathrm{loc}}(m\theta)$ on which $\ltwentyseveng\otimes t^m\C[t]$ acts by zero. In particular $M_{0,m}$ does not
depend on the choice of $\mathcal S$ within these bounds.
\end{lemma}

\textbf{Proof on paper.} This is not a Lean statement; it is the bridge between the Lean hypothesis and the published theorem.
Write $\mathcal I_m=\ltwentyseveng\otimes t^m\C[t]$. It suffices to show $\mathcal I_mw=0$: since $\mathcal I_m$ is an ideal of
$\ltwentyseveng[t]$, for $y\in\ltwentyseveng[t]$ and $i\in\mathcal I_m$ we have $i\,y=y\,i-[y,i]$ with $[y,i]\in\mathcal I_m$, so by
induction on the length of words $\mathcal I_m\,U(\ltwentyseveng[t])\subseteq U(\ltwentyseveng[t])\,\mathcal I_m$, and therefore
$\mathcal I_m\,U(\ltwentyseveng[t])w\subseteq U(\ltwentyseveng[t])\,\mathcal I_mw=0$. The space $\mathcal I_m$ is spanned by
$e_\alpha(s)$, $h(s)$ and $f_\alpha(s)$ with $s\ge m$.
\begin{itemize}
\item $e_\alpha(s)w=0$ for all $s$, and $h(s)w=0$ for $s\ge1$, are relations of $W_{\mathrm{loc}}(m\theta)$. (For $m=0$ the
  remaining claim $h(0)w=0$ is the eigenvalue relation $h(0)w=0\cdot\theta(h)w$.)
\item \emph{The bracket step.} For a positive root $\alpha$ choose $h$ in the Cartan subalgebra with $\alpha(h)\neq0$. In
  $\ltwentyseveng[t]$, $[h(1),f_\alpha(s)]=-\alpha(h)\,f_\alpha(s+1)$. Applying both sides to $w$ and using $h(1)w=0$ gives
  \begin{equation*}
  f_\alpha(s+1)\,w=-\alpha(h)^{-1}\,h(1)\,f_\alpha(s)\,w .
  \end{equation*}
  So $f_\alpha(s)w=0$ implies $f_\alpha(s+1)w=0$, and by induction $f_\alpha(s')w=0$ for all $s'\ge s$.
\item If $\langle\alpha,\theta^\vee\rangle=1$ or $\alpha=\theta$, then $f_\alpha\in\mathcal S$, so $f_\alpha(m)w=0$, and the
  bracket step gives $f_\alpha(s)w=0$ for all $s\ge m$.
\item If $\langle\alpha,\theta^\vee\rangle=0$, then also $\langle\theta,\alpha^\vee\rangle=0$. The degree-zero relations make
  $U(\ltwentyseveng)w$ a finite-dimensional highest-weight module of highest weight $m\theta$, and in such a module
  $f_\alpha(0)w=0$ exactly when $m\theta(h_\alpha)=m\langle\theta,\alpha^\vee\rangle=0$. So $f_\alpha(0)w=0$, and the bracket
  step gives $f_\alpha(s)w=0$ for all $s\ge0$.
\end{itemize}
Since $\theta$ is long, every positive root $\alpha\neq\theta$ has $\langle\alpha,\theta^\vee\rangle\in\{0,1\}$, so the cases are
exhaustive and $\mathcal I_mw=0$. Hence $M_{0,m}$ is a quotient of $W(m\theta,m)$. Conversely, the generator of $W(m\theta,m)$
satisfies all the relations (the extra ones $x(m)w=0$ hold because $x(m)\in\mathcal I_m$), so $W(m\theta,m)$ is a quotient of
$M_{0,m}$, and the two cyclic modules coincide. The same argument shows that the relations of the Lean family
$\mathrm{rav}(j,2j)$ (Definition~\ref{l27:def:relspec}) hold on the extremal vector of $D(1,j\theta)$ for the sectors whose
fusion set is larger than Ravinder's: the extra $f_\alpha$ have $\langle\alpha,\theta^\vee\rangle=0$ and are covered by the last
case. \qed

\begin{physics}
\emph{Heuristic.} Once the modes of degree $m$ of all currents annihilate the composite state of $m$ adjoint particles, so do
all higher modes: an energy cutoff on one mode propagates to all higher ones, because the higher modes are generated from the
lower ones by the Cartan current $h(1)$, which acts trivially on the highest-weight state.
\end{physics}

\begin{status}
\textbf{On paper only}; no Lean statement. The argument is the one recorded in the project's audit of the literature
inputs (\texttt{DELIGNE\_INPUTS\_PRECISION.md}, \S1).
\end{status}

\subsection{Input R: Demazure realizations}\label{l27:sec:realization}

\begin{definition}[Input R]\label{l27:def:realization}
Let $S$ be sector data, $\mathrm{spec}$ a family of relation specifications indexed by $\N\times\N$, and
$\mathrm{Dem}:\N\to\Q[[q]]$ a family of power series. For $j\in\N$, a \emph{Demazure realization at $j$} is a current module
$V$ and a vector $w\in V$ satisfying $\mathrm{spec}(j,0)$ such that, for every degree $N$,
\begin{equation*}
[q^N]\,\mathrm{Dem}_j\ \le\ \dim\bigl(\text{degree-}N\text{ part of the cyclic span of }w\bigr).
\end{equation*}
In the final theorems $\mathrm{spec}=$ \lean{ravFam S}, so $\mathrm{spec}(j,0)$ are the relations of $M_{j,0}$, and $\mathrm{Dem}_j$
is the type's Demazure table from Chapter~\ref{ch:L25}. Input R is ``a Demazure realization exists at every $j$''.
\end{definition}
\noindent File \leanpath{VOA/\allowbreak DeligneAudit/\allowbreak DeligneFusionUnivGSector.lean}.
\begin{leancode}
def DemazureRealizationG (spec : ℕ → ℕ → RelSpec S.n) (Dem : ℕ → PowerSeries ℚ) (n : ℕ) : Prop :=
  ∃ (V : Type) (_ : AddCommGroup V) (_ : Module ℂ V) (ρ : CurrentModule S.f V) (w : V),
    Sat S.f ρ w (spec n 0) ∧ ∀ N, coeff N (Dem n) ≤ (Module.finrank ℂ (cycDeg ρ w N) : ℚ)
\end{leancode}
\noindent File \leanpath{VOA/\allowbreak DeligneAudit/\allowbreak DeligneCarrierDemazure.lean}.
\begin{leancode}
def hilbSeries (hilb : ℕ → ℕ → ℕ) (n : ℕ) : PowerSeries ℚ :=
  PowerSeries.mk fun N => (hilb n N : ℚ)
\end{leancode}
\noindent File \leanpath{VOA/\allowbreak DeligneAudit/\allowbreak DeligneCharacterG2FusionZ.lean}.
\begin{leancode}
def g2HilbInt (m : ℕ) (N : ℤ) : ℤ := hilbIntG cartanG2 thetaG2 wordG2 m N
\end{leancode}
\begin{leancode}
def g2HilbSeriesZ (n : ℕ) : PowerSeries ℚ := PowerSeries.mk fun N => (g2HilbInt n N : ℚ)
\end{leancode}

\begin{dictionary}
\textbf{Decoding.} The existential packages a vector space \lean{V}, a current-module structure \lean{ρ} over the structure
constants \lean{S.f}, and a vector \lean{w}. \lean{Sat S.f ρ w (spec n 0)} says $w$ satisfies the relations of $M_{n,0}$
(Lean's \lean{n} is our $j$). \lean{coeff N (Dem n)} is $[q^N]\mathrm{Dem}_j$. The inequality is in $\Q$.

\textbf{It is a lower bound.} The statement asks for \emph{one} module. Since every module with these relations is a quotient
of the universal $M_{j,0}$, it says that $M_{j,0}$ is at least as large as $\mathrm{Dem}_j$ in every degree. It is the lower
bound (R) of the squeeze of Example~\ref{l27:ex:sl3}.

\textbf{Clamped and unclamped tables.} The Demazure algorithm of Chapter~\ref{ch:L25} produces integers. A true character has
nonnegative coefficients. The \emph{clamped} table replaces any negative output by $0$ (the last step of
\lean{levelOneDemazureHilb} is \lean{Int.toNat}); \lean{hilbSeries hilb n} turns such a table into the power series
$\sum_N\mathrm{hilb}(n,N)q^N$, and the tables \lean{f4DemazureHilb}, \lean{e6DemazureHilb}, \lean{e7DemazureHilb},
\lean{e8DemazureHilb} are clamped. The \emph{unclamped} tables keep the integers: \lean{g2HilbInt} (through \lean{hilbIntG}) and
\lean{D4Wp.hilbSeriesZ}. For Input R the two forms are \emph{equivalent}, not merely ordered: a requirement
$c\le\dim(\cdot)$ is vacuous both when $c<0$ and when $c=0$, because a dimension is never negative, and clamping changes only
negative entries into $0$. So Input R against the clamped table and Input R against the unclamped table are the same
condition on a module. The final theorems use the unclamped
tables for $G_2$ and $D_4$ and the clamped ones for $F_4$, $E_6$, $E_7$, $E_8$. For $G_2$ and $D_4$ the unclamped route has an
advantage: the library derives the nonnegativity of the algorithm's output from U5 (\lean{g2HilbInt_nonneg_of_U5Z} for $G_2$),
so no separate positivity hypothesis is needed. Elsewhere in this chapter we refer to this paragraph.

\textbf{Same as the literature?} With the true characters of $D(1,j\theta)$, the statement is ``$D(1,j\theta)$ satisfies Ravinder's
relations, and its graded dimension is the Demazure character''. Both halves are \emph{published theorems for every simple
$\ltwentyseveng$}. The relation half is Ravinder's Proposition~1 (after Chari--Venkatesh, Corollary~3.5): $D(1,j\theta)$ is
presented by $\mathfrak n^+[t]w=0$, $h(s)w=\delta_{s,0}\,j\theta(h)w$, $f_\alpha(0)w=0$ for $\langle\alpha,\theta^\vee\rangle=0$,
$f_\alpha(0)^{jd_\alpha+1}w=0$ (with $d_\alpha=2/(\alpha,\alpha)=\langle\theta,\alpha^\vee\rangle$) and $f_\alpha(j)w=0$ for
$\langle\alpha,\theta^\vee\rangle=1$, and $f_\theta(0)^{2j+1}w=0$. The Lean relations
$\mathrm{rav}(j,2j)$ follow from these: $f_\theta(2j)w=0$ holds already in $W_{\mathrm{loc}}(j\theta)$ (the $\theta$-triple sees
weight $2j$, and the $\mathfrak{sl}_2$ box kills $f_\theta(s)w$ for $s\ge2j$), and the extra fusion-set roots of the $D_4$, $G_2$,
$F_4$ sectors are covered by Lemma~\ref{l27:lem:truncweyl}. Only this direction (``$D(1,j\theta)$ satisfies the Lean relations'')
is needed, and it holds for every type, short roots included. The graded half is the Demazure character formula (Kumar 1987,
Mathieu 1988), together with the standard identification of $D(1,j\theta)$ with the Demazure module
$V_{(s_\theta s_0)^j}(\Lambda_0)$, which is how the Lean algorithm computes $\mathrm{Dem}_j$ (Chapter~\ref{ch:L25},
Caveat~\ref{l25:cav:U7}). What is \emph{not} established is formal: that the Lean algorithm (\lean{levelOneDemazureHilb},
\lean{hilbIntG}) implements the Demazure character formula correctly (Cartan-matrix orientation, the reduced word for $s_\theta$,
the step $D_0$, which is valid because $\theta$ is long, and the regrading by $j^2$). The evidence for the encoding is that the proved
totals agree with $(\ltwentysevendg+1)^j$ (for example \lean{g2_demazure_total}), which rules out a transposed Cartan matrix for
$G_2$ and $F_4$, together with numerical checks; this is evidence, not a proof.
\end{dictionary}

\textbf{What is proved for $A_2$.} For $\mathfrak{sl}_3$ the library proves the realization for every $j$:
\lean{EvalModel.demazureRealization_all} constructs an explicit evaluation model of $D(1,j\theta)$ on a tensor product of copies
of $D(1,\theta)=\mathfrak{sl}_3\oplus\C$ at distinct points, checks the relations, and bounds its graded dimension from below
through a filtration and Lagrange interpolation, with the matching upper bound from the Chari--Loktev spanning theorem, which is
special to $\mathfrak{sl}_{r+1}$ (Chapter~\ref{ch:L26}, Section~\ref{l26:sec:u7}).


\begin{physics}
\emph{Heuristic.} $D(1,j\theta)$ is a truncation of the level-one vacuum module of the affine Lie algebra (the WZW vacuum
module): the states reachable from the extremal vector of weight $j\theta$ by positive modes. A realization is an explicit
system of $j$ excitable adjoint particles whose energy spectrum is at least as rich as the Demazure character predicts.
\end{physics}

\begin{status}
Definitions. \lean{∀ n, DemazureRealizationG S (ravFam S) Dem n} is a \textbf{named hypothesis} for $G_2$, $F_4$, $E_6$,
$E_7$, $E_8$. For $A_2$ the corresponding statement \lean{∀ n, A2.DemazureRealization n} is a theorem
(\lean{EvalModel.demazureRealization_all}; axioms standard, re-checked at an earlier commit). For $D_4$ it is a theorem,
\lean{D4Char.d4_demazureRealization} (Theorem~\ref{l27:thm:d4R}; aggregate build green on 3 November 2026, standard axioms), and
the $D_4$ theorem \lean{D4_classicallyFree_of_inputT} no longer assumes it. In the literature Input R is a
\textbf{published theorem for all six types} (Ravinder, Proposition~1, with the Demazure character formula); for $G_2$, $F_4$,
$E_6$, $E_7$, $E_8$ only its formalization is missing. For $E_6$, $E_7$, $E_8$ the relation half holds in the lattice model for
every $n$.
\end{status}

\textbf{What is proved for $D_4$.} For $\mathfrak{so}_8$ the library also proves the realization for every $j$, by a different
route: a graded upper bound on the Levi side and the squeeze, rather than the Demazure character formula.

\begin{theorem}[Input R for $D_4$]\label{l27:thm:d4R}
For every $n$ there is a current module for the $D_4$ structure constants with a vector satisfying the relations
$\mathrm{rav}(n,2n)$ whose degree-$N$ piece has dimension at least $[q^N]$ of the unclamped $D_4$ Demazure series, for every $N$.
\end{theorem}
\noindent File \leanpath{VOA/\allowbreak DeligneAudit/\allowbreak D4CharDemazureRealization.lean}.
\begin{leancode}
theorem d4_demazureRealization :
    ∀ n, DemazureRealizationG D4sector (ravFam D4sector) D4Wp.hilbSeriesZ n :=
\end{leancode}

\begin{dictionary}
The full name is \lean{VOA.DeligneAudit.D4Char.d4_demazureRealization}; \lean{D4sector} is \lean{DeligneSector.D4.D4sector}. The type
is exactly the hypothesis \lean{hD} of \lean{D4_classicallyFree_final} (Theorem~\ref{l27:thm:D4}): same sector, same family
\lean{ravFam}, same unclamped series \lean{D4Wp.hilbSeriesZ}. It takes no arguments. The proof bounds every module with the
relations, degree by degree, by a Levi character (\lean{D4Char.d4_gradedLeviBound}), and closes with the squeeze.
\end{dictionary}

\begin{status}
Proved, \textbf{unconditional}. Machine-confirmed: aggregate build green on 3 November 2026, and \texttt{\#print axioms} reports
\lean{[propext, Classical.choice, Quot.sound]}. It discharges Input R for $D_4$ (Theorem~\ref{l27:thm:d4inputT}).
\end{status}

\begin{physics}
For $L_{-2}(\mathfrak{so}_8)$, the chiral algebra of $SU(2)$ gauge theory with four flavours, this supplies the lower
bound on the number of Schur operators in each fusion sector; with Input~T it pins the Schur index to the arc-space
count of the Higgs-branch chiral ring.
\end{physics}

\subsection{Input S: the star bounds}\label{l27:sec:starinput}

\begin{definition}[Input S]\label{l27:def:star}
For a Demazure table $\mathrm{hilb}$ and a target sequence $c:\N\to\Q$, the \emph{star bound} is $\mathcal D_N\le c_N$ for all
$N$, where $\mathcal D_N=\sum_{j\le N}[q^{N-j}]\bigl(\mathrm{Dem}_j/(q)_j\bigr)$ is the Demazure sum. For $X=G_2,F_4,E_6,E_7,E_8$,
\lean{XStarBound} is the star bound with the type's clamped Demazure table and $c_N=\mathtt{mldeCoeff}\ h^\vee\ N$, where
$h^\vee=4,9,12,18,30$.
\end{definition}
\noindent File \leanpath{VOA/\allowbreak DeligneAudit/\allowbreak DeligneCharacterStarBound.lean}.
\begin{leancode}
def StarBoundQ (hilb : ℕ → ℕ → ℕ) (c : ℕ → ℚ) : Prop :=
  ∀ N, (demazureSumCoeff hilb N : ℚ) ≤ c N
\end{leancode}
\begin{leancode}
abbrev G2StarBound : Prop := StarBoundQ g2DemazureHilb (mldeCoeff 4)
abbrev F4StarBound : Prop := StarBoundQ f4DemazureHilb (mldeCoeff 9)
abbrev E6StarBound : Prop := StarBoundQ e6DemazureHilb (mldeCoeff 12)
abbrev E7StarBound : Prop := StarBoundQ e7DemazureHilb (mldeCoeff 18)
abbrev E8StarBound : Prop := StarBoundQ e8DemazureHilb (mldeCoeff 30)
\end{leancode}

\begin{dictionary}
\lean{demazureSumCoeff hilb N} is $\mathcal D_N=\sum_{j\le N}\sum_i\mathrm{hilb}(j,i)\,p_{\le j}(N-j-i)$, where $p_{\le j}(r)$ is the
number of partitions of $r$ into parts $\le j$ (the coefficients of $1/(q)_j$); Definition~\ref{l25:def:dsum}. \lean{mldeCoeff h N}
is the rational number $c_N$ defined by the Arakawa--Kawasetsu recurrence with parameter $h$ (Definition~\ref{l25:def:mlde}); it is
\lean{VOA.DeligneAudit.DeligneCharacter.mldeCoeff}, not the homonym in \leanpath{VOA/Modularity/AKMLDE.lean}. The second arguments
$4,9,12,18,30$ are $h^\vee$ of $G_2,F_4,E_6,E_7,E_8$, as they must be. All five use the clamped tables; for $G_2$ the final theorem's
Input R uses the unclamped table, and the two agree under the other inputs (dictionary of Definition~\ref{l27:def:realization}).

\textbf{One-sided form.} The paper states the star \emph{identity} $\mathcal D_N=c_N$. Only $\le$ is used, because it is chained with
the proved $c_N\le\dim L_N$. The library's \lean{XStarIdentity} (equality) implies \lean{XStarBound} (\lean{g2StarBound_of_identity}
and its four siblings). Under the other inputs the chain closes and equality follows (Remark~\ref{l27:rem:arc}).
\end{dictionary}

\textbf{Literature status: a conjecture.} For $X\in\{G_2,F_4,E_6,E_7,E_8\}$ we know of \emph{no published proof} of the star
identity $\mathcal D_N=c_N$, nor of the bound $\mathcal D_N\le c_N$ alone (Section~\ref{l27:sec:starcore} explains why the bound
is no cheaper). We therefore state it as a conjecture, and the five final theorems that use it are conditional on it.

\begin{conjecture}[The exceptional star identities]\label{l27:conj:star}
For $X\in\{G_2,F_4,E_6,E_7,E_8\}$, with $h^\vee=4,9,12,18,30$ and $\mathrm{Dem}_j$ the graded character of the level-one Demazure
module $D(1,j\theta)$,
\begin{equation*}
\sum_{j\ge0}q^j\,\frac{\mathrm{Dem}_j(q)}{(q)_j}\ =\ \sum_{N\ge0}c_N\,q^N ,
\end{equation*}
where $c_N$ is the solution with $c_0=1$ of the Arakawa--Kawasetsu recurrence with parameter $h^\vee$. In particular
\lean{XStarBound} holds.
\end{conjecture}

\textbf{Evidence.} Exact-integer checks of the coefficientwise equality $\mathcal D_N=c_N$ hold with no discrepancy through $q^{60}$ for $G_2$, $q^{25}$ for $F_4$, $q^{20}$ for $E_6$, $q^{16}$ for $E_7$ and $q^{14}$ for $E_8$ (scoping pass of 28 October 2026, \texttt{Notes/deligne/STAR\_IDENTITY\_SCOPING.md}; cross-validated against the library's Demazure algorithm). No published proof of the identity or of the inequality is known to us. These are evidence, not proof.

\textbf{The $G_2$ reduction.} For $G_2$ the library has reduced Input S to one identity. The target half \lean{G2KWEqMLDE} (the
Kac--Wakimoto-type product $\theta_{A_2}(q)(q^3;q^3)_\infty^6/(q;q)_\infty^8$ solves the $h^\vee=4$ recurrence) is a theorem,
\lean{G2Star.g2KWEqMLDE}, and \lean{G2Star.g2StarIdentity_of_demazure} derives the star identity from the single named
statement \lean{G2DemazureEqKW} (Demazure sum $=$ that product); \lean{g2StarBound_of_identity} then gives \lean{G2StarBound}
(Chapter~\ref{ch:L25}, Definition~\ref{l25:def:starbounds}, where these declarations are quoted). \lean{G2DemazureEqKW} has no
published proof known to us either. The library splits it into two named halves (Theorem~\ref{l27:thm:G2F4ladder}).
\begin{itemize}
\item \lean{G2DemazureFermionic} (the $G_2$ Demazure table equals the HKOTY fermionic formula) \emph{follows from published
  results}, as recorded in the project note \texttt{Notes/deligne/G2\_STAR\_FERMIONIC.md} (28 October 2026). Fourier--Littelmann
  (Adv.\ Math.\ 211 (2007), Theorem~8 and Theorem~B) give $D(1,j\theta)\cong\mathrm{KR}(\omega_2)^{*j}$ as graded modules, and
  Chari--Moura identify $\mathrm{KR}(\omega_2)=W^{(2)}_1=V(\theta)\oplus\C$; Ardonne--Kedem (Theorem~5.2) bound the graded fusion
  multiplicities by the fermionic sums; Di Francesco--Kedem (Corollary~5.19, with the HKOTY and Hernandez results as they state
  them) give equality at $q=1$, hence coefficientwise equality by positivity. The conventions are those of Lin (IMRN 2021), and the
  note checks that they match the Lean definition with no shift. The note read the theorem statements in the arXiv versions of
  these sources; the primary texts of HKOTY and Hernandez were not opened. This is a reading of the literature, not a Lean proof:
  \lean{G2DemazureFermionic} remains a hypothesis, and the identification of \lean{g2DemazureHilb} with
  $\mathrm{Hilb}_q D(1,j\theta)$ is the same unformalized step (Demazure character formula) as for every Demazure table.
\item \lean{G2FermionicKW} (the Demazure sum of the fermionic table equals the Kac--Wakimoto-type product) is \emph{open}: a
  rank-two $q$-series identity with no known proof. It is checked numerically through $q^{60}$, and in a weight-refined form
  (every branching function) through $q^{50}$ (same note).
\end{itemize}
So the $G_2$ star identity remains a conjecture, now reduced to the single open identity \lean{G2FermionicKW}.

\begin{physics}
\emph{Heuristic.} For the types with a four-dimensional theory ($E_6$, $E_7$, $E_8$, and $D_4$, $A_2$ where the identity is proved),
the star identity would equate two counts of the Schur operators: the Demazure sum counts ``particle'' states ($j$ adjoint quanta with
their Demazure excitations, dressed by the factor $1/(q)_j$ of their relative motion), and the MLDE solution is the Schur index as
Beem and Rastelli computed it from a modular differential equation. For $G_2$ and $F_4$ no four-dimensional theory is known, and the
identity is a statement about $q$-series only.
\end{physics}

\begin{status}
Definitions. \lean{G2StarBound}, \lean{F4StarBound}, \lean{E6StarBound}, \lean{E7StarBound}, \lean{E8StarBound} are \textbf{named
hypotheses}, and in the literature they are \textbf{conjectural} (Conjecture~\ref{l27:conj:star}; numerical evidence only). For
$G_2$ the hypothesis is implied by the single named identity \lean{G2DemazureEqKW}. The $A_2$ star identity is a theorem (\lean{A2.starIdentity}, Chapter~\ref{ch:L25}); for $D_4$ the corresponding statement
is proved (\lean{D4Wp.d4SchurEqE4}, Theorem~\ref{l25:thm:d4}), so the $D_4$ final theorem has no star hypothesis.
\end{status}

\subsection{An alternative form of the \texorpdfstring{$A_2$}{A2} Input T}\label{l27:sec:a2alt}

Before Input T was proved for $A_2$, the library reduced it to a lower bound in a single model. That reduction is still the
hypothesis of one of the two $A_2$ wrappers (Theorem~\ref{l27:thm:A2}), so we state it.

\begin{definition}[Kernel lower bound]\label{l27:def:a2alt}
\lean{RavinderKernelLower}: for all $a,b\ge0$ there is a current $\mathfrak{sl}_3$-module with a generator $w$ satisfying the
relations of $M_{a+1,b}$, in which the cyclic span of the kernel vector $f_\theta(2a+b+1)w$ has total dimension at least $9^a8^b$.
\end{definition}
\noindent File \leanpath{VOA/\allowbreak DeligneAudit/\allowbreak DeligneFusionKernelLower.lean}.
\begin{leancode}
def RavinderKernelLower : Prop :=
  ∀ a b : ℕ, ∃ (V : Type) (_ : AddCommGroup V) (_ : Module ℂ V) (ρ : CurrentModule (bracketTable) V)
    (w : V), MabRel ρ w (a + 1) b ∧
    ∃ B, 9 ^ a * 8 ^ b ≤ ∑ n ∈ Finset.range B,
      Module.finrank ℂ (cycDeg ρ (ρ.op 5 (2 * a + b + 1) w) n)
\end{leancode}

\begin{dictionary}
\lean{ρ.op 5 (2 * a + b + 1) w} is $f_\theta(2a+b+1)w$ (index $5$ is $f_\theta$). \lean{cycDeg ρ v n} is graded by the words applied to
$v$, so the shift by $2a+b+1$ is not included; this does not affect a total dimension. $9^a8^b$ is the expected dimension of
$M_{a,b}$. In Ravinder's sequence (\ref{l27:eq:ravinder}) with $(a+1,b)$ in place of $(a,b)$, the cyclic span of
$f_\theta(2a+b+1)w$ is the kernel, and the statement says it is as large as $\tau M_{a,b}$: the injectivity half of the sequence,
realized in a model.
\end{dictionary}

\begin{physics}
\emph{Heuristic.} When one excitable adjoint particle loses its extra excitation, the states that disappear are as many as a
whole smaller system of particles.
\end{physics}

\begin{status}
Definition; \textbf{not proved}, and no longer needed: Input T for $A_2$ is proved by a different route. It is the hypothesis of
\lean{A2_classicallyFree_final_S2} only.
\end{status}

\begin{theorem}[The kernel lower bound implies Input T for $A_2$]\label{l27:thm:a2alt}
\lean{RavinderKernelLower} implies \lean{FusionTotalDim}.
\end{theorem}
\noindent File \leanpath{VOA/\allowbreak DeligneAudit/\allowbreak DeligneFusionKernelLower.lean}.
\begin{leancode}
theorem fusionTotalDim_of_kernelLower (hK : RavinderKernelLower) : FusionTotalDim := by
\end{leancode}

\begin{dictionary}
An implication between named propositions; no other hypothesis. \lean{FusionTotalDim} here is \lean{SectorWeyl.FusionTotalDim}
(the file opens \lean{SectorWeyl}). A second reduction, from an ``ideal lower bound'' \lean{IdealLower}, is presented in
Chapter~\ref{ch:L26} (Theorem~\ref{l26:thm:a2red}); no final theorem uses it.
\end{dictionary}

\textbf{Proof on paper.} Let $U_{a,b}$ be the universal module for the relations of $M_{a,b}$ and $u_{a,b}$ its total
dimension. The relations of $M_{a,b+1}$ imply those of $M_{a+1,b}$, so there is a surjection $U_{a+1,b}\to U_{a,b+1}$ that kills
$f_\theta(2a+b+1)w$ and hence its cyclic span. Rank--nullity gives $u_{a,b+1}+\dim(\text{cyclic span of }f_\theta(2a+b+1)w)\le
u_{a+1,b}$, and the hypothesis bounds the second term from below by $9^a8^b$ (it is at least its value in the model, of which the
universal module maps onto). With $u_{a+1,b}\le9^{a+1}8^b$, this gives $u_{a,b+1}\le9^{a+1}8^b-9^a8^b=9^a8^{b+1}$. The base case
$u_{a,0}\le9^a$ is the Chari--Loktev bound. Induction on $b$ gives $u_{0,m}\le8^m$. \qed

\begin{physics}
\emph{Heuristic.} A counting argument for the collision process, removing one excitation at a time.
\end{physics}

\begin{status}
Proved; \texttt{\#print axioms} standard (re-checked at an earlier commit, Section~\ref{l27:sec:verification}). Conditional on the
unproved \lean{RavinderKernelLower}.
\end{status}

%======================================================================
\section{The seven final theorems}\label{l27:sec:finals}
%======================================================================

\subsection{How every final theorem is proved}\label{l27:sec:commonproof}

The seven theorems are wrappers: \leanpath{VOA/DeligneAudit/DeligneSeries.lean} and
\leanpath{VOA/DeligneAudit/DeligneSeriesA2.lean} prove no new mathematics, and each proof is a one-line composition of theorems
from Chapters~\ref{ch:L24}--\ref{ch:L26} and Section~\ref{l27:sec:ladders}. Fix a type with Lie algebra $\ltwentyseveng$, VOA
$W=L_k(\ltwentyseveng)$, explicit presentation $P$ and sector data $S$.
\begin{enumerate}
\item \emph{From Input T to the total bound with finiteness.} The degree pieces of every module with the relations of $M_{0,m}$
  are finite-dimensional, proved for each type from a kernel-checked nilpotency certificate (\lean{X_fusionTotalBoundG_of_dim};
  Chapter~\ref{ch:L26}, Section~\ref{l26:sec:finiteG}). Together with Input T this is \lean{FusionTotalBoundG S}.
\item \emph{Ravinder's kernel lemma.} In a module with the relations of $M_{a+1,b}$, the vector $f_\theta(2a+b+1)w$ satisfies
  the relations of $M_{a,b}$; proved for every type with no hypothesis (Chapter~\ref{ch:L26}, Section~\ref{l26:sec:kernelG}). This
  gives the spanning inequality (K) of Example~\ref{l27:ex:sl3} for the universal modules.
\item \emph{The squeeze} (Chapter~\ref{ch:L26}, Section~\ref{l26:sec:squeeze}). Let $F=\mathtt{fusionOf}\ \mathrm{Dem}$ be the family
  defined from the Demazure table by the mixed recursion \emph{with equality}. Input R gives $\mathrm{Dem}_j\le\mathrm{ch}\,M_{j,0}$;
  (K) propagates it downwards to a coefficientwise lower bound $\mathrm{ch}\,M_{0,m}\ge F_m$ (the subtractions are by known upper
  bounds); the \emph{count} $F_m(1)=\ltwentysevendg^{\,m}$ is proved for every type (Chapter~\ref{ch:L25}, Section~\ref{l25:sec:counts});
  and Input T gives $\mathrm{ch}\,M_{0,m}(1)\le\ltwentysevendg^{\,m}$. A coefficientwise lower bound with the same total forces
  $\mathrm{ch}\,M_{0,m}=F_m$, as in Example~\ref{l27:ex:sl3}. This is U5 for $F$, and U6 holds for $F$ by its definition.
\item \emph{The carrier side} (Chapter~\ref{ch:L24}): U5 and U6, with the divisor comparison, give statement (B).
\item \emph{The character side} (Chapter~\ref{ch:L25}): $c_N\le\dim L_N$ for every type; with Input S this is statement (C).
\item \emph{The capstone} (Chapter~\ref{ch:L24}): (B) and (C) give $\mathrm{carrierDim}(N)\le\dim L_N$, and with current-jet
  surjectivity and the geometric carrier, classical freeness and $\ker=(\ltwentysevenIns)$ follow.
\end{enumerate}
For $A_2$ steps 1--3 run in the hand-written $\mathfrak{sl}_3$ setting (Chapter~\ref{ch:L26}, Sections~\ref{l26:sec:a2rel}
to~\ref{l26:sec:u7}), and all three inputs are theorems.

\subsection{The \texorpdfstring{$A_2$}{A2} theorem}

\begin{theorem}[$A_2$: classical freeness of $L_{-3/2}(\mathfrak{sl}_3)$]\label{l27:thm:A2}
Let $W=L_{-3/2}(\mathfrak{sl}_3)$ and let $P$ be its explicit current $C_2$ presentation, with generators the eight currents and
ideal $(\ltwentysevenIns)$ generated by the eight quadrics $P_0,\dots,P_7$ of Definition~\ref{l24:def:A2rel}. Then $W$ is classically
free for $P$, and $\ker(\operatorname{Sym}\mathfrak{sl}_3\to R_W)=(\ltwentysevenIns)$; that is, $R_W=\operatorname{Sym}\mathfrak{sl}_3/(\ltwentysevenIns)$
and $\ltwentysevenJ R_W\to\ltwentysevengr^{\mathrm{Li}}W$ is an isomorphism. The same conclusion is also stated under the (unproved)
hypothesis \lean{RavinderKernelLower}.
\end{theorem}
\noindent File \leanpath{VOA/\allowbreak DeligneAudit/\allowbreak A2Unconditional.lean}.
\begin{leancode}
theorem classicallyFree :
    CaseClassicallyFree .A2 W explicitPresentation ∧
      RingHom.ker (MvPolynomial.aeval fun i =>
        liDegreeZeroClass W 1 ⟨explicitPresentation.gen i,
          explicitPresentation.gen_weight_one i⟩).toRingHom = I2ns :=
  classicallyFree_of_totalDim' CVUniv.fusionTotalDim
\end{leancode}
\noindent File \leanpath{VOA/\allowbreak DeligneAudit/\allowbreak DeligneSeriesA2.lean}.
\begin{leancode}
theorem A2_classicallyFree_final :
    CaseClassicallyFree .A2 VOA.DeligneAudit.A2.W VOA.DeligneAudit.A2.explicitPresentation ∧
      RingHom.ker (MvPolynomial.aeval fun i =>
        liDegreeZeroClass VOA.DeligneAudit.A2.W 1 ⟨VOA.DeligneAudit.A2.explicitPresentation.gen i,
          VOA.DeligneAudit.A2.explicitPresentation.gen_weight_one i⟩).toRingHom =
        VOA.DeligneAudit.A2.I2ns :=
  VOA.DeligneAudit.A2.classicallyFree

/-- The older A₂ S2 interface, retained for reuse. -/
theorem A2_classicallyFree_final_S2 (hK : VOA.DeligneAudit.A2.RavinderKernelLower) :
    CaseClassicallyFree .A2 VOA.DeligneAudit.A2.W VOA.DeligneAudit.A2.explicitPresentation ∧
      RingHom.ker (MvPolynomial.aeval fun i =>
        liDegreeZeroClass VOA.DeligneAudit.A2.W 1 ⟨VOA.DeligneAudit.A2.explicitPresentation.gen i,
          VOA.DeligneAudit.A2.explicitPresentation.gen_weight_one i⟩).toRingHom =
        VOA.DeligneAudit.A2.I2ns :=
  VOA.DeligneAudit.A2.classicallyFree_of_kernelLower' hK
\end{leancode}

\begin{dictionary}
\textbf{Names.} \lean{classicallyFree} is \lean{VOA.DeligneAudit.A2.classicallyFree}; in its file the short names \lean{W},
\lean{explicitPresentation}, \lean{I2ns} are those of the namespace \lean{VOA.DeligneAudit.A2}. The two wrappers are in the namespace
\lean{VOA.DeligneAudit.DeligneSeries} and spell out the full names. \lean{A2_classicallyFree_final} is literally
\lean{A2.classicallyFree}: the proof term is that constant.

\textbf{The objects.} \lean{A2.W} is \lean{AffineSl3.simpleVOA (-3/2) _}, the quotient of the universal affine VOA of $\mathfrak{sl}_3$ at
level $-3/2$ (trace form, $(\theta,\theta)=2$) by its maximal graded ideal; the library proves it simple, nonzero, of central charge
$c=8k/(k+3)=-8$ (Chapter~\ref{ch:L24}, dictionary of Theorem~\ref{l24:thm:sevenpres}). \lean{A2.explicitPresentation} has generators the
eight Chevalley currents $J^a_{(-1)}\ltwentysevenvac$, which form a basis of $W_1$, and ideal \lean{A2.I2ns}
$=$ \lean{Ideal.span (Set.range relationPoly)}, generated by the eight quadrics (Warning ``What the relation lists are'' in
Section~\ref{l27:sec:conclusion}). The conclusion is that of Definition~\ref{l27:def:conclusion} with $I=$ \lean{A2.I2ns}.

\textbf{The hypotheses.} \lean{A2.classicallyFree} and \lean{A2_classicallyFree_final}: \textbf{none}. The statement is closed: it takes
no arguments and the file has no section variables. Internal definitions used in its proof (\lean{bracketTable}, the universal modules
of Chapter~\ref{ch:L26}) cannot weaken it, because the statement does not mention them; only axioms could, and the axiom report is
standard (status box). \lean{A2_classicallyFree_final_S2}: one hypothesis, \lean{RavinderKernelLower}
(Definition~\ref{l27:def:a2alt}); it is the older interface, kept for reuse, and is implied by the unconditional theorem in the sense
that its conclusion is already proved.

\textbf{The proof term.} \lean{classicallyFree_of_totalDim'} (\leanpath{VOA/DeligneAudit/DeligneFusionU7.lean}; Chapter~\ref{ch:L26},
Theorem~\ref{l26:thm:u7final}) proves the conclusion from Input T, with Input R supplied by \lean{EvalModel.demazureRealization_all}.
\lean{CVUniv.fusionTotalDim} is Input T (Chapter~\ref{ch:L26}, Theorem~\ref{l26:thm:cvtotal}).

\textbf{Same as the intended statement?} Yes. Classical freeness in the sense of van Ekeren--Heluani for the genuine simple affine VOA
$L_{-3/2}(\mathfrak{sl}_3)$, with $R_V$ presented by the explicit quadrics. The library proves no identification of
$(\ltwentysevenIns)$ with an ideal from the literature, and none is claimed.
\end{dictionary}

\paragraph{How $A_2$ was closed.} The $A_2$ chain of Chapter~\ref{ch:L26} reduced the theorem to Input T,
\lean{SectorWeyl.FusionTotalDim}, with Inputs R and S already proved. Input T is the theorem \lean{CVUniv.fusionTotalDim}: by
\lean{CVUniv.fusionTotalDim_of_cvBases} it follows from two families of dimension bounds for universal Chari--Venkatesh modules of
$\mathfrak{sl}_3$. The rectangular family (\lean{SharpRectangularBases}) asserts, for every $\nu,\rho$ and every
$\varepsilon\in\{(0,0),(1,0),(0,1)\}$, with $j=0$, that the module is finite-dimensional with
$\dim\le3^{|\varepsilon|}6^{|\nu|}10^{|\rho|}$. The residual family (\lean{ResidualTheta3Bound}) asserts, for every $\rho$ and with
$j=1$, $\varepsilon=\nu=0$, that the module is finite-dimensional with $\dim\le8\cdot10^{|\rho|}$. Both families are proved, \lean{CVUniv.sharpRectangularBases} and \lean{CVUniv.residualTheta3Bound} (Chapter~\ref{ch:L26},
Definition~\ref{l26:def:cvfamily} and Theorem~\ref{l26:thm:cvtotal}). \lean{A2.classicallyFree} is \lean{classicallyFree_of_totalDim'}
applied to that term.

\textbf{Proof on paper.} Section~\ref{l27:sec:commonproof} with all three inputs proved: Input T by the Chari--Venkatesh induction
(Chapter~\ref{ch:L26}), Input R by the evaluation model, and Input S because the Demazure sum equals $[q^N]\prod_{r\ge1}(1+q^r)^8$
(the star identity) and this product solves the $h^\vee=3$ Arakawa--Kawasetsu recurrence, while the vacuum quartic relation gives
$c_N\le\dim W_N$ (Chapter~\ref{ch:L25}). \qed

\begin{physics}
$L_{-3/2}(\mathfrak{sl}_3)$ is the chiral algebra of the Argyres--Douglas theory usually called $H_2$ (or $(A_1,D_4)$), whose Higgs
branch is the moduli space of one $SU(3)$ instanton. Granting the four-dimensional/two-dimensional correspondence of Beem et al.\
(a physics input, not part of the Lean statement), the theorem says that the Schur index of $H_2$ equals the Hilbert series of the
arc space of the scheme $\operatorname{Spec}\operatorname{Sym}\mathfrak{sl}_3/(\ltwentysevenIns)$. That scheme is non-reduced; its reduced
variety is the minimal nilpotent orbit closure of $\mathfrak{sl}_3$ (Arakawa--Moreau; not in Lean). Independently, the library proves $\dim W_N=c_N$ (\lean{A2.finrank_eq_mldeCoeff'}) and that the coefficients of
$\prod_{r\ge1}(1+q^r)^8$ satisfy the same recurrence; the composite $\dim W_N=[q^N]\prod_{r\ge1}(1+q^r)^8$ is derivable but not stated
(Chapter~\ref{ch:L25}, Remark~\ref{l25:rem:a2cons}). That $\prod(1+q^r)^8$ is the character of eight free fermions is a statement about
the $q$-series, not about the theory.
\end{physics}

\begin{status}
\lean{A2.classicallyFree} and \lean{A2_classicallyFree_final}: \textbf{no hypotheses}. Machine-confirmed on 2026-10-28: the aggregate build \texttt{bash scripts/campaign\_build.sh VOA} (the full \texttt{lake build VOA} under the project's memory-capped wrapper) was green on 2026-10-28 (a pre-release build): 18{,}901 jobs, and again 18{,}907 jobs after merging the then-current \texttt{origin/main} (the aggregate target contains the umbrella \lean{VOA.Deligne}); \texttt{\#print axioms} (via \texttt{lake env lean
-DmaxSynthPendingDepth=3}) on \lean{VOA.DeligneAudit.A2.classicallyFree} and
\lean{VOA.DeligneAudit.DeligneSeries.A2_classicallyFree_final} reports \lean{[propext, Classical.choice, Quot.sound]}. Earlier record (27 October 2026): targeted builds of \lean{VOA.DeligneAudit.A2Unconditional} (9755 jobs) and
\lean{DeligneSeriesA2} (9756 jobs) green, standard axioms. \lean{A2_classicallyFree_final_S2}: conditional on
\lean{RavinderKernelLower}; axioms standard (re-checked on a verbatim copy at an earlier commit).
\end{status}

\subsection{The six conditional theorems}

\begin{theorem}[$D_4$: the final theorem]\label{l27:thm:D4}
Let $W=L_{-2}(\mathfrak{so}_8)$, built from the certificate-basis Lie data, with its explicit presentation by $105$ quadrics. Suppose
Input T ($\dim\ltwentysevenF_m\le28^m$) and Input R for the unclamped $D_4$ Demazure series. Then $W$ is classically free for its
presentation and $\ker$ is the ideal generated by the $105$ quadrics.
\end{theorem}
\noindent File \leanpath{VOA/\allowbreak DeligneAudit/\allowbreak DeligneSeries.lean}.
\begin{leancode}
theorem D4_classicallyFree_final (hF : VOA.DeligneAudit.FiniteG.FusionTotalDimG D4.D4sector)
    (hD : ∀ n, DemazureRealizationG D4.D4sector (ravFam D4.D4sector)
      VOA.DeligneAudit.DeligneCharacter.D4Wp.hilbSeriesZ n) :
    CaseClassicallyFree .D4 VOA.DeligneAudit.D4Lie.W4 VOA.DeligneAudit.D4Lie.explicitPresentation4 ∧
      RingHom.ker (MvPolynomial.aeval fun i =>
        liDegreeZeroClass VOA.DeligneAudit.D4Lie.W4 1
          ⟨VOA.DeligneAudit.D4Lie.explicitPresentation4.gen i,
            VOA.DeligneAudit.D4Lie.explicitPresentation4.gen_weight_one i⟩).toRingHom =
        VOA.DeligneAudit.D4Lie.explicitPresentation4.ideal :=
  d4_classicallyFree_of_totalBound (VOA.DeligneAudit.FiniteG.D4.d4_fusionTotalBoundG_of_dim hF) hD
\end{leancode}

\begin{dictionary}
\textbf{The objects.} \lean{D4Lie.W4} is \lean{simpleVOA d4CertData (-2) _} (Warning~\ref{l27:warn:scope}).
\lean{explicitPresentation4.ideal} is \lean{DeligneCarrier.D4.D4data.I}, generated by the $105$ certified quadrics
(Theorem~\ref{l24:thm:sevenpres}). \lean{D4.D4sector} resolves to \lean{DeligneSector.D4.D4sector} through the \lean{open} line of the
file, with \lean{D4sector.n} $=28$. \lean{D4Wp.hilbSeriesZ n} is the unclamped $D_4$ Demazure series.

\textbf{The hypotheses.} \lean{hF}: Input T for $D_4$. \lean{hD}: Input R, for all $n$, with \lean{spec = ravFam D4sector}. Nothing
else: there is no star hypothesis, because the $D_4$ star identity is proved (\lean{D4Wp.d4SchurEqE4}, Theorem~\ref{l25:thm:d4}), and
no positivity hypothesis, because of the unclamped route (dictionary of Definition~\ref{l27:def:realization}).

\textbf{Same as the intended statement?} Yes, for the certificate-basis model of $L_{-2}(\mathfrak{so}_8)$ (Warning~\ref{l27:warn:scope}).
\end{dictionary}

\textbf{Proof on paper.} Section~\ref{l27:sec:commonproof}. Steps 1--3 are \lean{FiniteG.D4.d4_fusionTotalBoundG_of_dim} and
\lean{UnivG.d4_classicallyFree_of_totalBound} (\leanpath{VOA/DeligneAudit/D4FusionKernelTheta.lean}); steps 4--6 are the $D_4$ headline
of Section~\ref{l27:sec:ladders} (Theorem~\ref{l27:thm:D4ladder}) with the star identity, the Ramanujan identities, the divisor
comparison and the vacuum quartic relation plugged in as theorems. \qed

\begin{physics}
$L_{-2}(\mathfrak{so}_8)$ is the chiral algebra of $SU(2)$ gauge theory with four fundamental hypermultiplets, with $SO(8)$ flavour
symmetry. \emph{Heuristic:} its Higgs branch is the one-$SO(8)$-instanton moduli space, the minimal nilpotent orbit closure, the reduced
variety of the scheme defined by the $105$ quadrics. Given the two inputs, the theorem says the Schur index (via the 4d/2d
correspondence) is the arc-space Hilbert series of that scheme; its identification with $E_4'/(240\eta^{10})$ in Chapter~\ref{ch:L25}
is a statement about a $q$-series whose reading as the gauge-theory Schur index is modelling, not Lean.
\end{physics}

\begin{status}
Proved, \textbf{conditional} on \lean{FusionTotalDimG D4sector} and \lean{∀ n, DemazureRealizationG …}. \texttt{\#print axioms}
standard (campaign record; the two theorems it composes were re-checked at an earlier commit). Not unconditional. This is the
general form, with both inputs as hypotheses. The second hypothesis is now a theorem
(\lean{D4Char.d4_demazureRealization}, Theorem~\ref{l27:thm:d4R}), and Theorem~\ref{l27:thm:d4inputT} below is this theorem with it
plugged in: $D_4$ classical freeness is conditional on Input T alone.
\end{status}

\begin{theorem}[$D_4$: classical freeness from Input T alone]\label{l27:thm:d4inputT}
Let $W=L_{-2}(\mathfrak{so}_8)$, built from the certificate-basis Lie data, with its explicit presentation by $105$ quadrics. Suppose
Input T ($\dim\ltwentysevenF_m\le28^m$). Then $W$ is classically free for its presentation and $\ker$ is the ideal generated by the
$105$ quadrics.
\end{theorem}
\noindent File \leanpath{VOA/\allowbreak DeligneAudit/\allowbreak DeligneSeriesD4.lean}.
\begin{leancode}
theorem D4_classicallyFree_of_inputT (hF : VOA.DeligneAudit.FiniteG.FusionTotalDimG D4.D4sector) :
    CaseClassicallyFree .D4 VOA.DeligneAudit.D4Lie.W4 VOA.DeligneAudit.D4Lie.explicitPresentation4 ∧
      RingHom.ker (MvPolynomial.aeval fun i =>
        liDegreeZeroClass VOA.DeligneAudit.D4Lie.W4 1
          ⟨VOA.DeligneAudit.D4Lie.explicitPresentation4.gen i,
            VOA.DeligneAudit.D4Lie.explicitPresentation4.gen_weight_one i⟩).toRingHom =
        VOA.DeligneAudit.D4Lie.explicitPresentation4.ideal :=
\end{leancode}

\begin{dictionary}
\textbf{What it is.} The full name is \lean{VOA.DeligneAudit.DeligneSeries.D4_classicallyFree_of_inputT}. Its conclusion is word for
word that of \lean{D4_classicallyFree_final} (Theorem~\ref{l27:thm:D4}); the objects are the same. Its proof is one line: it applies
\lean{D4_classicallyFree_final} to \lean{hF} and to \lean{D4Char.d4_demazureRealization} (Theorem~\ref{l27:thm:d4R}) in place of
\lean{hD}.

\textbf{The hypothesis.} \lean{hF}: Input T for $D_4$, the bound $\dim W(m\theta,m)\le28^m$ for every $m$ (Definition~\ref{l27:def:totaldim}).
Nothing else: Input R is proved, the star identity is proved, and the unclamped route needs no positivity hypothesis.

\textbf{Same as the intended statement?} Yes, for the certificate-basis model of $L_{-2}(\mathfrak{so}_8)$ (Warning~\ref{l27:warn:scope}).
\end{dictionary}

\begin{status}
Proved, \textbf{conditional} on \lean{hF : FusionTotalDimG D4sector} (Input T). Machine-confirmed: the aggregate build was green
on 3 November 2026 (20{,}331 jobs), and \texttt{\#print axioms} reports \lean{[propext, Classical.choice, Quot.sound]}. Not
unconditional. Input T is a published theorem (Kus--Littelmann, Theorem~4.2), with the caveat of item~(i): one step of the
published proof is insufficient as written for $D_4$ at $m=2$, by a computational witness; the theorem is not known to fail.
In Lean, Input T for $D_4$ is proved for $m\le2$ and reduced, for all $m$, to $\mathfrak{sl}_4$ core bounds that remain open
(Section~\ref{l27:sec:partial} and the end of Section~\ref{l27:sec:feigin}); that progress is not part of the claim.
\end{status}

\begin{physics}
$L_{-2}(\mathfrak{so}_8)$ is the chiral algebra of four-dimensional $\mathcal N=2$ $SU(2)$ gauge theory with four
flavours, whose Higgs branch is the minimal nilpotent orbit of $\mathfrak{so}_8$. Classical freeness says that its
Schur operators are exactly the derivatives of the Higgs-branch chiral ring: the Schur index is the arc-space
Hilbert series of the $105$ moment-map relations. The theorem proves this assuming only a counting bound on
current-algebra modules (Input~T).
\end{physics}

\begin{theorem}[$G_2$ and $F_4$: the final theorems]\label{l27:thm:G2F4}
Let $X\in\{G_2,F_4\}$, $W=L_k(\ltwentyseveng_X)$ ($k=-5/3$, resp.\ $-5/2$) built from the certified structure-constant Lie data, with
its explicit presentation ($27$, resp.\ $324$ quadrics). Suppose Input T ($\ltwentysevendg=14$, resp.\ $52$), Input R (unclamped table for
$G_2$, clamped for $F_4$) and Input S ($h^\vee=4$, resp.\ $9$). Then $W$ is classically free for its presentation and $\ker$ is the ideal
generated by the listed quadrics.
\end{theorem}
\noindent File \leanpath{VOA/\allowbreak DeligneAudit/\allowbreak DeligneSeries.lean}.
\begin{leancode}
theorem G2_classicallyFree_final (hF : VOA.DeligneAudit.FiniteG.FusionTotalDimG G2.G2sector)
    (hD : ∀ n, DemazureRealizationG G2.G2sector (ravFam G2.G2sector)
      VOA.DeligneAudit.DeligneCarrier.g2HilbSeriesZ n)
    (hB : VOA.DeligneAudit.DeligneCharacter.G2StarBound) :
    CaseClassicallyFree .G2 VOA.DeligneAudit.G2Lie.g2W VOA.DeligneAudit.G2Lie.explicitPresentation ∧
      RingHom.ker (MvPolynomial.aeval fun i =>
        liDegreeZeroClass VOA.DeligneAudit.G2Lie.g2W 1
          ⟨VOA.DeligneAudit.G2Lie.explicitPresentation.gen i,
            VOA.DeligneAudit.G2Lie.explicitPresentation.gen_weight_one i⟩).toRingHom =
        VOA.DeligneAudit.G2Lie.explicitPresentation.ideal :=
  g2_classicallyFree_of_totalBound (VOA.DeligneAudit.FiniteG.G2.g2_fusionTotalBoundG_of_dim hF) hD hB
\end{leancode}
\begin{leancode}
theorem F4_classicallyFree_final (hF : VOA.DeligneAudit.FiniteG.FusionTotalDimG F4.F4sector)
    (hD : ∀ n, DemazureRealizationG F4.F4sector (ravFam F4.F4sector)
      (VOA.DeligneAudit.DeligneCarrier.hilbSeries VOA.DeligneAudit.DeligneCharacter.f4DemazureHilb) n)
    (hB : VOA.DeligneAudit.DeligneCharacter.F4StarBound) :
    CaseClassicallyFree .F4 VOA.DeligneAudit.F4Lie.f4W VOA.DeligneAudit.F4Lie.explicitPresentation ∧
      RingHom.ker (MvPolynomial.aeval fun i =>
        liDegreeZeroClass VOA.DeligneAudit.F4Lie.f4W 1
          ⟨VOA.DeligneAudit.F4Lie.explicitPresentation.gen i,
            VOA.DeligneAudit.F4Lie.explicitPresentation.gen_weight_one i⟩).toRingHom =
        VOA.DeligneAudit.F4Lie.explicitPresentation.ideal :=
\end{leancode}

\begin{dictionary}
\textbf{The objects.} \lean{G2Lie.g2W} and \lean{F4Lie.f4W} are \lean{simpleVOA} of \lean{g2CertData} at $-5/3$ and of \lean{f4CertData}
at $-5/2$ (Warning~\ref{l27:warn:scope}). \lean{explicitPresentation.ideal} is \lean{G2data.I}, resp.\ \lean{F4data.I}. \lean{G2.G2sector},
\lean{F4.F4sector} are the sector data (namespace \lean{DeligneSector}).

\textbf{The hypotheses.} \lean{hF}: Input T. \lean{hD}: Input R against \lean{g2HilbSeriesZ} (unclamped) for $G_2$ and against
\lean{hilbSeries f4DemazureHilb} (clamped) for $F_4$. \lean{hB}: Input S, \lean{G2StarBound}, resp.\ \lean{F4StarBound}. For $G_2$,
\lean{hD} and \lean{hB} refer to the unclamped and the clamped table respectively; inside the proof U5 is derived from \lean{hF} and
\lean{hD}, nonnegativity of the unclamped table follows (\lean{g2HilbInt_nonneg_of_U5Z}), and then the two tables agree
(\lean{g2FusionF_eq_Z}); see the dictionary of Definition~\ref{l27:def:realization}.

\textbf{Non-simply-laced types and the literature.} $G_2$ and $F_4$ are not simply laced, but this does not affect the literature
status of their inputs. The header of \leanpath{VOA/DeligneAudit/F4LieHeadline.lean} says the Ravinder/Chari--Venkatesh route to U5 is
``not literature-verified'' for $F_4$. That comment misplaces the issue: Ravinder's Theorem~1, Corollaries~1 and~2 and
Proposition~1 are stated for every simple $\ltwentyseveng$ (only his Theorem~2 and Corollary~3 assume simply laced), and
Kus--Littelmann's Theorem~4.2 with $\lambda_0=0$ has no type restriction. The live issues are the same as for every other type:
Feigin's gap, and the flagged step of the Kus--Littelmann repair (Input T); the unformalized Demazure character formula
(Input R); and Input S, which is conjectural. The short roots need no separate treatment (dictionary of
Definition~\ref{l27:def:relspec}). The Lean statements are unaffected in any case: the literature only bears on whether the
hypotheses are believed true.
\end{dictionary}

\textbf{Proof on paper.} Section~\ref{l27:sec:commonproof}. For $G_2$ steps 1--3 are \lean{g2_fusionTotalBoundG_of_dim} and
\lean{UnivG.g2_classicallyFree_of_totalBound} (\leanpath{VOA/DeligneAudit/G2FusionKernelTheta.lean}), whose last step is the concrete
headline \lean{G2Lie.g2_classicallyFree_concrete} (Theorem~\ref{l27:thm:G2F4ladder}); for $F_4$ they are \lean{f4_fusionTotalBoundG_of_dim}
and \lean{DeligneCarrier.f4_classicallyFree_of_totalBound'} (\leanpath{VOA/DeligneAudit/F4FusionCountZ.lean}). The target bound
$c_N\le\dim L_N$ is \lean{g2_targetLowerBound} and \lean{f4_targetLowerBound} (Chapter~\ref{ch:L25}); \lean{hB} closes (C). \qed

\begin{physics}
\emph{Heuristic.} The levels $-5/3$ and $-5/2$ continue the Deligne pattern $k=-1-h^\vee/6$, but no four-dimensional $\mathcal N=2$
SCFT with $G_2$ or $F_4$ flavour symmetry and these central charges is known; they are ``would-be'' theories. The VOAs exist regardless;
their associated varieties are again minimal nilpotent orbit closures (Arakawa--Moreau, for the Deligne series at these levels), and
the theorem says the vacuum character is an arc-space Hilbert series.
\end{physics}

\begin{status}
Proved, \textbf{conditional} on three named inputs each (\lean{hF}, \lean{hD}, \lean{hB}). \texttt{\#print axioms} standard (campaign
record; the composed theorems re-checked at an earlier commit). Not unconditional. Since \lean{hB} is conjectural
(Conjecture~\ref{l27:conj:star}), these theorems are conditional on a conjecture.
\end{status}

\begin{theorem}[$E_6$, $E_7$, $E_8$: the final theorems]\label{l27:thm:E}
Let $X\in\{E_6,E_7,E_8\}$, $W=L_k(\ltwentyseveng_X)$ ($k=-3,-4,-6$) with its explicit presentation ($650$, $1539$, $3875$ quadrics).
Suppose Input T ($\ltwentysevendg=78,133,248$), Input R for the clamped table, and Input S ($h^\vee=12,18,30$). Then $W$ is classically
free for its presentation and $\ker$ is the ideal generated by the listed quadrics.
\end{theorem}
\noindent File \leanpath{VOA/\allowbreak DeligneAudit/\allowbreak DeligneSeries.lean}.
\begin{leancode}
theorem E6_classicallyFree_final
    (hF : VOA.DeligneAudit.FiniteG.FusionTotalDimG VOA.DeligneAudit.DeligneCarrier.E6.E6sector)
    (hD : ∀ n, DemazureRealizationG VOA.DeligneAudit.DeligneCarrier.E6.E6sector
      (ravFam VOA.DeligneAudit.DeligneCarrier.E6.E6sector)
      (VOA.DeligneAudit.DeligneCarrier.hilbSeries VOA.DeligneAudit.DeligneCharacter.e6DemazureHilb) n)
    (hB : VOA.DeligneAudit.DeligneCharacter.E6StarBound) :
    CaseClassicallyFree .E6 VOA.DeligneAudit.DeligneCarrier.E6.W6
        VOA.DeligneAudit.DeligneCarrier.E6.explicitPresentation ∧
      RingHom.ker (MvPolynomial.aeval fun i =>
        liDegreeZeroClass VOA.DeligneAudit.DeligneCarrier.E6.W6 1
          ⟨VOA.DeligneAudit.DeligneCarrier.E6.explicitPresentation.gen i,
            VOA.DeligneAudit.DeligneCarrier.E6.explicitPresentation.gen_weight_one i⟩).toRingHom =
        VOA.DeligneAudit.DeligneCarrier.E6.explicitPresentation.ideal :=
\end{leancode}
\begin{leancode}
theorem E7_classicallyFree_final
    (hF : VOA.DeligneAudit.FiniteG.FusionTotalDimG VOA.DeligneAudit.DeligneCarrier.E7.E7sector)
    (hD : ∀ n, DemazureRealizationG VOA.DeligneAudit.DeligneCarrier.E7.E7sector
      (ravFam VOA.DeligneAudit.DeligneCarrier.E7.E7sector)
      (VOA.DeligneAudit.DeligneCarrier.hilbSeries VOA.DeligneAudit.DeligneCharacter.e7DemazureHilb) n)
    (hB : VOA.DeligneAudit.DeligneCharacter.E7StarBound) :
    CaseClassicallyFree .E7 VOA.DeligneAudit.E7Lie.W7 VOA.DeligneAudit.E7Lie.explicitPresentation7 ∧
      RingHom.ker (MvPolynomial.aeval fun i =>
        liDegreeZeroClass VOA.DeligneAudit.E7Lie.W7 1
          ⟨VOA.DeligneAudit.E7Lie.explicitPresentation7.gen i,
            VOA.DeligneAudit.E7Lie.explicitPresentation7.gen_weight_one i⟩).toRingHom =
        VOA.DeligneAudit.E7Lie.explicitPresentation7.ideal :=
\end{leancode}
\begin{leancode}
theorem E8_classicallyFree_final
    (hF : VOA.DeligneAudit.FiniteG.FusionTotalDimG VOA.DeligneAudit.DeligneCarrier.E8.E8sector)
    (hD : ∀ n, DemazureRealizationG VOA.DeligneAudit.DeligneCarrier.E8.E8sector
      (ravFam VOA.DeligneAudit.DeligneCarrier.E8.E8sector)
      (VOA.DeligneAudit.DeligneCarrier.hilbSeries VOA.DeligneAudit.DeligneCharacter.e8DemazureHilb) n)
    (hB : VOA.DeligneAudit.DeligneCharacter.E8StarBound) :
    CaseClassicallyFree .E8 VOA.DeligneAudit.E8Lie.W8 VOA.DeligneAudit.E8Lie.explicitPresentation8 ∧
      RingHom.ker (MvPolynomial.aeval fun i =>
        liDegreeZeroClass VOA.DeligneAudit.E8Lie.W8 1
          ⟨VOA.DeligneAudit.E8Lie.explicitPresentation8.gen i,
            VOA.DeligneAudit.E8Lie.explicitPresentation8.gen_weight_one i⟩).toRingHom =
        VOA.DeligneAudit.E8Lie.explicitPresentation8.ideal :=
\end{leancode}

\begin{dictionary}
\textbf{The objects.} \lean{DeligneCarrier.E6.W6} is \lean{simpleVOA e6LieData (-3) _} (structure-constant model); \lean{E7Lie.W7} and
\lean{E8Lie.W8} are \lean{simpleVOA} of \lean{e7CertData}, \lean{e8CertData} at $-4$, $-6$ (Warning~\ref{l27:warn:scope}). The sector data
\lean{E6sector}, \lean{E7sector}, \lean{E8sector} live in \lean{DeligneCarrier}; their fusion sets are exactly Ravinder's. The ideals are
\lean{E6data.I}, \lean{E7data.I}, \lean{E8data.I}.

\textbf{The hypotheses.} Three each: Input T (\lean{hF}), Input R against the clamped tables (\lean{hD}), and Input S (\lean{hB}).

\textbf{Proof terms.} Each omitted proof is
\lean{eX_classicallyFree_of_starBound (eX_fusionModuleBound_of_totalBound' (eX_fusionTotalBoundG_of_dim hF) hD) hB}, with the first
theorem from \leanpath{VOA/DeligneAudit/E6Final.lean} (resp.\ \texttt{E7Final}, \texttt{E8Final}) and the second from
\leanpath{VOA/DeligneAudit/DeligneCharacterEFusionCount.lean} ($E_6$, $E_7$) or \leanpath{VOA/DeligneAudit/DeligneCharacterE8FusionCount.lean}
($E_8$).
\end{dictionary}

\textbf{Proof on paper.} Section~\ref{l27:sec:commonproof}; step 6 is Theorem~\ref{l27:thm:Eladder}. For the $E$-types the divisor
comparison of step 4 is the largest certificate computation of the project (Chapter~\ref{ch:L24}, Section~\ref{l24:sec:certs}). \qed

\begin{physics}
$L_{-3}(\mathfrak e_6)$, $L_{-4}(\mathfrak e_7)$ and $L_{-6}(\mathfrak e_8)$ are the chiral algebras of the Minahan--Nemeschansky rank-one
theories. \emph{Heuristic:} their Higgs branches are the one-instanton moduli spaces of $E_6$, $E_7$, $E_8$. Given the three inputs, the
theorems say that the Schur indices of these non-Lagrangian theories (via the 4d/2d correspondence) are arc-space Hilbert series of the
schemes defined by the quadrics, and equal the MLDE solutions.
\end{physics}

\begin{status}
Proved, \textbf{conditional} on three named inputs each. \texttt{\#print axioms} standard: campaign record only for the three final
theorems and their $E_7$, $E_8$ ingredients; \lean{E6.e6_classicallyFree_of_starBound} was re-checked at an earlier commit. Not
unconditional. Since \lean{hB} is conjectural (Conjecture~\ref{l27:conj:star}), these theorems are conditional on a conjecture.
\end{status}

%======================================================================
\section{The input ledger, type by type}\label{l27:sec:ledger}
%======================================================================

The long form of the table at the start of the chapter. ``Hypothesis'' means: not proved in Lean, and passed as an explicit argument
of the final theorem. Table~\ref{l27:tab:precise} states, type by type, the precise mathematical content of each input, its status in
the literature, and what the library has proved toward it. Abbreviations: \textbf{KL} = Kus--Littelmann, Pacific J.\ Math.\ 278 (2015),
Theorem~4.2 at $\lambda=\theta$, $\lambda_0=0$, $N=m$ (through their Theorem~3.5 and Corollary~3.6(3)); \textbf{KL-step} = the
annihilator-transfer step on p.~437 of KL, flagged by the project as unjustified as written, and for $D_4$ at $m=2$ shown by
computation to be insufficient as written (no counterexample to the theorem is known);
\textbf{F08} = Feigin, \emph{SIGMA} 4 (2008) 070, Proposition~1 and Corollary~2 (proof has a gap); \textbf{R14} = Ravinder,
\emph{SIGMA} 10 (2014) 110; \textbf{DCF} = the Demazure character formula (Kumar 1987, Mathieu 1988); \textbf{FK} = Frenkel--Kac
(1980). $W(m\theta,m)$ is the truncated Weyl module, equal to the Lean universal module by Lemma~\ref{l27:lem:truncweyl}.
The status ``PROVED [ref]'' refers to the literature, never to Lean.
\begin{center}\footnotesize
\begin{longtable}{@{}p{0.05\textwidth}p{0.035\textwidth}p{0.25\textwidth}p{0.26\textwidth}p{0.29\textwidth}@{}}
\caption{The inputs of the six conditional final theorems: precise statement, literature status, library progress.}\label{l27:tab:precise}\\
\toprule
type & in. & precise statement & literature status & library progress (Lean)\\
\midrule\endfirsthead
\toprule
type & in. & precise statement & literature status & library progress (Lean)\\
\midrule\endhead
$G_2$ & T & $\dim W(m\theta,m)\le14^m$ for all $m$ (\lean{FusionTotalDimG G2sector}; fusion set: all negative roots) &
  PROVED with flagged step: KL (KL-step); earlier F08 (gap); R14 Cor.~2 at $a=0$ is F08, not independent &
  finiteness clause only (\lean{g2_fusionTotalBoundG_of_dim})\\
$G_2$ & R & for all $n$, a module with $\mathrm{rav}(n,2n)$ whose degree-$N$ piece has dimension $\ge[q^N]$ of the unclamped
  $G_2$ table &
  PROVED: R14 Prop.~1 (after Chari--Venkatesh Cor.~3.5; short roots included) $+$ DCF; unformalized: the Lean algorithm $=$ DCF &
  none (no lattice model)\\
$G_2$ & S & $\mathcal D_N\le c_N$ for all $N$, $h^\vee=4$ (\lean{G2StarBound}) &
  CONJECTURAL: no published source known; numerics through $q^{60}$ &
  \lean{G2KWEqMLDE} proved; S follows from the single named identity \lean{G2DemazureEqKW}, split as \lean{G2DemazureFermionic}
  (follows from published results) and \lean{G2FermionicKW} (open)\\
\midrule
$D_4$ & T & $\dim W(m\theta,m)\le28^m$ for all $m$ (fusion set: all negative roots) &
  PROVED with flagged step: KL (KL-step, insufficient as written at $m=2$ by computation, witness $\tilde a$; theorem not known to fail); F08 (gap) &
  $m=1,2$ proved (\lean{localWeyl_one_total_dim}, \lean{localWeyl_two_total_dim}: $\le28$, $\le784$); all-$m$ layers in progress
  (Section~\ref{l27:sec:partial}); presented-module route (\lean{D4_classicallyFree_of_presentedRankBoundE}): $m=1,2$
  proved; level route (\lean{d4_classicallyFree_of_twoAdjointCoreBounds}): reduced to $\mathfrak{sl}_4$ core bounds with $a\ge2$
  adjoint sites, and $a=2$ reduced on paper to $\tilde Q_\kappa\le10\cdot6^\kappa$ and $(C)_{\kappa+1}$ (end of
  Section~\ref{l27:sec:feigin}); all $m$ open. This is the one hypothesis of \lean{D4_classicallyFree_of_inputT}\\
$D_4$ & R & for all $n$, a module with $\mathrm{rav}(n,2n)$ and graded dimension $\ge$ the unclamped $D_4$ table &
  PROVED: R14 Prop.~1 $+$ DCF; FK realizes $D(1,n\theta)$ in $V_{D_4}$ &
  $n=0,1,2$ proved (\lean{D4Small.}\allowbreak\lean{demazureRealization_}\allowbreak\lean{zero}, \lean{…_one}, \lean{D4Two.}\allowbreak\lean{demazureRealization_}\allowbreak\lean{two}); relation half
  in the lattice model for all $n$ (\lean{d4_lattice_sat}); \textbf{all $n$ proved}:
  \lean{D4Char.d4_demazureRealization} (Theorem~\ref{l27:thm:d4R}; aggregate build green, 3 November 2026)\\
$D_4$ & S & none: the identity is proved & --- & \lean{D4Wp.d4SchurEqE4} with the unclamped Demazure $=$ Schur identity
  (\lean{d4DemazureEqSchur_of_positivity}, positivity from \lean{D4Wp.d4DemazurePositivity_of_U5Z})\\
\midrule
$F_4$ & T & $\dim W(m\theta,m)\le52^m$ (fusion set: all negative roots) &
  PROVED with flagged step: KL (KL-step); F08 (gap) & finiteness clause only\\
$F_4$ & R & for all $n$, $\mathrm{rav}(n,2n)$ and graded dimension $\ge$ the clamped $F_4$ table (equivalent to the unclamped form) &
  PROVED: R14 Prop.~1 $+$ DCF ($\theta$ long; short roots included) & none (no lattice model)\\
$F_4$ & S & $\mathcal D_N\le c_N$, $h^\vee=9$ (\lean{F4StarBound}) &
  CONJECTURAL: no published source known; numerics through $q^{25}$ & none\\
\midrule
$E_6$ & T & $\dim W(m\theta,m)\le78^m$ (fusion set: exactly Ravinder's) &
  PROVED with flagged step: KL (KL-step); F08 (gap) & finiteness clause only\\
$E_6$ & R & for all $n$, $\mathrm{rav}(n,2n)$ and graded dimension $\ge$ the clamped $E_6$ table &
  PROVED: R14 Prop.~1 $+$ DCF; FK & relation half in the lattice model (\lean{e6_lattice_sat}); dimension half named
  (\lean{E6LatticeDim})\\
$E_6$ & S & $\mathcal D_N\le c_N$, $h^\vee=12$ (\lean{E6StarBound}) &
  CONJECTURAL: no published source known; numerics through $q^{20}$ & none\\
\midrule
$E_7$ & T & $\dim W(m\theta,m)\le133^m$ (Ravinder's fusion set) &
  PROVED with flagged step: KL (KL-step); F08 (gap) & finiteness clause only\\
$E_7$ & R & for all $n$, $\mathrm{rav}(n,2n)$ and graded dimension $\ge$ the clamped $E_7$ table &
  PROVED: R14 Prop.~1 $+$ DCF; FK & relation half in the lattice model (\lean{e7_lattice_sat}); dimension half named\\
$E_7$ & S & $\mathcal D_N\le c_N$, $h^\vee=18$ (\lean{E7StarBound}) &
  CONJECTURAL: no published source known; numerics through $q^{16}$ & none\\
\midrule
$E_8$ & T & $\dim W(m\theta,m)\le248^m$ (Ravinder's fusion set) &
  PROVED with flagged step: KL (KL-step); F08 (gap) & finiteness clause only\\
$E_8$ & R & for all $n$, $\mathrm{rav}(n,2n)$ and graded dimension $\ge$ the clamped $E_8$ table &
  PROVED: R14 Prop.~1 $+$ DCF; FK & relation half in the lattice model (\lean{e8_lattice_sat}); dimension half named\\
$E_8$ & S & $\mathcal D_N\le c_N$, $h^\vee=30$ (\lean{E8StarBound}) &
  CONJECTURAL: no published source known; numerics through $q^{14}$ & none\\
\bottomrule
\end{longtable}
\end{center}
For $A_2$ all three inputs are \textbf{proved in Lean}: T by \lean{CVUniv.fusionTotalDim}, R by
\lean{EvalModel.demazureRealization_all}, S by \lean{A2.starIdentity} (Chapter~\ref{ch:L25}). The older $A_2$ interface
\lean{A2.RavinderKernelLower} (the injectivity half of Ravinder's sequence in a model) is not proved and not needed.

\paragraph{Reading the table.} The final theorems of $G_2$, $F_4$, $E_6$, $E_7$, $E_8$ depend on Input S, which is
\emph{conjectural}: they are theorems conditional on a conjecture. For $D_4$ Input R is proved in Lean
(\lean{D4Char.d4_demazureRealization}), so the $D_4$ theorem \lean{D4_classicallyFree_of_inputT} depends on Input T alone. Input T
is a published theorem (with the flagged KL-step); what is missing for $D_4$ is its formalization and a complete proof of the
KL-step, whose written argument is insufficient at $m=2$ (the statement it is meant to prove is not known to fail). In Lean,
Input T for $D_4$ has the cases $m\le2$ and a reduction for all $m$ (end of Section~\ref{l27:sec:feigin}); it stays a hypothesis.
Everything else in the chain is proved for all seven types: the explicit presentation (as a containment), jet surjectivity, the
geometric carrier, B1, the divisor comparison, the squeeze, Ravinder's kernel lemma, the finiteness clause of Input T, the $q=1$
counts and energies, and the target lower bound $c_N\le\dim L_N$. The numerical checks are evidence, not proof, and they concern the
identities, not their Lean statements. For $G_2$ the star identity is also split into two named halves, checked separately (the
fermionic half for $j\le40$ in all degrees $\le40-j$, and it follows from published results; the $q$-series half through $q^{60}$,
weight-refined through $q^{50}$; Theorem~\ref{l27:thm:G2F4ladder} and Section~\ref{l27:sec:starinput}); the $q^{60}$ figure in the
table is for the combined identity.

\begin{warning}[title={Source citations do not discharge hypotheses}]
A named input with a published proof is still a hypothesis of the Lean theorem. The six conditional final theorems are
\texttt{\#print axioms}-clean, conditional on the inputs above; they are not unconditional, and a citation in a docstring changes
nothing about that.
\end{warning}

\subsection{Partial results toward Inputs T and R for \texorpdfstring{$D_4$}{D4}}\label{l27:sec:partial}

The library proves some cases of the $D_4$ inputs outright. None of them is used by a final theorem (the final theorems need all
$m$ and all $n$), but each is an exact instance of the corresponding hypothesis, so they show that the hypotheses are satisfiable in
the cases covered and that the Lean statements mean what they are said to mean. All are in the files assigned to
Chapter~\ref{ch:L26} and are imported by the \lean{VOA.Deligne} umbrella. Every quotation was checked against the source at
the time of the 28 October revision. Input R for $D_4$ is now proved for every $n$ (Theorem~\ref{l27:thm:d4R}); the cases
$n\le2$ below were proved first and are kept as worked instances.

\begin{theorem}[Input T for $D_4$ at $m=1$ and $m=2$]\label{l27:thm:d4Tsmall}
Let $V$ be a current module for the $D_4$ structure constants and $w\in V$ satisfy the local Weyl and fusion relations of $D_4$ for
$m\theta$. If $m=1$, the cyclic span of $w$ has total dimension at most $28$; if $m=2$, at most $784=28^2$.
\end{theorem}
\noindent File \leanpath{VOA/\allowbreak DeligneAudit/\allowbreak D4FusionContactOne.lean}.
\begin{leancode}
theorem localWeyl_one_total_dim (h : LocalWeylRelG D4sector R w 1) (B : ℕ) :
    ∑ d ∈ Finset.range B, Module.finrank ℂ (cycDeg R w d) ≤ 28 := by
\end{leancode}
\noindent File \leanpath{VOA/\allowbreak DeligneAudit/\allowbreak D4FusionContactTwo.lean}.
\begin{leancode}
theorem localWeyl_two_total_dim (h : LocalWeylRelG D4sector R w 2) (B : ℕ) :
    ∑ d ∈ Finset.range B, Module.finrank ℂ (cycDeg R w d) ≤ 784 :=
\end{leancode}

\begin{dictionary}
Both are in the namespace \lean{VOA.DeligneAudit.D4FusionContact}. The section variables are
\lean{\{V : Type\} [AddCommGroup V] [Module ℂ V]} and \lean{\{R : CurrentModule D4sector.f V\} \{w : V\}}, the same universe and
data over which \lean{FusionTotalDimG} quantifies (Definition~\ref{l27:def:totaldim}); \lean{D4sector} is
\lean{DeligneSector.D4.D4sector}, the sector of the $D_4$ final theorem. So each theorem is literally the body of
\lean{FusionTotalDimG D4sector} at $m=1$, resp.\ $m=2$, with $28^1=28$ and $28^2=784$. The finiteness of the degree pieces, needed to
read the bound as a bound on dimension (dictionary of Definition~\ref{l27:def:totaldim}), is proved alongside
(\lean{localWeyl_one_finite}, \lean{localWeyl_two_finite}). By Lemma~\ref{l27:lem:truncweyl} the statements say
$\dim W(\theta,1)\le28$ and $\dim W(2\theta,2)\le784$; the first is the classical fact $W(\theta,1)=\mathfrak{so}_8$ as a
$\mathfrak{so}_8$-module (the relations kill all positive modes), the second is the first nontrivial case of Kus--Littelmann's
Theorem~4.2 for $D_4$. The docstrings say explicitly that these are not the all-$m$ input.
\end{dictionary}

\textbf{Proof on paper.} For $m=1$ every positive mode kills $w$ (Lemma~\ref{l27:lem:truncweyl} with $\mathcal I_1$), so the
cyclic span is $U(\mathfrak{so}_8)w$, an irreducible module of highest weight $\theta$, the adjoint representation, of dimension
$28$. For $m=2$ the library constructs an explicit finite-dimensional space (\lean{M2Model}) of dimension at most $784$
(\lean{m2Model_bound}) together with a surjective linear map from it onto the universal module with the $m=2$ relations; every module
with the relations is a quotient of the universal one, so its degree pieces sum to at most $784$. The proof is in the two files above
and their imports. \qed

\begin{physics}
\emph{Heuristic.} One adjoint particle of $SO(8)$ flavour carries $28$ states; two of them fused carry at most $28^2$: bringing two
particles together creates no new states.
\end{physics}

\begin{status}
Proved, \textbf{unconditional} (no hypothesis beyond the relations). \texttt{\#print axioms} standard per the campaign record
(entry ``D4 total bound'': the $m=0,1,2$ cases closed, targets green); not re-run by the author of this chapter.
For $m\ge3$ the bound $28^m$ is \textbf{open} in Lean.
\end{status}

\begin{theorem}[Input R for $D_4$ at $n=0,1,2$]\label{l27:thm:d4Rsmall}
For $n=0,1,2$ there is a current module for the $D_4$ structure constants with a vector satisfying the relations
$\mathrm{rav}(n,2n)$ whose degree-$N$ piece has dimension at least $[q^N]$ of the unclamped $D_4$ Demazure series, for every $N$.
\end{theorem}
\noindent File \leanpath{VOA/\allowbreak DeligneAudit/\allowbreak D4FusionRealizationSmall.lean}.
\begin{leancode}
theorem demazureRealization_zero :
    DemazureRealizationG D4sector (ravFam D4sector) D4Wp.hilbSeriesZ 0 := by
\end{leancode}
\begin{leancode}
theorem demazureRealization_one :
    DemazureRealizationG D4sector (ravFam D4sector) D4Wp.hilbSeriesZ 1 := by
\end{leancode}
\noindent File \leanpath{VOA/\allowbreak DeligneAudit/\allowbreak D4FusionRealizationTwoRealization.lean}.
\begin{leancode}
theorem demazureRealization_two :
    DemazureRealizationG D4sector (ravFam D4sector) D4Wp.hilbSeriesZ 2 := by
\end{leancode}

\begin{dictionary}
The first two are \lean{VOA.DeligneAudit.D4Small.demazureRealization_zero} and \lean{…_one} (not to be confused with the $A_2$
\lean{demazureRealization_zero} of \leanpath{VOA/DeligneAudit/DeligneFusionDemazure.lean}); the third is
\lean{VOA.DeligneAudit.D4Two.demazureRealization_two}. The type is exactly the hypothesis \lean{hD} of the $D_4$ final theorem
(Theorem~\ref{l27:thm:D4}) at \lean{n}$\,=0,1,2$: same sector, same family \lean{ravFam}, same unclamped series
\lean{D4Wp.hilbSeriesZ}. None takes an argument.
\end{dictionary}

\textbf{Proof on paper.} $n=0$: the trivial one-dimensional module, with $\mathrm{Dem}_0=1$. $n=1$: the module $D(1,\theta)=
\mathfrak{so}_8\oplus\C$ of the primer (Section~\ref{l27:sec:fusionprimer}), with character $28+q$, built as an explicit one-site
module. $n=2$: the library builds a cyclic graded module from an explicit source representation, checks Ravinder's relations on its
generator, and computes its graded dimension degree by degree; it matches every coefficient of the unclamped Demazure series at
$n=2$ (total $29^2=841$). \qed

\begin{physics}
\emph{Heuristic.} Explicit systems of zero, one and two excitable adjoint particles of $SO(8)$ whose energy spectra are exactly as rich
as the level-one WZW vacuum module predicts.
\end{physics}

\begin{status}
Proved, \textbf{unconditional}. \texttt{\#print axioms} standard per the campaign record (entries ``D4 ORIGINAL n=2
REALIZATION CLOSED'' and ``D4 realizations: actual cases 0,1,2 are now closed''); not re-run by the author of this chapter. For
$n\ge3$ a total-dimension count $29^n$ alone would not suffice, since the squeeze needs the graded inequality; the statement for
every $n$ is \lean{D4Char.d4_demazureRealization} (Theorem~\ref{l27:thm:d4R}), which is proved and confirmed by the aggregate build.
\end{status}

\paragraph{Work in progress toward all $m$ (Input T for $D_4$).} The library contains research layers toward the bound $28^m$ for every
$m$, in the files \leanpath{VOA/DeligneAudit/D4FusionContactAllMCountResearchCubeCurrentOrders.lean} and
\leanpath{VOA/DeligneAudit/D4FusionContactAllMCountResearchCubeQuinticLeading.lean}. What Lean states there
(\lean{raw_order_chain_counts_mem_ann}, \lean{quinticHead_mem_ann_sup_low}) is only this: for every $m\ge1$ and every module with
\lean{LocalWeylRelG D4sector R x m}, certain explicit coefficients of products of currents lie in the annihilator ideal of the generator
modulo an explicit span of lower ``colour'' layers. These are steps of a reduction, not a dimension bound; no Lean statement in these
files bounds a dimension for $m\ge3$, and we describe them no further.

%======================================================================
\section{The inputs that remain open in Lean}\label{l27:sec:cores}
%======================================================================

For $G_2$, $D_4$, $F_4$, $E_6$, $E_7$, $E_8$ the remaining hypotheses are of two different kinds, and the distinction matters. Inputs T
and R are \emph{published theorems} (T with one proof step flagged by the project); the library has formalized R for $D_4$ but
neither input for the other five types, nor T for $D_4$, which is the only remaining hypothesis of the $D_4$ theorem. Input S is,
as far as we know, an \emph{open mathematical problem}: a conjecture with numerical evidence. This section says what is known about
each, and what made the $A_2$ case tractable.

\subsection{Input T beyond \texorpdfstring{$A_2$}{A2}}\label{l27:sec:a2core}

\paragraph{The statement.} $\dim M_{0,m}\le\ltwentysevendg^{\,m}$ for every $m$; the lower bound is free (the evaluation fusion product is
a quotient). By Lemma~\ref{l27:lem:truncweyl}, $M_{0,m}=W(m\theta,m)$, and the bound is Kus--Littelmann's Theorem~4.2 at $\lambda=\theta$,
$\lambda_0=0$, $N=m$, which for $\lambda_0=0$ carries no restriction on $\ltwentyseveng$; the specialization is identical for all six types.
The earlier sources are Feigin's Corollary~2 (\emph{SIGMA} 4 (2008) 070), whose proof has the gap of Section~\ref{l27:sec:feigin}, and
Ravinder's Corollary~2 at $a=0$, which is Feigin's statement imported as Ravinder's Proposition~5. \emph{Status:} published, with the
flagged step on p.~437 of Kus--Littelmann, described at the end of Section~\ref{l27:sec:feigin}; it applies equally to all six types,
and for $D_4$ at $m=2$ the step as written is shown by computation to be insufficient (the theorem is not known to fail).
In Lean: $D_4$ for $m\le2$ (Theorem~\ref{l27:thm:d4Tsmall}); for $D_4$ and all $m$, a reduction to $\mathfrak{sl}_4$ core bounds
with at least two adjoint sites, reduced further on paper (end of Section~\ref{l27:sec:feigin}); otherwise open.

\paragraph{What made $A_2$ tractable.} The $A_2$ proof (Chapter~\ref{ch:L26}, Section~\ref{l26:sec:cv}) uses three ingredients that are
special to $\mathfrak{sl}_3$: an enlarged family of modules with Chari--Venkatesh relations along the three positive roots, whose
members are expected to be fusion products of the $\mathfrak{sl}_3$ Kirillov--Reshetikhin modules $V(\omega_i)$, $V(2\omega_i)$, $V(3\omega_i)$
of dimensions $3,6,10$; explicit spanning arguments in the associated graded of these modules; and dimension splittings ($8=5+3$,
$5=3+2$) of the adjoint site. For the other types an analogous family would involve many more Kirillov--Reshetikhin modules, and
neither the family nor the base bounds are formalized.

\subsubsection{Why the published proof needed repair: 55 versus 63}\label{l27:sec:feigin}

Feigin realizes $D(1,m\theta)\cong D(1,\theta)^{*m}$ on the tensor product $T=\bigotimes_{j=1}^mD(1,\theta)(c_j)$ with its fusion filtration
$\Phi_\bullet$ (by degree in $t$). Let $f(t)=\prod_j(t-c_j)$. On $T$ the image of $\ltwentyseveng\otimes f(t)\C[t]$ is exactly the subspace $K$
of tensors with at least one factor in the trivial summand $\C$, and $T/K\cong\ltwentyseveng^{\otimes m}$. Input T is equivalent to
$\ltwentysevengr K\subseteq(t^m)\cdot\ltwentysevengr T$: the associated graded of $K$ is contained in the image of the ideal $\mathcal I_m$.
Feigin's Proposition~1 (\emph{SIGMA} 4 (2008) 070; Proposition~2.1 in arXiv v1) asserts that the associated graded of $T$ is a \emph{bigraded} fusion product in which a second degree $u$
counts the trivial factors.

Take $\mathfrak{sl}_3$ and $m=2$. Then $T=D(1,\theta)(c_1)\otimes D(1,\theta)(c_2)$ has dimension $81$, and by the number $u$ of trivial
factors it splits as $81=64+16+1$ ($u=0$: $\mathfrak{sl}_3\otimes\mathfrak{sl}_3$; $u=1$: $\mathfrak{sl}_3\otimes\C\oplus\C\otimes\mathfrak{sl}_3$;
$u=2$: $\C\otimes\C$). The fusion filtration has cumulative dimensions $\dim\Phi_0,\dots,\dim\Phi_4=27,55,72,80,81$ (the coefficients
of $\mathrm{Dem}_2=27+28q+17q^2+8q^3+q^4$ summed up; Example~\ref{l27:ex:sl3}). If $\Phi_1$ were the direct sum of its projections to the
three $u$-pieces, as the natural reading of Feigin's argument requires, the dimensions of those projections would add up to
$\dim\Phi_1=55$. The project's exact computation gives $63$. For $m=3$ the computation gives, for $d=1,2$, the pair $(\dim\Phi_d,\ \text{sum of the dimensions of the projections of }\Phi_d\text{ to the }u\text{-pieces})$ equal to $(161,188)$ and $(313,403)$; the two entries of a pair would agree if the filtration were homogeneous for $u$. So the fusion
filtration is not homogeneous for $u$, and the printed argument (labelled ``the idea of the proof'') does not prove Proposition~1; a rigorous
Proposition~1 is as strong as the statement wanted.

\paragraph{The repair and its flagged step.} Kus and Littelmann (Pacific J.\ Math.\ 278 (2015); the remark on p.~438 says their
Theorem~3.5 generalizes Feigin's case $\ell=1$, $\lambda_0=0$, $\lambda_1=\theta$ and fixes its gap) prove the bigraded statement by
handling all consequences of the defining relations in the associated graded. The upper bound needs the hard direction of their
Theorem~3.5: the annihilator of the generator must pass to the associated graded. On p.~437 the written proof passes from ``the
difference $M-M'$ lies in the left ideal generated by the relations'' to ``its symbol kills the fusion vector''. That step controls only
the initial forms of the \emph{generators} of the ideal, not of every element of it. The project's source audits (for $A_2$,
\texttt{Notes/deligne/campaign/A2\_SOURCE\_AUDIT\_2026\_10\_27.md}, and the audit of the literature inputs for all types) flag this step as
unjustified as written. This is a gap in the written argument, not a counterexample: no counterexample to the theorem is known. Since the
specialization $\lambda=\theta$, $\lambda_0=0$ is the same for every type, the caveat applies uniformly to $G_2$, $D_4$, $F_4$, $E_6$,
$E_7$, $E_8$. That initial forms can misbehave is a real phenomenon: in a filtered algebra the initial forms of all consequences of a set of
relations can differ from those of the relations themselves (in the Heisenberg algebra $[x,y]=z$ with filtration weights $0,0,1$ the
left ideal $(x,y)$ contains $z$, but the symbols of $x$ and $y$ do not generate $z$ in the associated graded).

\paragraph{The flagged step tested for $D_4$ (October 2026).} For $\mathfrak{so}_8$, $\lambda=\theta$, $\lambda_0=0$, $\ell=1$, the
project tested both the statement that the p.~437 step is meant to establish and the step itself, by computation (Python, not
Lean; \texttt{Notes/deligne/campaign/D4\_KL\_SYMBOL\_LEMMA\_TEST.md}, with the reproducible script
\texttt{Notes/deligne/campaign/generators/D4KLSymbolLemmaTest.py}). The statement, that the symbols of all elements of the
annihilator of the generator of $D(1,m\theta)$ kill the generator of the bigraded fusion product (so that the associated graded
of $D(1,m\theta)$ is that fusion product), holds at $m=2$ (exact arithmetic over $\mathbb Q$) and at $m=3$ (modulo two primes),
block by block. The step does not prove it: at $m=2$ the left ideal generated by the symbols of the \emph{defining} relations has
codimension $842=29^2+1$, one more than the $29^2$ required (computed modulo a prime). The extra class is spanned by the symbol of
\begin{equation*}
  \tilde a=f_{\alpha_2}\,f_{\theta-\alpha_2}\,(f_\theta\otimes t^3)+(f_{\alpha_2}\otimes t)\,f_{\theta-\alpha_2}\,(f_\theta\otimes t^2),
\end{equation*}
which annihilates the generator of $D(1,2\theta)$ (checked both in the fusion model and in the module presented by the defining
relations alone) while its symbol is not in the ideal generated by the symbols of the relations; that symbol does kill the fusion
vector, as the theorem predicts. So the written step is insufficient at $m=2$: a proof must control symbols of consequences of the
relations, not only of the relations. Nothing here contradicts Kus--Littelmann's Theorem~4.2; what is missing is a complete proof
of that one step. The extra relations were identified by the same computations: they are the symbols of the Garland coefficients of
the $\mathfrak{sl}_2(\theta)$-string.

\paragraph{An independent route for $D_4$.} The library pursues a route to Input~T for $D_4$ that does not use the p.~437 step.
An explicit family $J'_m$ of symbol relations (the symbols of the defining relations together with the Garland symbols above) is
proved to lie in both annihilators, that of the associated graded of the universal module (\lean{Jprime_le_annGr}) and that of
the fusion vector (\lean{StarB2.Jprime_le_annF}). This reduces the $D_4$ final theorem, with Input~R unchanged, to a rank bound
$\dim\operatorname{PTE}_m\le29^m$ for every $m\ge1$ on an explicitly presented toroidal module
(\lean{D4_classicallyFree_of_presentedRankBoundE}). The bound is proved for $m=1$ (\lean{presentedRankBoundE_one'}) and for
$m=2$ (\lean{presentedRankBoundE_two}, which discharges the two layer bounds of \lean{presentedRankBoundE_two_of'}); for $m\ge3$ it
is open on this route. Targeted builds are green and the printed axioms standard, and the files are compiled by the
aggregate build of 3 November 2026. These results do not change the hypothesis of \lean{D4_classicallyFree_of_inputT}.

\paragraph{The level route for $D_4$ (state on 2 November 2026).} A second route organizes $W(m\theta,m)$ by levels for a
$\mathfrak{gl}_4$ Levi subalgebra and bounds each level by a truncated $\mathfrak{sl}_4$ module. We record its state, separating Lean
from paper. Each Lean declaration named here was checked by a targeted build of its audit module (no
\lean{sorryAx}; \texttt{\#print axioms} within \lean{[propext, Classical.choice, Quot.sound]}) and is compiled by the aggregate
build of 3 November 2026. Only the first item changes a final theorem: it removes Input R from the $D_4$ theorem
(Theorem~\ref{l27:thm:d4inputT}). The rest is progress toward Input T, which stays a hypothesis.
\begin{itemize}
\item \emph{Lean, Input R.} \lean{D4Char.d4_demazureRealization} proves
  \lean{∀ n, DemazureRealizationG D4sector (ravFam D4sector) D4Wp.hilbSeriesZ n}, the hypothesis \lean{hD} of the final theorem,
  with no hypotheses of its own, from a graded bound on the Levi side (\lean{D4Char.d4_gradedLeviBound}) and the squeeze
  (Theorem~\ref{l27:thm:d4R}; \texttt{\#print axioms} standard, re-checked after the aggregate build of 3 November 2026).
\item \emph{Lean, the reduction of Input T.} \lean{D4LevelReduction.d4_classicallyFree_of_coreBounds} derives the conclusion of the
  final theorem from \lean{TruncatedCoreRankBound a b c} for all $a,c\ge1$, a dimension bound for finite-dimensional
  $\mathfrak{sl}_4$ current modules with three proved relation families, truncated at $a+b+c$; Input R is supplied by
  \lean{D4Char.d4_demazureRealization}. The levels with one adjoint site are proved (\lean{truncatedCoreRankBound_singleAdjoint}),
  so \lean{d4_classicallyFree_of_twoAdjointCoreBounds} needs only the core bounds with $a\ge2$.
\item \emph{Lean, toward $a=2$.} The levels with $a=2$ use the bound at one adjoint site,
  \begin{equation*}
    (\mathrm{III})_B:\qquad \dim W(\theta+B\omega_2,B)/\langle f_2(B-1)v\rangle\le300\cdot6^{B-2}.
  \end{equation*}
  The flat relation it needs holds for every truncation (\lean{SL4T.adjointSiteFlatTopRelation_holds},
  \lean{SL4T.zeroAdjointFlatRelation_all}), so \lean{SL4T.adjointSiteLawB_of_stratumCount} derives the bound at one adjoint site
  from a stratum count of a coefficient ring alone, and \lean{SL4T.oneAdjointStratumCount_of_low_strata} reduces that count to its
  low strata. The marked full-flag rings behind the count have their exact dimensions and stratum bounds in Lean; the step still
  missing is the descent from those rings to the coefficient ring.
\item \emph{On paper (refereed by a Claude session, not by a human referee, and not in Lean).} $(\mathrm{III})_B$ holds for every
  $B$ (a gap in the first version was found by the Lean formalization and repaired). Given $(\mathrm{III})$, every level with $a=2$
  follows from the single family
  \begin{equation*}
    (\mathrm{II})_\kappa:\qquad \dim W(2\theta+\kappa\omega_2,\kappa+2)/\langle f_1(1)f_3(1)v\rangle\le205\cdot6^\kappa\qquad(\kappa\ge1),
  \end{equation*}
  proved at $\kappa=0$ and checked by machine for $\kappa\le3$. A same-point splitting reduces $(\mathrm{II})_\kappa$ to two
  one-node statements: $\tilde Q_\kappa:=\dim W(2\omega_3+\kappa\omega_2,\kappa+1)/\langle f_3(1)v\rangle\le10\cdot6^\kappa$, and
  $(C)_{\kappa+1}$, where $(C)_n$ says that in the same-point product $V(\omega_2)_0\otimes W(n\omega_2,n)$ the tensor of the
  highest-weight vectors generates a module of dimension $20\cdot6^{n-1}$. Both are open in general; $(C)_n$ is proved for
  $n\le2$, and both hold by machine in the ranges checked ($\kappa\le3$, $n\le4$). The levels with $a\ge3$ are reduced, by a
  peeling argument, to local repairs and three base families, which are open.
\end{itemize}
So Input~T for $D_4$ remains open, in Lean and on paper; what is new is the precise form of what remains.

An alternative route to the $A_2$ bound, through a \emph{slice lemma} in mixed evaluation models, is recorded on paper in
Chapter~\ref{ch:L26}, Section~\ref{l26:sec:open}; it is not proved, and it is not used.

\subsection{Input R beyond \texorpdfstring{$A_2$}{A2}}\label{l27:sec:demcore}

\paragraph{What is needed.} For each $j$, a module with the relations of $M_{j,0}$ whose graded dimension is at least $\mathrm{Dem}_j$
\emph{in every degree}. It has two halves: a model satisfying the relations (the \emph{relation half}), and a graded lower bound for its
dimension (the \emph{graded half}).

\paragraph{The literature: a theorem for every type.} Both halves are published for every simple $\ltwentyseveng$, including $G_2$ and
$F_4$: the relation half is Ravinder's Proposition~1 (after Chari--Venkatesh, Corollary~3.5) together with Lemma~\ref{l27:lem:truncweyl},
and the graded half is the Demazure character formula of Kumar and Mathieu (dictionary of Definition~\ref{l27:def:realization}). Input R
is therefore \emph{not} an open mathematical problem. What is open is its formalization: a Lean model of $D(1,j\theta)$ with a proved
graded dimension, and a proof that the Lean algorithm implements the Demazure character formula.

\paragraph{What is formalized.}
\begin{itemize}
\item \emph{The relation half.} In the lattice VOA of the root lattice the vector $e^{j\theta}$ satisfies the $M_{j,0}$ relations, for $D_4$,
  $E_6$, $E_7$ and $E_8$ (Frenkel--Kac: the level-one vacuum module is the lattice VOA; Chapter~\ref{ch:L26},
  Section~\ref{l26:sec:lattice}). $G_2$ and $F_4$ are not simply laced and have no lattice model in the library; a formalization would need
  another construction of $D(1,j\theta)$ (the mathematics is not in doubt).
\item \emph{The graded half.} For $D_4$ the cases $j\le2$ are proved (Theorem~\ref{l27:thm:d4Rsmall}), and every $j$ is covered by
  \lean{D4Char.d4_demazureRealization}, through a graded Levi bound and the squeeze rather than the Demazure character formula
  (Theorem~\ref{l27:thm:d4R}; confirmed by the aggregate build of 3 November 2026). Numerically the lattice model is
  graded-sharp ($A_2$ and $D_4$ for $j\le3$). For the other types a Lean proof for all $j$ would need a formalization of the Demazure character formula in
  this setting, or an analogue for types $D$ and $E$ of the Chari--Loktev spanning basis, which for $\mathfrak{sl}_{r+1}$ gives the matching
  upper bound; the library has neither.
\item \emph{Total versus graded.} A lower bound on the \emph{total} dimension of $M_{j,0}$ (expected value $(\ltwentysevendg+1)^j$) would
  not be enough: the squeeze uses the inequality $\mathrm{Dem}_j\le\mathrm{ch}\,M_{j,0}$ coefficient by coefficient (Example~\ref{l27:ex:sl3}).
  For $A_2$ the passage from total to graded used the Chari--Loktev upper bound, special to type $A$.
\end{itemize}

\subsection{Input S for \texorpdfstring{$G_2$}{G2}, \texorpdfstring{$F_4$}{F4}, \texorpdfstring{$E_6$}{E6}, \texorpdfstring{$E_7$}{E7}, \texorpdfstring{$E_8$}{E8}}\label{l27:sec:starcore}

\paragraph{Why the bound is no cheaper than the identity.} Write $\mathcal D=\sum_jq^j\mathrm{Dem}_j/(q)_j$, a sum of series with nonnegative
coefficients. Suppose a proof of $\mathcal D\le c$ produced series $U_j\ge\mathrm{Dem}_j$ coefficientwise with $\sum_jq^jU_j/(q)_j=c$. If
equality $\mathcal D=c$ holds (as it does numerically), then $\sum_jq^j(U_j-\mathrm{Dem}_j)/(q)_j=0$ with every term nonnegative, so
$U_j=\mathrm{Dem}_j$. A domination proof therefore has to compute the Demazure characters exactly (Chapter~\ref{ch:L25},
Proposition~\ref{l25:prop:rigidity}).

\paragraph{What is known.} No published proof of the identity, or of the bound, is known to us at these five types: Input S is a
conjecture (Conjecture~\ref{l27:conj:star}), supported by numerical checks only. The structural route is Macdonald theory at $t=0$: Demazure characters are specializations of nonsymmetric
Macdonald polynomials (Sanderson; Ion), and the identity would be a constant-term summation of Cherednik type; we know of no reference
proving the needed summation at the exceptional types. For $G_2$ the library has reduced the problem to one identity: the target half
\lean{G2KWEqMLDE} is proved (\lean{G2Star.g2KWEqMLDE}), so the star identity follows from \lean{G2DemazureEqKW} alone
(\lean{G2Star.g2StarIdentity_of_demazure}; Section~\ref{l27:sec:starinput}). That identity is in turn split into two named halves
(Theorem~\ref{l27:thm:G2F4ladder}): the statement that the $G_2$ Demazure table equals a fermionic formula, and a pure $q$-series
identity between the fermionic sum and the Kac--Wakimoto-type product. The first, \lean{G2DemazureFermionic}, follows from published results
(Fourier--Littelmann, Chari--Moura, Ardonne--Kedem, Di Francesco--Kedem; conventions as in Lin, IMRN 2021; see
Section~\ref{l27:sec:starinput} for the strength of this reading); the second, \lean{G2FermionicKW}, is open.

%======================================================================
\section{The per-type headline theorems}\label{l27:sec:ladders}
%======================================================================

The final theorems are assembled from per-type headline theorems, each with the same conclusion and a longer list of hypotheses. We
present those the final theorems use, one for each mechanism, and list the others. Every theorem here has the conclusion of
Definition~\ref{l27:def:conclusion} for its type.

\subsection{\texorpdfstring{$A_2$}{A2}: classical freeness from U5 and U6}

The $A_2$ headline files, in logical order (each hypothesis that disappears from one row to the next became a theorem):
\begin{center}\small
\begin{longtable}{p{0.20\textwidth}p{0.30\textwidth}p{0.40\textwidth}}
\toprule
file & theorem & named hypotheses\\
\midrule\endhead
\texttt{A2Assembly} & \lean{classicallyFree_}\allowbreak\lean{of_residualInputs} & carrier $\le$ Demazure sum, \lean{StarIdentity},
  \lean{KacWakimotoLowerBound}\\
\texttt{A2SectorAssembly} & \lean{classicallyFree_}\allowbreak\lean{of_sectorInputs} & \lean{CasimirLayerBound}, \lean{SectorFusionBound F},
  U6, \lean{StarIdentity}, \lean{KacWakimotoLowerBound}\\
\texttt{A2SectorHeadline} & \lean{classicallyFree_}\allowbreak\lean{of_weylInputs} & \lean{CasimirLayerBound}, \lean{CarrierGeneration},
  U5, U6, \lean{StarIdentity}, \lean{KacWakimotoLowerBound}\\
\texttt{A2Headline} & \lean{classicallyFree_}\allowbreak\lean{of_openInputs} and four siblings & \lean{MFreeness},
  \lean{DivisorComparison}, U5, U6, \lean{KacWakimotoLowerBound}\\
\texttt{A2SectorTop} & \lean{classicallyFree_}\allowbreak\lean{of_divisorAndLiterature} & \lean{DivisorComparison}, U5, U6, KW\\
\texttt{A2Top} & \lean{classicallyFree_}\allowbreak\lean{of_pairAndLiterature} & \lean{PairComparison}, U5, U6, KW\\
\texttt{A2DivisorTop} & \lean{classicallyFree_}\allowbreak\lean{of_literature} & U5, U6, KW\\
\texttt{A2Final} & \lean{classicallyFree_}\allowbreak\lean{of_fusion} & U5, U6\\
\bottomrule
\end{longtable}
\end{center}
Here KW is \lean{KacWakimotoLowerBound} ($[q^N]\prod(1+q^r)^8\le\dim W_N$), proved from the $A_2$ product identity and the vacuum quartic
relation (Chapter~\ref{ch:L25}); \lean{StarIdentity} is \lean{A2.starIdentity}; \lean{CasimirLayerBound}, \lean{MFreeness},
\lean{CarrierGeneration}, \lean{PairComparison} and \lean{DivisorComparison} are steps B1, B6, B8 and B7 of the carrier chain, proved in
Chapter~\ref{ch:L24}. After \texttt{A2Final}, Chapter~\ref{ch:L26} replaces U5 and U6 by Inputs T and R.

\begin{theorem}[$A_2$ from U5 and U6]\label{l27:thm:A2ladder}
For $W=L_{-3/2}(\mathfrak{sl}_3)$ with its explicit presentation, classical freeness and $\ker=(\ltwentysevenIns)$ follow from U5 and U6 for
any family $F$.
\end{theorem}
\noindent File \leanpath{VOA/\allowbreak DeligneAudit/\allowbreak A2Final.lean}.
\begin{leancode}
theorem classicallyFree_of_fusion (F : ℕ → PowerSeries ℚ)
    (hU5 : SectorWeyl.FusionModuleBound F)
    (hU6 : MixedFusionRecursion F demazureHilbSeries) :
    CaseClassicallyFree .A2 W explicitPresentation ∧
      RingHom.ker (MvPolynomial.aeval fun i =>
        liDegreeZeroClass W 1 ⟨explicitPresentation.gen i,
          explicitPresentation.gen_weight_one i⟩).toRingHom = I2ns :=
\end{leancode}

\begin{dictionary}
Namespace \lean{VOA.DeligneAudit.A2}. \lean{SectorWeyl.FusionModuleBound F} is U5 in the $\mathfrak{sl}_3$ form
(\leanpath{VOA/DeligneAudit/A2SectorWeyl.lean}): for every module with the local Weyl and fusion relations for $m\theta$,
$\dim(\text{degree }b)\le[q^b]F_m$ and total $\le8^m$. \lean{MixedFusionRecursion F Dem} is the structure with the single field
\lean{eq : ∀ n, Dem n = VOA.DeligneGaussian.mixed (X : PowerSeries ℚ) F n 0} (U6; \leanpath{VOA/DeligneAudit/A2CarrierChain.lean}), where
\lean{mixed q F a b} is defined by $M_{0,b}=F_b$ and $M_{a+1,b}=M_{a,b+1}+q^{2a+b+1}M_{a,b}$
(\leanpath{VOA/DeligneAudit/GaussianTransform.lean}); so the condition is $M_{j,0}=\mathrm{Dem}_j$. \lean{demazureHilbSeries} is the power
series of the Gelfand--Tsetlin table \lean{demazureHilb} (\leanpath{VOA/DeligneAudit/A2SectorAssembly.lean}). $F$ is a free variable: U5
and U6 together pin it down to the fusion characters. The file has no \lean{variable} lines.
\end{dictionary}

\textbf{Proof on paper.} The carrier chain of Chapter~\ref{ch:L24} gives $\mathrm{carrierDim}(N)\le\mathcal D_N$ from U5 and U6 (B1
proved, B6 from U5, the divisor comparison proved, D3 proved). The star identity gives $\mathcal D_N=[q^N]\prod(1+q^r)^8$, and KW gives
$[q^N]\prod(1+q^r)^8\le\dim W_N$; the capstone (Theorem~\ref{l24:thm:upper}) concludes. \qed

\begin{physics}
\emph{Heuristic.} With the fusion of adjoint particles described by its graded character, the count of classical states (the carrier) is
bounded by the particle count (the Demazure sum), which equals the vacuum character; so no classical state fails to lift.
\end{physics}

\begin{status}
Proved; axioms standard (campaign record). Conditional on U5 and U6 for $F$; Chapter~\ref{ch:L26} discharges both for $F=$ \lean{fusionF}.
\end{status}

\begin{lemma}[Two bookkeeping identities of the $A_2$ assembly]\label{l27:lem:A2book}
$\mathrm{(i)}$ For every $j$, $1/(q)_j=\sum_Np_{\le j}(N)q^N$ in $\Q[[q]]$, where $p_{\le j}(N)$ is the number of partitions of $N$ into parts at most
$j$. $\mathrm{(ii)}$ The concrete Demazure sum of Chapter~\ref{ch:L25} equals the coefficient $\sum_j[q^N]\,q^j\mathrm{Dem}_j/(q)_j$ used by the
carrier chain, for $\mathrm{Dem}=$ \lean{demazureHilbSeries}.
\end{lemma}
\noindent File \leanpath{VOA/\allowbreak DeligneAudit/\allowbreak A2SectorAssembly.lean}.
\begin{leancode}
theorem partitionsLe_series (n : ℕ) : (qpoch n)⁻¹ = partitionSeries n :=
\end{leancode}
\begin{leancode}
theorem demazureCoeff_eq_demazureCoeffOf (N : ℕ) :
    (demazureCoeff N : ℚ) = demazureCoeffOf demazureHilbSeries N := by
\end{leancode}

\begin{dictionary}
\lean{qpoch n} is $(q)_n$ in $\Q[[q]]$, invertible because its constant term is $1$; \lean{partitionSeries n} is
\lean{PowerSeries.mk fun N => (partitionsLe n N : ℚ)}. \lean{demazureCoeffOf Dem N} is the chain's coefficient
$\sum_{j\le N}[q^N]\,q^j\mathrm{Dem}_j/(q)_j$ (Chapter~\ref{ch:L24}); \lean{demazureCoeff N} is the natural number
$\sum_j(\mathrm{hilb}_j\ast p_{\le j})(N-j)$ of Chapter~\ref{ch:L25}. The two lemmas say that two independently written definitions of
$\mathcal D_N$ agree.
\end{dictionary}

\textbf{Proof on paper.} (i) Split partitions into parts $\le j+1$ by whether a part $j+1$ occurs:
$p_{\le j+1}(N)=p_{\le j}(N)+p_{\le j+1}(N-j-1)$. Hence $P_{j+1}(1-q^{j+1})=P_j$ for the series $P_j$, and induction from $P_0=1$ gives
$P_j(q)_j=1$. (ii) Expand $q^j\mathrm{Dem}_j/(q)_j$ with (i) and compare coefficients. \qed

\begin{physics}
$1/(q)_j$ is the partition function of the relative motion of $j$ identical bosonic particles on a line. The lemma checks that two
programs counting the same particle states agree.
\end{physics}

\begin{status}
Proved, unconditional; axioms standard (campaign record).
\end{status}

\subsection{\texorpdfstring{$D_4$}{D4}: the headline with every input visible}

\begin{theorem}[$D_4$ headlines]\label{l27:thm:D4ladder}
For $W=L_{-2}(\mathfrak{so}_8)$ (certificate basis) with its explicit presentation:
\begin{enumerate}
\item[(a)] $c(W)=-14=-2(6+1)$;
\item[(b)] classical freeness and $\ker=$ the ideal of the $105$ quadrics follow from U5 and U6 for a family $F$, the divisor comparison,
  Demazure positivity, the Schur-index evaluation, Ramanujan's identities and a vacuum quartic relation on $W$;
\item[(c)] they follow from U5, U6 and Demazure positivity alone;
\item[(d)] they follow from the squeeze inputs for a family $M$ together with the universal bound.
\end{enumerate}
\end{theorem}
\noindent File \leanpath{VOA/\allowbreak DeligneAudit/\allowbreak D4Top.lean}.
\begin{leancode}
theorem W4_centralCharge : W4.centralCharge = -2 * (((6 : ℕ) : ℂ) + 1) := by
\end{leancode}
\begin{leancode}
theorem d4_headline (F : ℕ → PowerSeries ℚ) (hU5 : FusionModuleBoundG D4sector F)
    (hD : DivisorComparisonG D4sector) (hU6 : MixedFusionRecursion F (hilbSeries d4DemazureHilb))
    (hpos : D4K.D4DemazurePositivity) (hS : D4MLDE.D4SchurEqE4) (hR : D4MLDE.RamanujanIdentities)
    {κ : ℂ} (hrel : VacuumQuarticRelation W4 κ) :
    CaseClassicallyFree .D4 W4 explicitPresentation4 ∧
      RingHom.ker (MvPolynomial.aeval fun i =>
        liDegreeZeroClass W4 1 ⟨explicitPresentation4.gen i,
          explicitPresentation4.gen_weight_one i⟩).toRingHom = explicitPresentation4.ideal :=
\end{leancode}
\noindent File \leanpath{VOA/\allowbreak DeligneAudit/\allowbreak D4Final.lean}.
\begin{leancode}
theorem classicallyFree_of_fusion (F : ℕ → PowerSeries ℚ) (hU5 : FusionModuleBoundG D4sector F)
    (hU6 : MixedFusionRecursion F (hilbSeries d4DemazureHilb))
    (hpos : D4K.D4DemazurePositivity) :
    CaseClassicallyFree .D4 W4 explicitPresentation4 ∧
      RingHom.ker (MvPolynomial.aeval fun i =>
        liDegreeZeroClass W4 1 ⟨explicitPresentation4.gen i,
          explicitPresentation4.gen_weight_one i⟩).toRingHom = explicitPresentation4.ideal :=
\end{leancode}
\noindent File \leanpath{VOA/\allowbreak DeligneAudit/\allowbreak D4Squeeze.lean}.
\begin{leancode}
theorem classicallyFree_of_squeeze {M : ℕ → ℕ → PowerSeries ℚ}
    (h : FusionSqueezeInputsG D4Wp.hilbSeriesZ 28 M)
    (huniv : ∀ (m : ℕ) (V : Type) [AddCommGroup V] [Module ℂ V]
      (ρ : CurrentModule VOA.DeligneAudit.DeligneSector.D4.D4sector.f V) (w : V),
      LocalWeylRelG VOA.DeligneAudit.DeligneSector.D4.D4sector ρ w m →
        ∀ b, (Module.finrank ℂ (cycDeg ρ w b) : ℚ) ≤ coeff b (M 0 m)) :
    CaseClassicallyFree .D4 W4 explicitPresentation4 ∧
      RingHom.ker (MvPolynomial.aeval fun i =>
        liDegreeZeroClass W4 1 ⟨explicitPresentation4.gen i,
          explicitPresentation4.gen_weight_one i⟩).toRingHom = explicitPresentation4.ideal :=
\end{leancode}

\begin{dictionary}
All in the namespace \lean{VOA.DeligneAudit.D4Lie}. \lean{W4.centralCharge} is the Sugawara central charge
$k\ltwentysevendg/(k+h^\vee)=-2\cdot28/4=-14$, written as $-2(h^\vee+1)$ because that is the form the MLDE argument uses. In
\lean{d4_headline}: \lean{FusionModuleBoundG D4sector F} is U5 (Definition~\ref{l24:def:U5}); \lean{DivisorComparisonG D4sector} is the
divisor comparison (Chapter~\ref{ch:L24}); \lean{D4K.D4DemazurePositivity} says the algorithm's integer $D_4$ Demazure coefficients are
nonnegative; \lean{D4MLDE.D4SchurEqE4} says the series \lean{d4SchurQ} equals the coefficients of $q^{-7/12}E_4'/(240\eta^{10})$, where
$E_4'=q\,\frac{d}{dq}E_4$ and the factor $q^{-7/12}$ removes the leading power $q^{-c/24}=q^{7/12}$ of the normalised character;
\lean{D4MLDE.RamanujanIdentities} are Ramanujan's three differential identities for $E_2,E_4,E_6$; and \lean{VacuumQuarticRelation W4 κ} is
the weight-four identity of Chapter~\ref{ch:L25} (Definition~\ref{l25:def:vqr}) for some constant $\kappa$ (implicit argument
\lean{\{κ : ℂ\}}). In \lean{classicallyFree_of_squeeze}: \lean{FusionSqueezeInputsG Dem 28 M} is the generic squeeze-input structure of
Chapter~\ref{ch:L26} (Definition~\ref{l26:def:squeezeinputs}), and \lean{huniv} says every module with the $M_{0,m}$ relations has degree
pieces bounded by $M_{0,m}$.

\textbf{Which hypotheses became theorems.} Between (b) and (c): the divisor comparison (\lean{DeligneDivisor.D4.d4_divisorComparison}),
the Ramanujan identities (\lean{D4MLDE.ramanujanIdentities}), the vacuum quartic relation (\lean{vacuumQuarticRelation4}) and the Schur
evaluation (\lean{D4Wp.d4SchurEqE4}). Between (c) and the final theorem: U5 and U6 from Inputs T and R, and positivity, absorbed by the
unclamped series (Chapter~\ref{ch:L26}). Last step: Input R itself (\lean{D4Char.d4_demazureRealization}), which leaves Input T as the
only hypothesis (Theorem~\ref{l27:thm:d4inputT}).
\end{dictionary}

\textbf{Proof on paper.} (a) Substitute $k=-2$, $\ltwentysevendg=28$, $h^\vee=6$. (b) The vacuum quartic relation with $c=-14$ gives
$c_N\le\dim W_N$ for $h^\vee=6$; with Ramanujan's identities and the Schur evaluation the Demazure sum equals the MLDE solution
(Theorem~\ref{l25:thm:d4}), using positivity to read the algorithm's table as nonnegative; so $\mathcal D_N\le\dim W_N$. U5, U6 and the
divisor comparison give the carrier bound (Chapter~\ref{ch:L24}), and the capstone concludes. (c) Plug in the four proved inputs.
(d) The squeeze gives U5 for the unclamped fusion series (Chapter~\ref{ch:L26}), and the unclamped form of (c) applies. \qed

\begin{physics}
(a) is the relation $c_{2d}=-12c_{4d}$ of Beem et al.\ with $c_{4d}=7/6$ for $SU(2)$ with four flavours. The Kac--Wakimoto character formula
used for $A_2$ is unavailable here, because $-2$ is not an \emph{admissible} level for $D_4$ (admissible levels are those whose vacuum
module has a Kac--Wakimoto character formula; for simply laced $\ltwentyseveng$ they are $k+h^\vee=p/u$ with coprime positive integers
$p\ge h^\vee$, $u$, and for $D_4$ at $k=-2$ one has $k+h^\vee=4<6$); the modular differential equation replaces it.
\end{physics}

\begin{status}
All proved; axioms standard (campaign record). (a) unconditional; (b)--(d) conditional on the listed hypotheses. The files
\texttt{D4LieTop} (\lean{d4_classicallyFree_of_sector}) and \texttt{D4Top} (\lean{W4_demazureBound}, \lean{W4_vacuum_ne_zero}) contain the
intermediate steps.
\end{status}

\subsection{\texorpdfstring{$G_2$}{G2} and \texorpdfstring{$F_4$}{F4}: the concrete headlines and the two halves of the \texorpdfstring{$G_2$}{G2} star identity}

\begin{theorem}[$G_2$ and $F_4$ concrete headlines]\label{l27:thm:G2F4ladder}
For $W=L_{-5/3}(\mathfrak g_2)$, classical freeness and $\ker=$ the ideal of the listed quadrics follow from U5 for the clamped fusion series
and Input S; or from U5 and the two halves of the star identity. For $W=L_{-5/2}(\mathfrak f_4)$ they follow from U5 and Input S.
\end{theorem}
\noindent File \leanpath{VOA/\allowbreak DeligneAudit/\allowbreak G2LieHeadline.lean}.
\begin{leancode}
theorem g2_classicallyFree_concrete (hU5 : FusionModuleBoundG G2sector g2FusionF)
    (hB : G2StarBound) :
    CaseClassicallyFree .G2 g2W explicitPresentation ∧
      RingHom.ker (MvPolynomial.aeval fun i =>
        liDegreeZeroClass g2W 1 ⟨explicitPresentation.gen i,
          explicitPresentation.gen_weight_one i⟩).toRingHom = explicitPresentation.ideal :=
\end{leancode}
\begin{leancode}
theorem g2_classicallyFree_concrete_of_fermionic (hU5 : FusionModuleBoundG G2sector g2FusionF)
    (hF1 : G2Star.G2DemazureFermionic) (hF2 : G2Star.G2FermionicKW) :
    CaseClassicallyFree .G2 g2W explicitPresentation ∧
      RingHom.ker (MvPolynomial.aeval fun i =>
        liDegreeZeroClass g2W 1 ⟨explicitPresentation.gen i,
          explicitPresentation.gen_weight_one i⟩).toRingHom = explicitPresentation.ideal :=
\end{leancode}
\noindent File \leanpath{VOA/\allowbreak DeligneAudit/\allowbreak F4LieHeadline.lean}.
\begin{leancode}
theorem f4_classicallyFree_concrete' (hU5 : FusionModuleBoundG F4sector f4FusionF)
    (hB : F4StarBound) :
    CaseClassicallyFree .F4 f4W explicitPresentation ∧
      RingHom.ker (MvPolynomial.aeval fun i =>
        liDegreeZeroClass f4W 1 ⟨explicitPresentation.gen i,
          explicitPresentation.gen_weight_one i⟩).toRingHom = explicitPresentation.ideal :=
\end{leancode}
\noindent File \leanpath{VOA/\allowbreak DeligneAudit/\allowbreak G2StarFermionic.lean}.
\begin{leancode}
structure G2DemazureFermionic : Prop where
  eq : ∀ n N, g2DemazureHilb n N = g2FermHilb n N

/-- **Named (one `q`-series identity)**: the fermionic Demazure sum is the Kac–Wakimoto form. -/
structure G2FermionicKW : Prop where
  eq : ∀ N, ((demazureSumCoeff g2FermHilb N : ℕ) : ℤ) = coeff N g2KW
\end{leancode}

\begin{dictionary}
\lean{g2FusionF} and \lean{f4FusionF} are \lean{fusionOf (hilbSeries g2DemazureHilb)}, resp.\ \lean{fusionOf (hilbSeries f4DemazureHilb)}:
the fusion characters determined by the clamped Demazure tables through the mixed recursion (Chapter~\ref{ch:L26},
Definition~\ref{l26:def:fusionOf}). U6 holds for them by construction, so U5 is the only fusion hypothesis.

\lean{G2DemazureFermionic} says that the $G_2$ Demazure table equals \lean{g2FermHilb}, the fermionic formula of Hatayama, Kuniba, Okado,
Takagi and Yamada, $\sum_\lambda\dim V(\lambda)\,X_{j,\lambda}(q)$, for $j$ copies of the Kirillov--Reshetikhin module $W^{(2)}_1=V(\theta)\oplus\C$.
On paper it follows from published results: $D(1,j\theta)\cong(W^{(2)}_1)^{*j}$ as graded modules (Fourier--Littelmann, Theorems~8
and~B; Chari--Moura for $W^{(2)}_1$), and the graded multiplicities of this fusion product are the fermionic sums (Ardonne--Kedem's upper
bound with Di Francesco--Kedem's equality at $q=1$; conventions as in Lin, IMRN 2021). Section~\ref{l27:sec:starinput} records the
strength of this reading. In Lean it is a single named statement, a hypothesis of \lean{g2_classicallyFree_concrete_of_fermionic}. \lean{G2FermionicKW} says that the Demazure sum of
the fermionic table equals the coefficients of \lean{g2KW}, the product
$\theta_{A_2}(q)\,(q^3;q^3)_\infty^6/(q;q)_\infty^8$, where $(x;q)_\infty=\prod_{i\ge0}(1-xq^i)$ (so $(q;q)_\infty=\prod_{i\ge1}(1-q^i)$) and
$\theta_{A_2}$ is the theta series of the $A_2$ root lattice; this product is the Kac--Wakimoto-type closed form of the target of
Chapter~\ref{ch:L25} (Theorem~\ref{l25:thm:g2target}). It is a pure $q$-series identity, and it is open (checked through $q^{60}$).
\end{dictionary}

\textbf{Proof on paper.} Section~\ref{l27:sec:commonproof}, steps 4--6. For the fermionic form: the two halves give
$\mathcal D_N=[q^N]$\lean{g2KW} (\lean{g2DemazureEqKW_of_fermionic}), and \lean{g2KW} equals the MLDE solution by a modular-forms argument
(Lorenz--Borwein, Chapter~\ref{ch:L25}, Theorem~\ref{l25:thm:g2target}), which gives the star bound. \qed

\begin{physics}
\emph{Heuristic.} The fermionic formula is the Bethe-ansatz count of states of a spin chain built from Kirillov--Reshetikhin modules; its
equality with a product formula is the ``bosonic $=$ fermionic'' kind of identity familiar from conformal field theory characters.
\end{physics}

\begin{status}
Proved; axioms standard (\lean{g2_classicallyFree_concrete} and \lean{f4_classicallyFree_concrete'} re-checked at an earlier commit; the
others campaign record). Conditional: U5 and the star input. The squeeze forms (\lean{g2_classicallyFree_concrete_of_squeeze},
\lean{…_Z}, \lean{…_Z_of_fermionic}, \lean{f4_classicallyFree_concrete_of_squeeze} and its primed form) replace U5 by the squeeze inputs
and the universal bound.
\end{status}

\subsection{\texorpdfstring{$E$}{E}-types}

\begin{theorem}[$E_6$, $E_7$, $E_8$ headlines]\label{l27:thm:Eladder}
For $X\in\{E_6,E_7,E_8\}$ and $W=L_k(\ltwentyseveng_X)$ with its explicit presentation, classical freeness and $\ker=$ the ideal of the listed
quadrics follow from U5 for the clamped fusion series and Input S.
\end{theorem}
\noindent File \leanpath{VOA/\allowbreak DeligneAudit/\allowbreak E6Final.lean}.
\begin{leancode}
theorem e6_classicallyFree_of_starBound
    (hU5 : FusionModuleBoundG E6sector (fusionOf (hilbSeries e6DemazureHilb)))
    (hB : E6StarBound) :
    CaseClassicallyFree .E6 W6 explicitPresentation ∧
      RingHom.ker (MvPolynomial.aeval fun i =>
        liDegreeZeroClass W6 1 ⟨explicitPresentation.gen i,
          explicitPresentation.gen_weight_one i⟩).toRingHom = explicitPresentation.ideal :=
\end{leancode}
\noindent File \leanpath{VOA/\allowbreak DeligneAudit/\allowbreak E7Final.lean}.
\begin{leancode}
theorem e7_classicallyFree_of_starBound
    (hU5 : FusionModuleBoundG E7sector (fusionOf (hilbSeries e7DemazureHilb)))
    (hB : E7StarBound) :
    CaseClassicallyFree .E7 W7 explicitPresentation7 ∧
      RingHom.ker (MvPolynomial.aeval fun i =>
        liDegreeZeroClass W7 1 ⟨explicitPresentation7.gen i,
          explicitPresentation7.gen_weight_one i⟩).toRingHom = explicitPresentation7.ideal :=
\end{leancode}
\noindent File \leanpath{VOA/\allowbreak DeligneAudit/\allowbreak E8Final.lean}.
\begin{leancode}
theorem e8_classicallyFree_of_starBound
    (hU5 : FusionModuleBoundG E8sector (fusionOf (hilbSeries e8DemazureHilb)))
    (hB : E8StarBound) :
    CaseClassicallyFree .E8 W8 explicitPresentation8 ∧
      RingHom.ker (MvPolynomial.aeval fun i =>
        liDegreeZeroClass W8 1 ⟨explicitPresentation8.gen i,
          explicitPresentation8.gen_weight_one i⟩).toRingHom = explicitPresentation8.ideal :=
\end{leancode}

\begin{dictionary}
Namespaces \lean{DeligneCarrier.E6}, \lean{E7Lie}, \lean{E8Lie}. U5 is stated directly for \lean{fusionOf (hilbSeries …)}. The same files
contain siblings \lean{eX_classicallyFree} with the star \emph{identity} \lean{EXStarIdentity} in place of the bound; the bound is weaker
and suffices. Inside, the divisor comparisons \lean{eX_divisorComparison}, jet surjectivity and the geometric carrier are plugged in as
theorems, and the target bounds \lean{e6_targetLowerBound}, \lean{e7_targetLowerBound}, \lean{targetLowerBoundQ8} are unconditional.
\end{dictionary}

\textbf{Proof on paper.} Steps 4--6 of Section~\ref{l27:sec:commonproof}, through the generic
\lean{DeligneCarrier.classicallyFree_of_sectorFusion} (Chapter~\ref{ch:L24}). \qed

\begin{physics}
\emph{Heuristic.} The Minahan--Nemeschansky statements of Theorem~\ref{l27:thm:E}, one step earlier: the fusion of adjoint particles is
still assumed in the form of its graded character.
\end{physics}

\begin{status}
Proved; conditional on U5 and \lean{EXStarBound}. Axioms standard: $E_6$ re-checked at an earlier commit; $E_7$, $E_8$ campaign record
only.
\end{status}

%======================================================================
\section{Approaches known to fail}\label{l27:sec:deadends}
%======================================================================

For the record, the following natural approaches to bounds of the kind in Inputs T and S are known to fail. Each comes with the reason.
They were found while working on $A_2$, and they apply equally to attempts at the other types.

\begin{warning}[title={Approaches to the total bound (Input T) that fail}]
\emph{(i) Circular.}
\begin{enumerate}
\item \emph{The presentation route} (``$M_{0,m}$ is the associated graded of the fusion model''). The surjection from $M_{0,m}$ onto the
  associated graded is free; its injectivity is exactly the statement to be proved.
\end{enumerate}
\emph{(ii) Numerically false.}
\begin{enumerate}\setcounter{enumi}{1}
\item \emph{The naive truncated box.} Imposing $f(k)w=0$ for $k\ge m$ together with ``box'' bounds on monomials overcounts: $72$ instead of
  $64$ for $\mathfrak{sl}_3$ at $m=2$, $656$ instead of $512$ at $m=3$. The truncations do not capture all relations.
\item \emph{Short truncation identities.} Reducing the truncation in $M_{0,m}$ to a single-letter recursion reproduces the correct count
  for $m\le3$ but fails in $28$ of $256$ cases at $m=4$: the truncation is a normal-form problem for the whole relation ideal.
\item \emph{Leading-term independence of Chari--Loktev words} against monomials of the lattice (Fock) model: the words are not independent
  under any tested monomial order ($77$ of $81$ at $N=2$, $547$ of $729$ at $N=3$).
\item \emph{Separating functionals in the lattice VOA $V_{A_2}$} by charge, $L_0$-weight or $\theta$-degree: these separate the pieces of
  Ravinder's chain for $N=2$ but not for $N=3$, and the $\theta$-degree is not preserved by $U(\ltwentyseveng[t])$.
\end{enumerate}
\emph{(iii) False on an algebraic technicality.}
\begin{enumerate}\setcounter{enumi}{5}
\item \emph{Plain tensor induction} $T_m=T_{m-1}\otimes D(1,\theta)$: the image of the fusion filtration of $T_m$ in $T_{m-1}\otimes\ltwentyseveng$ is
  a \emph{mixed} filtration, not a shifted tensor filtration (fusion products are not associative on the nose); one must work with mixed
  models.
\item \emph{Ungraded counting in the evaluation model.} At nonzero points $\mathcal I_m=\ltwentyseveng\otimes t^m\C[t]$ acts as all of
  $\ltwentyseveng[t]$, so only the graded structure carries information.
\item \emph{Treating a splitting of the modes as an algebra map before taking the associated graded.} Commutators do not transfer: with
  $[h_\theta,f_\theta]=-2f_\theta$ the relation $h_\theta(1)f_\theta(m-1)-f_\theta(m-1)h_\theta(1)+2f_\theta(m)=0$ has terms of filtration degree
  $0,0,1$; in the split algebra the commutator vanishes while the last term acts nontrivially. Checking only the initial forms of the
  defining generators does not prove that annihilators transfer.
\item \emph{Treating the image of the site-removing shift map of the polynomial model} (Feigin, Jimbo, Kedem, Loktev and Miwa) as an ideal.
  That image is the linear subspace of polynomials independent of one zero-mode variable; it contains $1$, so an ideal with that image
  would be everything. One must keep the actual target ideal modulo this subspace.
\end{enumerate}
\end{warning}

\begin{warning}[title={Approaches to the star bounds (Input S) that fail}]
\begin{enumerate}
\item \emph{Positivity or domination proofs.} Excluded by rigidity (Section~\ref{l27:sec:starcore}): a dominating family must equal the
  Demazure characters.
\item \emph{Positive expansions of the MLDE solution.} The exponents of the product expansions of the MLDE solutions alternate in sign, so
  there is no positive product to dominate.
\item \emph{A rational certificate for the hardest step of the $G_2$ Macdonald route}: the needed identities have binomial denominators of
  size $17$ to $41$, and an exhaustive search for a small local cancellation found none.
\end{enumerate}
\end{warning}

%======================================================================
\section{Verification record}\label{l27:sec:verification}
%======================================================================

\textbf{The commit.} Every \texttt{leancode} block of this chapter occurs verbatim in the sources of the state this chapter was written against
(\texttt{sync.py check}).

\textbf{What the machine confirmed (campaign records, 27 October 2026).}
\begin{plaincode}
VOA.Deligne umbrella      earlier: green, 11099 jobs, no sorryAx;
                          every #print axioms in the campaign files standard
lake build VOA            last green on 26 Oct (13868 jobs)
A2Unconditional           targeted build green, 9755 jobs, max RSS 6.13 GB;
                          CVUniv.fusionTotalDim, A2.classicallyFree:
                          [propext, Classical.choice, Quot.sound]
DeligneSeriesA2           targeted build green, 9756 jobs, standard axioms
CVUniv.sharpRectangularBases  targeted build green, 8863 jobs, standard axioms
CVUniv.residualTheta3Bound    targeted build green, 8887 jobs, standard axioms
\end{plaincode}
\textbf{What the machine confirmed (28 October 2026).}
\begin{plaincode}
bash scripts/campaign_build.sh VOA    (full lake build VOA, memory-capped wrapper)
  pre-release: green, 18901 jobs
  after merging the then-current origin/main: green, 18907 jobs
python3 scripts/check_repository.py         PASS: 10431 library files, all reachable,
                                            0 project axioms
python3 scripts/audit/run_axiom_audits.py   PASS: 5958 reports across 24 probes
#print axioms (lake env lean -DmaxSynthPendingDepth=3):
  VOA.DeligneAudit.A2.classicallyFree
  VOA.DeligneAudit.DeligneSeries.A2_classicallyFree_final
    both: [propext, Classical.choice, Quot.sound]
\end{plaincode}
Both declarations have no hypotheses. The aggregate target contains the umbrella \lean{VOA.Deligne}, which imports
\leanpath{VOA/DeligneAudit/A2Unconditional.lean}. The final theorems of the six conditional types are in
\leanpath{VOA/DeligneAudit/DeligneSeries.lean}, which is built within that aggregate build (earlier, within the umbrella build at
an earlier commit); the statements quoted here are unchanged since. A source scan of the import closure of \lean{A2Unconditional} (release audit,
\texttt{Notes/Textbook/referee/RELEASE\_AUDIT\_L24\_L25\_A2.md}) found no \lean{sorry}, \lean{admit}, \lean{axiom}, \lean{native_decide} or
\lean{implemented_by} outside comments.

\textbf{What the machine confirmed (3 November 2026): $D_4$ Input R.}
\begin{plaincode}
lake build VOA (aggregate)                  green, 20331 jobs
python3 scripts/check_repository.py         11859 library files, all reachable
                                            except 4 held from the release;
                                            0 axiom declarations
python3 scripts/audit/run_axiom_audits.py   PASS: 5958 reports across 24 probes
#print axioms:
  VOA.DeligneAudit.D4Char.d4_demazureRealization
  VOA.DeligneAudit.DeligneSeries.D4_classicallyFree_of_inputT
    both: [propext, Classical.choice, Quot.sound]
\end{plaincode}
The first declaration has no hypotheses (Theorem~\ref{l27:thm:d4R}); the second has the single hypothesis \lean{hF}, Input T for
$D_4$ (Theorem~\ref{l27:thm:d4inputT}). The four files held from the release are not imported by the $D_4$ theorems.

\textbf{What this chapter's author re-checked.} At two earlier commits (not re-run for this edition), on scratch files importing a
single module each, with \texttt{lake env lean -DmaxSynthPendingDepth=3} from the repository root (about 4\,s each, peak memory 5.8--5.95\,GB):
\begin{plaincode}
import VOA.DeligneAudit.DeligneFusionU7
  A2.classicallyFree_of_totalDim', A2.classicallyFree_of_kernelLower',
  A2.fusionTotalDim_of_kernelLower, EvalModel.demazureRealization_all,
  (verbatim copies of) A2_classicallyFree_final_S2
import VOA.DeligneAudit.G2FusionKernelTheta
  UnivG.g2_classicallyFree_of_totalBound, G2Lie.g2_classicallyFree_concrete,
  UnivG.g2_ravinderKernelRel
import VOA.DeligneAudit.DeligneFusionFiniteGG2 / GD4 / GF4
  FiniteG.G2.g2_fusionTotalBoundG_of_dim, FiniteG.fusionTotalBoundG_of_dim,
  FiniteG.D4.d4_fusionTotalBoundG_of_dim, FiniteG.F4.f4_fusionTotalBoundG_of_dim
import VOA.DeligneAudit.D4FusionKernelTheta
  UnivG.d4_classicallyFree_of_totalBound, UnivG.d4_ravinderKernelRel'
import VOA.DeligneAudit.F4FusionCountZ
  DeligneCarrier.f4_classicallyFree_of_totalBound', F4Lie.f4_classicallyFree_concrete'
import VOA.DeligneAudit.E6Final
  DeligneCarrier.E6.e6_classicallyFree_of_starBound
\end{plaincode}
Every one of these (names prefixed by \lean{VOA.DeligneAudit.}) printed \lean{[propext, Classical.choice, Quot.sound]}. The statements
concerned are unchanged in the state this chapter was written against. The $D_4$, $G_2$ and $F_4$ final theorems are compositions of two re-checked theorems each,
so their axioms are the union (standard). The author did not re-run any check at that state: this checkout has no compiled
objects for \lean{A2Unconditional}, \lean{DeligneSeriesA2} or the Chari--Venkatesh modules, and building them is outside the scope of a
chapter revision. The $E_7$/$E_8$ modules were not imported (they load the large certificate libraries); for the three $E$-type final
theorems the axiom status is the campaign record. The author did not run \texttt{lake build}.

\textbf{The literature-precision revision (28 October 2026).} The statements of the literature inputs were corrected following
\texttt{Notes/Textbook/referee/DELIGNE\_INPUTS\_PRECISION.md}. The Lean names cited in that revision (Section~\ref{l27:sec:partial},
the $G_2$ reduction of Input S, the lattice relation halves) were checked by reading the source at the time of that revision, and the six new
quotations occur verbatim there; their files are in the import closure of \leanpath{VOA/Deligne.lean}. Their axiom status is the
campaign record; no Lean check was re-run for this revision, and no \texttt{lake build} was run. The literature
claims are those of the audit, which read the papers' own text (Feigin arXiv v1/v2, Ravinder arXiv v4, Kus--Littelmann published
PDF); the Kumar, Mathieu, Frenkel--Kac and Chari--Venkatesh citations were not re-fetched.

%======================================================================
\section{Remarks on the Lean sources}\label{l27:sec:defects}
%======================================================================

The statements of the final theorems say what the mathematics says, about genuine simple affine VOAs, with every hypothesis explicit.
Some docstrings in the sources describe an earlier state and contradict the theorems; a reader of the source should rely on the
statements.
\begin{warning}[title={Docstrings that are out of date}]
\begin{enumerate}
\item \leanpath{VOA/DeligneAudit/DeligneFusionCVTotal.lean}: \lean{SharpRectangularBases} and \lean{ResidualTheta3Bound} are labelled
  ``UNPROVED \dots\ explicit mathematical input'', and the module docstring says the file ``does not assert unconditional $A_2$ classical
  freeness''. Both base families are now proved (\lean{CVUniv.sharpRectangularBases}, \lean{CVUniv.residualTheta3Bound}); the file itself
  is correct as an implication.
\item \leanpath{VOA/DeligneAudit/DeligneFusionU7.lean}, module docstring: ``$A_2$ classical freeness rests on the single named residual
  \lean{RavinderKernelLower}''. It no longer does.
\item \leanpath{VOA/DeligneAudit/A2Instance.lean}: ``The three remaining packages \dots\ as explicit hypotheses''; all are now discharged
  for $A_2$.
\item \leanpath{VOA/DeligneAudit/A2Assembly.lean}, \leanpath{VOA/DeligneAudit/A2Headline.lean} and \leanpath{VOA/DeligneAudit/A2Top.lean} list
  as ``remaining'' inputs (\lean{StarIdentity}, \lean{KacWakimotoLowerBound}, \lean{PairComparison}) that are proved in the library.
\end{enumerate}
\end{warning}
Two remarks on formulation. There is no Lean theorem from \lean{IdealLower} directly to the $A_2$ conclusion (the composite is a one-line
term, and the unconditional theorem makes it unnecessary). The short names \lean{explicitPresentation}, \lean{classicallyFree_of_fusion}
and \lean{mldeCoeff} denote different declarations in different namespaces; use full names.

%======================================================================
\section{Referee's checklist}\label{l27:sec:checklist}
%======================================================================

\begingroup\small
\begin{longtable}{p{0.30\textwidth}p{0.30\textwidth}p{0.22\textwidth}p{0.10\textwidth}}
\toprule
Lean declaration & Intended meaning & Hypotheses / status & Fidelity check\\
\midrule\endhead
\lean{A2.}\allowbreak\lean{classicallyFree}, \lean{A2_classicallyFree_}\allowbreak\lean{final} & $L_{-3/2}(\mathfrak{sl}_3)$ classically free,
  $R_V=\operatorname{Sym}\mathfrak{sl}_3/(\ltwentysevenIns)$ & \textbf{none}; aggregate build green 2026-10-28;
  axioms standard & Thm.~\ref{l27:thm:A2}\\
\lean{A2_classicallyFree_}\allowbreak\lean{final_S2} & same & \lean{Ravinder}\allowbreak\lean{KernelLower} & Thm.~\ref{l27:thm:A2}\\
\lean{CaseClassicallyFree} & classical freeness of a presentation & definition & Def.~\ref{l27:def:conclusion}\\
\lean{RelSpec}, \lean{Sat}, \lean{ravSpec}, \lean{ravFam}, \lean{MabRel} & relations of $M_{a,b}$ & definitions & Def.~\ref{l27:def:relspec}\\
\lean{FusionTotalDimG}, \lean{SectorWeyl.}\allowbreak\lean{FusionTotalDim} & Input T & $A_2$: proved; others: \textbf{hypothesis} & Def.~\ref{l27:def:totaldim}\\
\lean{DemazureRealizationG}, \lean{hilbSeries}, \lean{g2HilbSeriesZ} & Input R & $A_2$, $D_4$: proved; others: \textbf{hypothesis} & Def.~\ref{l27:def:realization}\\
\lean{StarBoundQ}, \lean{G2StarBound} \dots\ \lean{E8StarBound} & Input S & \textbf{hypothesis} ($A_2$, $D_4$: proved) & Def.~\ref{l27:def:star}\\
\lean{RavinderKernelLower} & $A_2$ kernel lower bound in a model & not proved; not needed & Def.~\ref{l27:def:a2alt}\\
\lean{fusionTotalDim_}\allowbreak\lean{of_kernelLower} & reduction of the $A_2$ Input T & implication; proved & Thm.~\ref{l27:thm:a2alt}\\
\lean{D4_classicallyFree_final} & $D_4$ target, general form & T, R & Thm.~\ref{l27:thm:D4}\\
\lean{D4_classicallyFree_}\allowbreak\lean{of_inputT} & $D_4$ target & T only; aggregate build green 2026-11-03; axioms standard & Thm.~\ref{l27:thm:d4inputT}\\
\lean{D4Char.}\allowbreak\lean{d4_demazureRealization} & Input R for $D_4$, all $n$ & unconditional; aggregate build green 2026-11-03; axioms standard & Thm.~\ref{l27:thm:d4R}\\
\lean{G2_classicallyFree_final}, \lean{F4_classicallyFree_final} & $G_2$, $F_4$ targets & T, R, S & Thm.~\ref{l27:thm:G2F4}\\
\lean{E6_}/\lean{E7_}/\lean{E8_}\allowbreak\lean{classicallyFree_}\allowbreak\lean{final} & $E$-type targets & T, R, S & Thm.~\ref{l27:thm:E}\\
\lean{A2.}\allowbreak\lean{classicallyFree_of_fusion} and the $A_2$ headline table & $A_2$ from U5, U6 & listed & Thm.~\ref{l27:thm:A2ladder}\\
\lean{partitionsLe_series}, \lean{demazureCoeff_}\allowbreak\lean{eq_}\allowbreak\lean{demazureCoeffOf} & bookkeeping & unconditional & Lem.~\ref{l27:lem:A2book}\\
\lean{W4_centralCharge}, \lean{d4_headline}, \lean{D4Lie.}\allowbreak\lean{classicallyFree_of_fusion}, \lean{classicallyFree_}\allowbreak\lean{of_squeeze} & $D_4$ headlines & listed & Thm.~\ref{l27:thm:D4ladder}\\
\lean{g2_classicallyFree_}\allowbreak\lean{concrete}, \lean{…_of_fermionic}, \lean{f4_classicallyFree_}\allowbreak\lean{concrete'}, \lean{G2DemazureFermionic}, \lean{G2FermionicKW} & $G_2$, $F_4$ headlines & U5, star input & Thm.~\ref{l27:thm:G2F4ladder}\\
\lean{eX_classicallyFree_}\allowbreak\lean{of_starBound} ($X=6,7,8$) & $E$-type headlines & U5, star & Thm.~\ref{l27:thm:Eladder}\\
\lean{D4FusionContact.}\allowbreak\lean{localWeyl_one_total_dim}, \lean{localWeyl_two_total_dim} & Input T for $D_4$ at $m=1,2$ & unconditional & Thm.~\ref{l27:thm:d4Tsmall}\\
\lean{D4Small.}\allowbreak\lean{demazureRealization_zero}, \lean{…_one}, \lean{D4Two.}\allowbreak\lean{demazureRealization_two} & Input R for $D_4$ at $n=0,1,2$ & unconditional & Thm.~\ref{l27:thm:d4Rsmall}\\
(paper only) & $M_{0,m}=W(m\theta,m)$ & on paper; no Lean statement & Lem.~\ref{l27:lem:truncweyl}\\
(paper only) & exceptional star identities & \textbf{conjecture} & Conj.~\ref{l27:conj:star}\\
\bottomrule
\end{longtable}
\endgroup

%======================================================================
\section{Notes and references}\label{l27:sec:notes}
%======================================================================

\textbf{Classical freeness, Li's filtration, the $C_2$-algebra.} Y.~Zhu, ``Modular invariance of characters of vertex operator
algebras'', \emph{J.\ Amer.\ Math.\ Soc.} 9 (1996), 237--302 (the $C_2$-algebra); H.~Li, ``Abelianizing vertex algebras'',
\emph{Comm.\ Math.\ Phys.} 259 (2005), 391--411 (Li's filtration and the surjection from the arc algebra); J.~van~Ekeren and R.~Heluani,
``Chiral homology of elliptic curves and the Zhu algebra'', \emph{Comm.\ Math.\ Phys.} 386 (2021), 495--550, arXiv:1804.00017 (classical
freeness); T.~Arakawa and A.~Moreau, ``Joseph ideals and lisse minimal $W$-algebras'', \emph{J.\ Inst.\ Math.\ Jussieu} 17 (2018), 397--417,
arXiv:1506.00710 (the associated variety of $L_k(\ltwentyseveng)$ for the Deligne series at $k=-1-h^\vee/6$ is the minimal nilpotent orbit
closure); P.~Deligne, ``La s\'erie exceptionnelle de groupes de Lie'', \emph{C.\ R.\ Acad.\ Sci.\ Paris} 322 (1996), 321--326.

\textbf{Fusion products, Weyl and Demazure modules.} B.~Feigin and S.~Loktev, ``On generalized Kostka polynomials and the quantum Verlinde
rule'', Amer.\ Math.\ Soc.\ Transl.\ Ser.~2, vol.~194 (1999), 61--79, arXiv:math/9812093; V.~Chari and A.~Pressley, ``Weyl modules for classical
and quantum affine algebras'', \emph{Represent.\ Theory} 5 (2001), 191--223; V.~Chari and S.~Loktev, ``Weyl, Demazure and fusion modules for the
current algebra of $\mathfrak{sl}_{r+1}$'', \emph{Adv.\ Math.} 207 (2006), 928--960 (the spanning theorem used for the $A_2$ realizations);
G.~Fourier and P.~Littelmann, ``Weyl modules, Demazure modules, KR-modules, crystals, fusion products and limit constructions'',
\emph{Adv.\ Math.} 211 (2007), 566--593, arXiv:math/0509276 ($D(1,j\theta)\cong D(1,\theta)^{*j}$; Theorem~8 is stated for every
simple $\ltwentyseveng$, and Theorem~B identifies $D(1,\omega_i^\vee)$ with a Kirillov--Reshetikhin module);
V.~Chari and A.~Moura, ``Kirillov--Reshetikhin modules associated to $G_2$'', \emph{Contemp.\ Math.} 442 (2007), arXiv:math/0604281;
E.~Ardonne and R.~Kedem, ``Fusion products of Kirillov--Reshetikhin modules and fermionic multiplicity formulas'', \emph{J.\ Algebra}
308 (2007), arXiv:math/0602177 (Theorem~5.2); P.~Di~Francesco and R.~Kedem, ``Proof of the combinatorial Kirillov--Reshetikhin
conjecture'', \emph{Int.\ Math.\ Res.\ Not.} 2008, arXiv:0710.4415 (Corollary~5.19); M.~S.~Lin, ``Quantum $Q$-systems and fermionic
sums---the non-simply laced case'', \emph{Int.\ Math.\ Res.\ Not.} 2021, 805--854, arXiv:1903.11492 (conventions); D.~Hernandez,
\emph{J.\ reine angew.\ Math.} 596 (2006) (the $Q$-system, cited as stated by Di Francesco--Kedem). The theorem statements of the first five of these were read in their arXiv
versions for \texttt{Notes/deligne/G2\_STAR\_FERMIONIC.md}; Hernandez and the HKOTY paper were not opened; V.~Chari, P.~Shereen,
R.~Venkatesh and J.~Wand, arXiv:1408.4090 (cited by Kus--Littelmann, with Fourier--Littelmann, for the decomposition at
$\lambda_0=0$ in all types; this attribution was not verified by the project); V.~Chari and R.~Venkatesh, ``Demazure modules, fusion products and
$Q$-systems'', \emph{Comm.\ Math.\ Phys.} 333 (2015), 799--830, arXiv:1305.2523 (Corollary~3.5, on which Ravinder's Proposition~1
rests; not re-checked by the project); B.~Ravinder, ``Demazure modules, Chari--Venkatesh modules and
fusion products'', \emph{SIGMA} 10 (2014), 110, arXiv:1409.0274 (Proposition~1: the presentation of $D(1,j\theta)$; Theorem~1 and
Corollaries~1--2: the exact sequences and $M_{0,m}\cong\ltwentyseveng^{*m}$, all stated for every simple $\ltwentyseveng$; only
Theorem~2 and Corollary~3 assume $\ltwentyseveng$ simply laced; Proposition~5 quotes Feigin's Corollary~2, so Corollary~2 at $a=0$ is
not independent of Feigin; numbering of arXiv v4); H.~Garland, ``The arithmetic theory of loop algebras'',
\emph{J.\ Algebra} 53 (1978), 480--551 (Garland's lemma).

\textbf{The total bound.} E.~Feigin, ``The PBW filtration, Demazure modules and toroidal current algebras'', \emph{SIGMA} 4 (2008), 070,
arXiv:0806.4851 (Proposition~1 and Corollary~2 in the \emph{SIGMA} and arXiv v2 numbering; Proposition~2.1 and Corollary~2.2 in
arXiv v1; the proof of Proposition~1 has a gap); D.~Kus and P.~Littelmann, ``Fusion products and toroidal algebras'', \emph{Pacific J.\
Math.} 278 (2015), 427--445, arXiv:1411.5272 (Definition~4.1: the truncated Weyl module; Theorem~3.5, Corollary~3.6(3) and
Theorem~4.2: the total bound for every simple $\ltwentyseveng$ when $\lambda_0=0$; p.~437: the step flagged by the project; p.~438:
the remark on Feigin's gap. The hypothesis of Theorem~4.2 was read from a PDF extraction that loses some glyphs); S.~Setia, S.~Rani and T.~Khandai, ``Filtration of tensor
product of local Weyl modules for $\mathfrak{sl}(n+1)[t]$'', arXiv:2405.11944 (2024), \S8 (the route formalized for $A_2$); B.~Feigin, M.~Jimbo,
R.~Kedem, S.~Loktev and T.~Miwa, arXiv:math/0208156, \S2 (the polynomial realization of fusion products).

\textbf{Demazure characters and Macdonald theory.} S.~Kumar, ``Demazure character formula in arbitrary Kac--Moody setting'', \emph{Invent.\
Math.} 89 (1987), 395--423; O.~Mathieu, ``Formules de caract\`eres pour les alg\`ebres de Kac--Moody g\'en\'erales'', \emph{Ast\'erisque}
159--160 (1988); Y.~Sanderson, ``On the connection between Macdonald polynomials and Demazure characters'', \emph{J.\ Algebraic Combin.} 11
(2000), 269--275; B.~Ion, ``Nonsymmetric Macdonald polynomials and Demazure characters'', \emph{Duke Math.\ J.} 116 (2003), 299--318;
I.~Cherednik, ``Nonsymmetric Macdonald polynomials'', \emph{Int.\ Math.\ Res.\ Not.} 1995, 483--515; G.~Hatayama, A.~Kuniba, M.~Okado,
T.~Takagi and Y.~Yamada, ``Remarks on fermionic formula'', \emph{Contemp.\ Math.} 248 (1999), 243--291 (the fermionic formula and the
$X=M$ conjecture).

\textbf{Lattice VOAs and the character side.} I.~Frenkel and V.~Kac, ``Basic representations of affine Lie algebras and dual resonance
models'', \emph{Invent.\ Math.} 62 (1980), 23--66; T.~Arakawa and K.~Kawasetsu, ``Quasi-lisse vertex algebras and modular linear
differential equations'', in \emph{Lie Groups, Geometry, and Representation Theory}, Progr.\ Math.\ 326 (2018), arXiv:1610.05865.

\textbf{Physics.} C.~Beem, M.~Lemos, P.~Liendo, W.~Peelaers, L.~Rastelli and B.~C.~van~Rees, ``Infinite chiral symmetry in four
dimensions'', \emph{Comm.\ Math.\ Phys.} 336 (2015), 1359--1433, arXiv:1312.5344 (including $c_{2d}=-12c_{4d}$); C.~Beem and L.~Rastelli,
``Vertex operator algebras, Higgs branches, and modular differential equations'', \emph{JHEP} 2018(8), 114, arXiv:1707.07679;
J.~A.~Minahan and D.~Nemeschansky, ``An $N=2$ superconformal fixed point with $E_6$ global symmetry'', \emph{Nucl.\ Phys.\ B} 482 (1996),
142--152; A.~D.~Shapere and Y.~Tachikawa, ``Central charges of $\mathcal N=2$ superconformal field theories in four dimensions'',
\emph{JHEP} 2008(09), 109, arXiv:0804.1957.

\endgroup
